% !TeX root = ../Book3.tex
\BookPart{调和函数与势论}{b3:part:03}
% 本编导读（章级标题，不计入编号）
\BookPartGuide

调和函数是全纯函数的实部，也是二维 Laplace 方程的解。前一种身份提供共形不变性，后一种身份允许我们处理边值、正性与测度。本编把这两种语言接在一起：先从平均值与最大值原理出发，再解 Dirichlet 问题，随后用 Green 核记录一个单位点源对整个区域的影响，最后回到全纯函数的零点。
\[
\begin{tikzcd}[column sep=large,row sep=large]
\text{第十五章：调和与次调和函数}\arrow[r]\arrow[dr]&
\text{第十六章：Dirichlet 问题}\arrow[d]\\
\text{第十八章：对数势与零点}&
\text{第十七章：Green 函数}\arrow[l]
\end{tikzcd}
\]
第十五章到第十七章的斜箭头表示最大值原理、分布 Laplace 算子和 Riesz 测度的直接输入；第十八章的 Jensen 公式也可以先从圆盘内的因子分解学习，再回读一般 Green 表述。第十七章关于能量的讨论是势论支线，不是所有零点估计的前提。
\ProseParagraphBreak
阅读时应反复检验三个条件：边界点是否正则，核的奇性是否局部可积，测度或函数列的极限能否与积分交换。光滑边界下的图像并不能自动推广到任意区域。特别是，覆盖意义的双曲型只描述普适覆盖；平面 Green 函数的存在还包含势论条件，两者在多连通区域上不能混为同一概念。
\ProseParagraphBreak
完成本编后，读者应能从 \(\log|f|\) 的平均值读出零点信息，解释 Poisson 延拓为什么必须扣除零点的 Green 项，并把唯一性、存在性和边界连续性分别证明。这样的训练也为下一编从零点到整函数增长的讨论准备了工具。

% !TeX root = ../../Book3.tex
\chapter{调和函数与次调和函数}
\label{b3:ch:15}
\BookChapterNumber{b3:ch:15}
\begin{chapterabstract}
全纯函数的实部满足 Laplace 方程，圆周平均则把这个微分方程转化为比较和估计。本章由平均值性质建立最大值原理、Poisson 表示、Harnack 控制与局部紧致性，再用周期说明调和共轭何时能够拼成单值函数。为了保留全纯函数零点处的对数奇性，随后把平均等式放宽为次调和不等式，并以正分布和 Riesz 测度记录非调和部分。Weyl 引理与局部对数势分解将微分方程、点态代表和测度联系起来，为后续边值问题提供工具。
\end{chapterabstract}

\begin{learninggoals}
\begin{itemize}
\item 从平均值性质恢复调和方程与光滑性，并用 Poisson 核推出内部估计和正调和函数的边界测度表示。
\item 检查 Harnack 比较中的正性和区域条件，说明局部有界性怎样产生紧致性，以及单调性怎样控制调和函数列。
\item 计算共轭微分形式的周期，区分单值调和共轭与全局全纯对数模的存在条件。
\item 在允许负无穷值的情形使用次调和比较、最大值原理及稳定运算，辨明上半连续性与平均不等式各自的作用。
\item 用测试函数计算分布 Laplace 算子，通过正性识别 Riesz 测度，并准确使用分布判据、Weyl 引理和局部 Riesz 分解。
\end{itemize}
\end{learninggoals}

函数 \(\log|z|\) 给出了进入本章的一条线索。在不含原点的小圆盘上，它是某个全纯对数的实部；绕原点一周，共轭函数却增加 \(2\pi\)，因而不能在整个去心平面上取成单值函数。若把原点处的值设为 \(-\infty\)，它仍局部可积，并满足次调和平均不等式；它的分布 Laplace 算子又在原点留下一个点质量。局部方程、全局周期和奇性的测度描述，需要分别建立。

\ProseParagraphBreak
\secref{b3:sec:1501}从圆周平均出发。Poisson 核既构造连续边值的调和延拓，也给出 Harnack 不等式、导数估计和局部紧致性。正调和函数的边界测度表示进一步容纳不连续甚至集中的边界数据；它也提示，几乎处处的径向边值并不总能确定一个调和函数。\secref{b3:sec:1502}再把调和函数接回全纯函数：方程保证局部共轭存在，闭曲线周期决定这些局部共轭能否拼合，而周期属于 \(2\pi\mathbb Z\) 已足以使指数单值。

\ProseParagraphBreak
\secref{b3:sec:1503}用上半连续性和平均不等式处理零点处的对数奇性。比较原理使这一概念可以在小圆盘上验证，也使最大值、拼接和极限运算有统一的依据。\secref{b3:sec:1504}转到分布语言：先确定 \(\Delta\log|z|\) 的常数，再由正性把 Laplace 分布识别为测度。分布只看见几乎处处等价类，因此恢复次调和函数时还须选定点态代表；局部 Riesz 分解中的处处等式也依赖这一步。

\ProseParagraphBreak
本章使用平面微积分、前面的全纯函数局部理论，以及第一卷的局部可积函数、近似单位和正泛函的测度表示。阅读时可对照 Axler、Bourdon、Ramey 的《Harmonic Function Theory》\cite{Axler2001Harmonic}第4--8页的均值与最大值原理、第12--18页的球上 Poisson 表示与均值逆命题，以及第47--50页的 Harnack 理论；比较维数变化时，应重新核对核函数及估计的常数。Ransford 的《Potential Theory in the Complex Plane》\cite{Ransford1995Potential}可作为平面势论的进一步阅读方向。

% !TeX root = ../../Book3.tex
\section{调和函数}
\label{b3:sec:1501}

全纯函数的实部不只是满足一条微分方程。它在每个圆盘中心都等于边界的平均值；反过来，这种平均性质足以恢复光滑性和方程。先把这两个描述联系起来，再从正性得到定量控制。

\subsection{Laplace 方程}
\label{b3:subsec:150101}

设 \(D\subseteq\mathbb C=\mathbb R^2\) 为开集。实值函数 \(u\in C^2(D)\) 若满足 \(\Delta u=u_{xx}+u_{yy}=0\)，称为\term[diaohe-hanshu]{调和函数}[harmonic function]。复值函数调和是指实部、虚部均调和；最大值和正性结论以下只用于实值函数。\(x\)、\(x^2-y^2\)、\(\log|z|\) 分别在全平面、全平面和去心平面上调和。最后一个例子提醒我们：定义域的缺口可能承载奇性，不能在原点直接补值。

由 Cauchy--Riemann 方程和全纯函数的光滑性，\(f=u+\mathrm{i}v\) 全纯时 \(u_{xx}+u_{yy}=v_{yx}-v_{xy}=0\)。实变量的方程在二维保留了复变量的局部结构，但尚未保证全局调和共轭存在。

\subsection{平均值性质}
\label{b3:subsec:150102}

\term[pingjunzhi-xingzhi]{平均值性质}[mean-value property]把方程变成与尺度有关的等式。
\begin{theorem}[调和函数的平均值性质]\label{b3:thm:15-mean-value}
若 \(u\) 在 \(D\) 调和，且 \(\overline{B(a,r)}\subset D\)，则
\[
u(a)=\frac1{2\pi}\int_0^{2\pi}u(a+r e^{\mathrm{i}t})\,\mathrm{d}t
=\frac1{\pi r^2}\int_{B(a,r)}u\,\mathrm{d}A.
\]
反之，连续函数在每个内部圆盘上满足圆周平均等式，则为调和函数。
\end{theorem}
\begin{proof}
圆周平均记作 \(m(r)\)。微分和散度公式给出
\[
 m'(r)=\frac1{2\pi r}\int_{\partial B(a,r)}\partial_nu\,\mathrm{d}s
 =\frac1{2\pi r}\int_{B(a,r)}\Delta u\,\mathrm{d}A=0.
\]
由连续性 \(m(r)\to u(a)\)，得第一式；乘 \(2r\) 并积分得面积版本。
反向证明不能预先假定函数可微。在相对紧的小开集上取光滑径向非负核 \(\eta_\varepsilon\)，积分为一。极坐标和假设给 \(u*\eta_\varepsilon=u\)，故 \(u\) 局部光滑。对光滑函数在中心作二阶 Taylor 展开，圆周平均为 \(u(a)+r^2\Delta u(a)/4+o(r^2)\)。假设使二阶系数为零，故 \(\Delta u=0\)。
\end{proof}
径向磨光同时说明平均值性质为何能恢复所有阶导数：不是从连续性直接求导，而是先将函数识别成它的光滑卷积。

\subsection{最大值原理}
\label{b3:subsec:150103}

\begin{theorem}[强最大值原理]\label{b3:thm:15-maximum}
区域 \(D\) 上的实值调和函数若在内部取得全域最大值或最小值，则为常数。若 \(D\) 有界、\(u\in C(\overline D)\) 且在 \(D\) 调和，则最大值和最小值均可在边界取得。
\end{theorem}
\begin{proof}
设最大值为 \(M=u(a)\)。每个以 \(a\) 为中心的内部圆周上都有 \(u\le M\)，而平均等于 \(M\)。连续非负函数 \(M-u\) 的积分为零，故它在圆周恒零。所有充分小的半径合起来说明 \(u=M\) 的集合开；它由连续性又是闭集。连通性使其等于 \(D\)。最小值情形应用于 \(-u\)。有界域闭包紧，最大值存在；若极值只在内部出现，则上面的结论使函数处处相等，边界亦取得该值。
\end{proof}
有界性与边界条件各有作用。上半平面上的 \(u(x+\mathrm{i}y)=y\) 在实轴为零，却不恒零；没有无穷远控制，边界最大值原理和 Dirichlet 解的唯一性都会失败。

\subsection{Harnack 不等式}
\label{b3:subsec:150104}

为了比较不同点的正调和函数，先建立圆盘表示。在 \(|z|<R\) 记
\[
P_R(z,t)=\frac{R^2-|z|^2}{|Re^{\mathrm{i}t}-z|^2}.
\]
它是 \((Re^{\mathrm{i}t}+z)/(Re^{\mathrm{i}t}-z)\) 的实部，故关于 \(z\) 调和；几何级数表明其圆周平均为一。若 \(\varphi\) 连续，积分 \((2\pi)^{-1}\int P_R(z,t)\varphi(Re^{\mathrm{i}t})\,\mathrm{d}t\) 在内部调和。靠近边界一点时，远离该点的核一致趋零，近处由 \(\varphi\) 的一致连续性控制，所以积分连续取边值 \(\varphi\)。与 \(u\) 比较并用最大值原理，便得到 \(u\) 的 Poisson 表示。

\begin{theorem}[Harnack]\label{b3:thm:15-harnack}
若 \(u\ge0\) 在 \(B(0,R)\) 调和，则 \(|z|=r<R\) 时
\[
\frac{R-r}{R+r}\,u(0)\le u(z)\le\frac{R+r}{R-r}\,u(0).
\]
\end{theorem}
\begin{proof}
先在 \(r<S<R\) 的圆盘用上述表示。核满足 \((S-r)/(S+r)\le P_S(z,t)\le(S+r)/(S-r)\)，对非负边值积分后利用平均值性质得到不等式，再令 \(S\uparrow R\)。
\end{proof}
有限条相交圆盘构成的链把估计传到区域中任意固定紧集。这给出 Harnack 收敛原理：递增调和函数列若在一点有界，就局部一致收敛到调和函数；否则处处趋于正无穷。证明时把差函数代入链式估计，中心处的 Cauchy 性即变成紧集上的一致 Cauchy 性，再用平均值性质的逆命题。


平均值公式还给出内部导数估计。取 \(\overline{B(a,2r)}\subset D\)，在 \(B(a,r)\) 上用半径 \(2r\) 的 Poisson 表示。对核关于 \(z\) 的各阶导数求界，分母距离至少 \(r\)，得到
\[
\sup_{B(a,r)}|\partial^\alpha u|\le C_\alpha r^{-|\alpha|}\sup_{B(a,2r)}|u|.
\]
常数只依赖多重指标，因为把变量缩放到单位圆盘之后，核导数在半径二分之一的闭圆盘一致有界。估计同时控制导数与局部尺度；区域靠近边界时可用的 \(r\) 变小，不能把同一个内部常数当作边界导数界。

\ProseParagraphBreak

若调和函数族在每个紧集上统一有界，以上估计给导数的统一界和等度连续性。对一列这样的函数，先在第一个紧集抽出一致收敛子列，再逐层抽取并取对角列；所得局部一致极限保留平均等式，因此调和。对更高导数重用同一估计，可得到局部各阶导数收敛。这里的紧性机制与全纯正规族相似，但使用的是实值核和平均，而不是复曲线积分。

\ProseParagraphBreak

具体计算可以检验这些结论。边值 \(\varphi(t)=1+\cos t\) 的单位圆盘延拓为 \(1+r\cos\theta\)，非负，且在边界点 \(-1\) 为零。内部仍严格正，因为正调和函数只有在处处零时才可在内部为零。相反，\(u=x\) 在直径上有很多内部零点，却不违背强最小值原理，因为它在圆盘中不非负。不能把“调和函数无内部最小值”误读成“调和函数不能有内部零点”。

\ProseParagraphBreak

Harnack 原理中一点有界的量词也值得区分。若递增列 \(u_n\) 全部调和但不非负，可以先减 \(u_1\)，使差列非负；在一点上界便由 Harnack 链传到任何固定紧集。若没有单调性，一点有界不足以控制全列，例如 \(u_n(z)=n\operatorname{Re}z\) 在原点恒零，在别处发散。真正传播的是非负差函数的中心值，而不是任意函数的一个点值。


\ProseParagraphBreak

有界的整个调和函数也满足 Liouville 性质。设 \(|u|\le M\) 于平面，内部一阶导数估计在任意点、任意半径给 \(|\nabla u(a)|\le CM/R\)，令 \(R\to\infty\) 得梯度零，所以常数。实际上只要求单侧有界即可：若 \(u\ge0\)，对固定 \(z\) 在任意大圆盘应用 Harnack，不等式两端的系数均趋一，故 \(u(z)=u(0)\)。若只要求在某条直线有界，结论不成立，\(u(x+\mathrm{i}y)=y\) 已给反例。全平面的范围在扩大圆盘这一证明中被明确使用。

\ProseParagraphBreak

圆盘上的正调和函数还有边界测度表示，不要求其连续达到边界。对 \(0<s<R\)，把非负测度 \(u(se^{\mathrm{i}t})\mathrm{d}t/(2\pi)\) 放在参数圆周，它的总质量恒为 \(u(0)\)。紧圆周上的有限测度族具有弱收敛子列；此处可用连续函数可数稠密集上的对角抽取构造正泛函，再用正泛函表示得到极限测度 \(\nu\)。固定 \(|z|<R\)，在 \(s>|z|\) 的圆盘使用 Poisson 表示，并让子列 \(s\uparrow R\)，核关于参数一致收敛，得到
\[
u(z)=\int_0^{2\pi}P_R(z,t)\,\mathrm{d}\nu(t),\qquad \nu([0,2\pi))=u(0).
\]
傅里叶展开确定了 \(\nu\) 的全部三角矩，再由三角多项式稠密性确定测度，故表示唯一。这是边界数据由函数扩大到测度的一种自然方式。

\ProseParagraphBreak

特别地，取 \(\nu\) 为一点质量，就得到 \(P_R(z,t_0)\)。它在几乎所有边界点趋零，却在 \(Re^{\mathrm{i}t_0}\) 趋于无穷，内部仍严格正。因而“边界几乎处处为零”不等于经典连续边值为零；若不规定解的函数类，不能从几乎处处边值推出唯一性。Poisson 核的这个例子会在 Hardy 空间中再次出现，届时可积边界代表和测度奇异部分必须分别处理。


\ProseParagraphBreak

正调和函数的导数还可用自身点值控制。若 \(u\ge0\) 调和于 \(B(a,2r)\)，Harnack 给 \(\sup_{B(a,r)}u\le C u(a)\)，再在较小圆盘使用一阶内部估计，得到 \(|\nabla u(a)|\le C u(a)/r\)。当 \(u(a)>0\)，除以自身值可得 \(|\nabla\log u(a)|\le C/r\)。所以正调和函数在一个内部尺度上不能变化得任意快；常数只依赖尺度比而不依赖函数大小。若允许变号，\(u_n=x/n+y\) 等函数在某个点可有零值而梯度非零，便不可能由自身值给同样的界。

\ProseParagraphBreak

平均值与面积测度的归一化在缩放下保持一致。将 \(z=a+r\zeta\) 后，面积变为 \(r^2\mathrm{d}A(\zeta)\)，Laplace 算子变为 \(r^{-2}\Delta_\zeta\)，圆周长度变为 \(r\mathrm{d}s\)。因此圆周平均用 \(2\pi r\) 除，面积平均用 \(\pi r^2\) 除；分布基本解中的点质量则在对应变量替换时保持单位总质量。这些尺度检查会在后面的 Green 核、Riesz 测度与零点计数中反复出现。只看公式的函数部分而忽略测度变量，是丢失二维常数与半径幂的常见原因。

% !TeX root = ../../Book3.tex
\section{调和共轭}
\label{b3:sec:1502}

在二维，寻找调和共轭等价于对一个闭微分形式积分。方程解决局部存在性，沿闭曲线的积分则决定局部对象能否拼成一个全局函数。

\subsection{Cauchy–Riemann 方程}
\label{b3:subsec:150201}

若希望 \(u+\mathrm{i}v\) 全纯，必须满足 \(v_x=-u_y\)、\(v_y=u_x\)。所需的一形式为 \(\omega=-u_y\,\mathrm{d}x+u_x\,\mathrm{d}y\)，其外微分为 \(\Delta u\,\mathrm{d}x\wedge\mathrm{d}y\)。所以调和方程恰好消去局部积分的路径依赖。这里的符号由 Cauchy--Riemann 方程确定，改成 \(u_y\,\mathrm{d}x-u_x\,\mathrm{d}y\) 会得到共轭的负号。

\subsection{局部调和共轭}
\label{b3:subsec:150202}

\begin{proposition}\label{b3:prop:15-local-conjugate}
实值调和函数在每点的某个圆盘邻域具有调和共轭，且共轭只差实常数。
\end{proposition}
\begin{proof}
以圆盘中心 \((a,b)\) 为起点，令
\[
v(x,y)=-\int_a^x u_y(t,b)\,\mathrm{d}t+\int_b^y u_x(x,t)\,\mathrm{d}t.
\]
得到 \(v_y=u_x\)，又由 \(u_{xx}=-u_{yy}\) 得 \(v_x=-u_y(x,b)-\int_b^yu_{yy}(x,t)\,\mathrm{d}t=-u_y(x,y)\)。因此 \(u+\mathrm{i}v\) 全纯，\(v\) 也调和。两个共轭的偏导相同，差在连通圆盘上为常数。
\end{proof}
例如 \(u=x^2-y^2\) 对应 \(v=2xy\)；\(u=\log|z|\) 的局部共轭为某个辐角分支。后一例在绕原点一周后增加 \(2\pi\)，已显示全局障碍。

\subsection{单连通区域}
\label{b3:subsec:150203}

\begin{theorem}\label{b3:thm:15-global-conjugate}
区域 \(D\) 上的实值调和函数 \(u\) 有单值调和共轭，当且仅当每条分段光滑闭曲线 \(\gamma\subset D\) 均满足 \(\int_\gamma\omega=0\)。特别地，单连通区域上一定存在共轭。
\end{theorem}
\begin{proof}
若 \(\omega=\mathrm{d}v\)，闭曲线积分为零。反之，从固定点到 \(z\) 积分 \(\omega\)；两条路径之差是闭曲线，故积分独立于路径，局部微分得到上一小节的偏导。单连通时，每条闭曲线可缩到一点；把同伦像用有限个局部原函数圆盘覆盖、细分参数方格，相邻边积分相消，故积分在同伦中不变。
\end{proof}
去心圆盘上的 \(\log|z|\) 给出 \(\omega=\mathrm{d}\arg z\)，单位圆上积分为 \(2\pi\)，所以单连通条件不能随意删去。

\subsection{全纯函数}
\label{b3:subsec:150204}

如果只需要全纯函数的模，可以先构造局部共轭。不同共轭的差为常数，故局部函数 \(e^{u+\mathrm{i}v}\) 在交集只差模一常数；它们不一定拼成单值函数。若所有周期落在 \(2\pi\mathbb Z\)，指数却能消除差异，得到全局无零点全纯函数。

\begin{proposition}\label{b3:prop:15-log-modulus}
区域上无零点全纯函数 \(f\) 的 \(\log|f|\) 调和。反之，调和函数 \(u\) 能写成 \(\log|f|\)，其中 \(f\) 无零点全纯，当且仅当 \(-u_y\,\mathrm{d}x+u_x\,\mathrm{d}y\) 的所有周期属于 \(2\pi\mathbb Z\)。
\end{proposition}
\begin{proof}
局部对数使第一结论成立，并使必要性化为辐角沿闭曲线的总变化。充分性由上面的路径积分定义 \(e^{u+\mathrm{i}\int\omega}\)，不同路径改变整数倍 \(2\pi\)，故指数相同；局部表示保证全纯。
\end{proof}


可直接从实部计算全纯函数的例子能帮助核对共轭符号。设
\(u=x^3-3xy^2+xy\)。先算 \(\Delta u=0\)，再由 \(v_y=u_x=3x^2-3y^2+y\) 积分得
\(v=3x^2y-y^3+y^2/2+h(x)\)。代入 \(v_x=-u_y=6xy-x\)，得到 \(h'(x)=-x\)，故
\[
v=3x^2y-y^3+\frac{y^2-x^2}{2}+C,
\quad f(z)=z^3-\frac{\mathrm{i}}2z^2+\mathrm{i}C.
\]
这个计算同时检查了调和条件、共轭的相容性和纯虚常数的唯一自由度。若一开始给出的 \(u\) 不调和，例如 \(x^2+y^2\)，两偏导方程会在混合导数处矛盾；这比形式地积分其中一条方程更能发现问题。

\ProseParagraphBreak

不同边界孔的周期可以各自贡献。在去掉 \(a_1,\ldots,a_k\) 的区域上，函数
\(u(z)=\sum c_j\log|z-a_j|\) 调和。沿只绕 \(a_j\) 一周的曲线，周期为 \(2\pi c_j\)。因此全局单值共轭存在当且仅当这些基本周期全部零；全局无零点全纯函数的对数模存在则要求 \(c_j\in\mathbb Z\)。后一情况由 \(f=\prod(z-a_j)^{c_j}\) 明确实现，在当前去点区域内负指数不会引入内部极点。两种存在问题的差异来自指数函数的周期，不是来自不同的调和方程。

\ProseParagraphBreak

单连通性是保证所有调和函数都有共轭的条件，却不是某个特定函数有共轭的必要条件。环域上 \(u=\operatorname{Re}z\) 有共轭 \(\operatorname{Im}z\)，而 \(\log|z|\) 没有。对一个给定 \(u\)，应检查它的闭形式周期；不能只看到域有孔就拒绝任何共轭构造。


\ProseParagraphBreak

局部共轭还使调和函数具有实解析性。在小圆盘写 \(u=\operatorname{Re}f\)，全纯函数的 Taylor 级数局部一致收敛，因此
\[
u(a+re^{\mathrm{i}\theta})=c_0+\sum_{n\ge1}r^n(A_n\cos n\theta+B_n\sin n\theta).
\]
这里 \(c_0\) 为实数，系数由 \(f\) 的复 Taylor 系数拆成实部、虚部得到。每个非恒定项本身调和，且圆周平均零。若 \(u\) 在某个非空开集恒零，则其局部全纯函数实部恒零，Cauchy--Riemann 方程使它为纯虚常数；沿重叠圆盘传播后，\(u\) 在连通区域恒零。这个唯一延拓结论比最大值原理要求的单侧极值条件不同：前者使用开集上的全部零值，后者使用一点极值。

\ProseParagraphBreak

若两个调和函数在一条曲线上相等，唯一延拓却未必适用。\(u=x\) 与零函数在虚轴相等但在两侧不同。若额外给出该直线上的法向导数相等，在具有解析延拓的局部情形，Taylor 系数才能被同时限制。只由曲线上的函数值不能恢复二维的全部局部数据；全纯函数沿有聚点的零集恒等定理不能直接挪给实值调和函数。

\ProseParagraphBreak

Harnack 不等式还把局部正性传遍连通区域。固定紧集 \(K\Subset D\) 和一点 \(z_0\in D\)，用有限条圆盘链把 \(z_0\) 连到覆盖 \(K\) 的小圆盘；每一步比较的常数只依赖圆盘的半径比例。相乘后得到 \(\sup_Ku\le C_{K,z_0}u(z_0)\)，对所有正调和函数使用同一个常数。连通性在这里转化为有限次局部比较，而非直接给一个全域常数。若区域无界或紧集逐渐碰到边界，常数一般随之增大；圆盘 Poisson 核靠近其边界质量所在点时的增长就是一个具体检验。

% !TeX root = ../../Book3.tex
\section{次调和函数}
\label{b3:sec:1503}

全纯函数在零点处的对数模趋于负无穷，不能作为通常的调和函数处理。把平均等式放宽为不等式，并允许这个无穷值，便能保留最大值原理和零点信息。

\subsection{次调和函数}
\label{b3:subsec:150301}

\begin{definition}\label{b3:def:15-subharmonic}
设 \(D\) 为区域，\(u:D\to[-\infty,\infty)\) 上半连续、局部可积且不恒为 \(-\infty\)。若每个内部闭圆盘满足
\[
u(a)\le\frac1{2\pi}\int_0^{2\pi}u(a+r e^{\mathrm{i}t})\,\mathrm{d}t,
\]
则称 \(u\) 为\term[cidiaohe-hanshu]{次调和函数}[subharmonic function]。圆周积分取扩展意义。函数 \(-u\) 称为\term[chaodiaohe-hanshu]{超调和函数}[superharmonic function]。
\end{definition}
\(C^2\) 函数次调和当且仅当 \(\Delta u\ge0\)：必要性用二阶平均展开，充分性用圆周平均导数的面积公式。因此 \(|z|^2\) 次调和而非调和，\(-|z|^2\) 超调和。局部可积条件保证后面的分布运算有定义，也排除处处无穷的退化情形。

\subsection{平均值不等式}
\label{b3:subsec:150302}

圆盘上的比较描述同样有用。若 \(h\) 在圆盘调和、连续到边界，且 \(u\le h\) 于边界，则 \(u\le h\) 于内部。证明可先对 \(u-h+\varepsilon|z|^2\) 使用严格平均不等式：它不能有内部最大值，随后令 \(\varepsilon\downarrow0\)。一般上半连续边值用连续上逼近处理。反过来，连续边值的 Poisson 解与 \(u\) 比较，就给中心平均不等式。

\begin{proposition}\label{b3:prop:15-subharmonic-operations}
有限个次调和函数的最大值次调和；非负线性组合只要不恒为负无穷也次调和。若 \(u_j\downarrow u\) 且极限局部可积，则 \(u\) 次调和。
\end{proposition}
\begin{proof}
取最大值的函数在中心选出达到最大的一项，将其平均用所有项的最大值控制。和的情形由积分线性性。下降极限上半连续；在固定圆周上减去局部共同上界，再用单调收敛，得到平均不等式。局部可积保证其为允许的对象。
\end{proof}
任意上确界一般不保持上半连续；这正是 Perron 法必须处理而不能略去的一点。

\subsection{最大值原理}
\label{b3:subsec:150303}

\begin{proposition}[次调和函数的最大值原理]\label{b3:prop:15-subharmonic-maximum}
次调和函数若在区域内部取得有限最大值，则恒等于该值。若 \(D\) 有界且每个边界点满足 \(\limsup_{z\to\xi}u(z)\le0\)，则 \(u\le0\) 于 \(D\)。
\end{proposition}
\begin{proof}
在最大点，圆周平均与中心值相同。若圆周某处小于最大值，上半连续性给一段严格下降的弧，使平均严格下降，矛盾。故最大值集合既开又闭。边界条件加有限覆盖使某个边界邻域中 \(u\le\varepsilon\)；若内部仍有值大于 \(\varepsilon\)，最大值将在余下的紧集取得，第一结论迫使常数，和边界矛盾。令 \(\varepsilon\downarrow0\)。
\end{proof}

\subsection{\texorpdfstring{\(\log|f|\)}{log|f|}}
\label{b3:subsec:150304}

\begin{proposition}\label{b3:prop:15-log-subharmonic}
非恒零全纯函数 \(f\) 的 \(\log|f|\) 次调和，零点处取 \(-\infty\)。
\end{proposition}
\begin{proof}
局部写 \(f(z)=(z-a)^m g(z)\)，\(g\) 无零点。函数 \(\log|g|\) 调和，\(\log|z-a|\) 局部可积。后者以连续函数 \(\tfrac12\log(|z-a|^2+\varepsilon^2)\) 下降逼近；其 Laplace 算子为 \(2\varepsilon^2/(|z-a|^2+\varepsilon^2)^2\ge0\)。于是它次调和；加上调和函数保持不等式。局部性质通过圆盘比较描述给全域性质。
\end{proof}
这一证明保留零点重数，后面会把 \(m\) 解释为 Laplace 测度在 \(a\) 的质量。不能在零点任意补一个有限值；那会破坏平均不等式；仅仅保持上半连续性并不足够。


严格平均不等式的比较证明可以避免对 \(-\infty\) 点求导。设圆盘 \(B\) 上 \(u\le h\) 于边界，若内部存在 \(u-h>0\)，选充分小 \(\varepsilon>0\)，使 \(u-h+\varepsilon(|z|^2-R^2)\) 仍在某个内部点为正，而边界非正。上半连续性使它在闭圆盘取得有限正最大值。以最大点为中心取一个内部小圆，\(|z|^2\) 的圆周平均比中心多 \(r^2\)，所以中心值严格小于平均，和最大性矛盾。该证明只使用上半连续性、平均不等式及一个严格次调和修正，适用于对数奇性。

\ProseParagraphBreak

在圆盘边界上，上半连续且有上界的数据可由连续函数从上逼近。具体可先将数据截断到 \([-N,M]\)，再用小邻域上确界和平滑分片构造连续上界；截断后的单调逼近回到扩展值。Poisson 核非负且总质量一，使这种上逼近与积分相容。这个过程说明比较描述和圆周平均描述为什么一致，也解释了为什么不能用任意的点态修改代替上半连续代表。

\ProseParagraphBreak

一个常见构造是凸函数作用于调和函数。若 \(u\) 调和、\(\Phi:\mathbb R\to\mathbb R\) 凸且连续，则 Jensen 不等式与平均值性质给
\[
\Phi(u(a))\le\frac1{2\pi}\int_0^{2\pi}\Phi(u(a+re^{\mathrm{i}t}))\,\mathrm{d}t.
\]
故 \(\Phi\circ u\) 次调和。取 \(\Phi(t)=|t|\)、\(e^t\) 可得到许多不是调和的函数；在 \(C^2\) 情况，\(\Delta\Phi(u)=\Phi''(u)|\nabla u|^2\) 显示非调和性集中在凸性的二阶项。凸性不能删，例如 \(\Phi(t)=-t^2\) 与 \(u=x\) 给 Laplace 为负。

\ProseParagraphBreak

两个实部组合成复全纯函数时，\(|f|^p\) 对每个 \(p>0\) 次调和。在无零点局部，用调和函数 \(\log|f|\) 与凸函数 \(e^{pt}\) 即得；在零点处，函数连续、中心值零，平均非负，故平均不等式直接成立。这个结论不需要 \(p\ge1\)，因为这里使用的是指数作用于对数模，而不是直接声称 \(t\mapsto|t|^p\) 在所有实数上凸。

% !TeX root = ../../Book3.tex
\section{分布观点}
\label{b3:sec:1504}

调和与次调和的差别还可以用一个正测度表达。它记录函数偏离平均等式的来源，而不要求函数具有通常的二阶导数。

\subsection{分布意义下的 Laplace 算子}
\label{b3:subsec:150401}

对于 \(u\in L^1_{\mathrm{loc}}(D)\)，定义 \(\langle\Delta u,\varphi\rangle=\int_Du\Delta\varphi\,\mathrm{d}A\)，其中 \(\varphi\in C_c^\infty(D)\)。分布配对不加复共轭，二阶积分分部使符号为正。
\begin{proposition}[对数基本解]\label{b3:prop:15-log-fundamental}
在复平面分布意义下，\(\Delta\log|z|=2\pi\delta_0\)。
\end{proposition}
\begin{proof}
删去半径 \(\varepsilon\) 的圆盘，在余域用 Green 公式。小圆的余域外法向朝向原点，边界项为
\[
\int_0^{2\pi}\bigl(-\varepsilon\log\varepsilon\,\partial_r\varphi(\varepsilon e^{\mathrm{i}t})+\varphi(\varepsilon e^{\mathrm{i}t})\bigr)\,\mathrm{d}t.
\]
第一项趋零，第二项趋 \(2\pi\varphi(0)\)。被删圆盘上 \(\log|z|\) 的可积性保证原积分极限成立。
\end{proof}

\subsection{正分布}
\label{b3:subsec:150402}

实分布 \(T\) 对每个非负测试函数取非负值时，称为\term[zheng-fenbu]{正分布}[positive distribution]。取 \(\chi\in C_c^\infty(D)\)，\(\chi=1\) 于紧集 \(K\) 邻域，则支集在 \(K\) 的实测试函数满足
\[
|T(\varphi)|\le T(\chi)\cdot\|\varphi\|_\infty.
\]
这是由 \(-\|\varphi\|_\infty\chi\le\varphi\le\|\varphi\|_\infty\chi\) 得到的。光滑逼近因此把 \(T\) 唯一延拓到连续紧支集函数，并保持正性。第一卷定理11.4.1（正泛函的 Riesz--Markov 表示）给出局部有限正 Radon 测度。这里调用的是表示定理，不把分布的高阶连续性直接误当作测度性。

\subsection{Riesz 测度}
\label{b3:subsec:150403}

\begin{proposition}\label{b3:prop:15-riesz-measure}
次调和函数 \(u\) 的分布 \(\Delta u\) 为正分布。故有唯一局部有限正测度
\(\mu_u=(2\pi)^{-1}\Delta u\)，称为其\term[riesz-cedu]{Riesz 测度}[Riesz measure]。
\end{proposition}
\begin{proof}
径向非负磨光 \(u_\varepsilon\) 在缩小的区域中仍满足平均不等式；它光滑，故 \(\Delta u_\varepsilon\ge0\)。又 \(u_\varepsilon\to u\) 于局部 \(L^1\)，对非负测试函数取极限得到 \(\int u\Delta\varphi\ge0\)。上一小节把该正分布识别为测度，唯一性由测试函数稠密性得到。
\end{proof}
例如 \(u=|z|^2\) 的测度为 \((2/\pi)\mathrm{d}A\)，而 \(u=\log|z-a|\) 的测度为 \(\delta_a\)。连续分布与点质量都可以出现，不能把 Riesz 测度只理解为离散零点。

\subsection{势论解释}
\label{b3:subsec:150404}

对光滑紧支集非负密度 \(\rho\)，对数势
\[
V(z)=\int_{\mathbb C}\log|z-w|\cdot\rho(w)\,\mathrm{d}A(w)
\]
局部可积，且 \(\Delta V=2\pi\rho\)。换序由紧支集上的对数核局部可积保证。因而若 \(\mu_u=\rho\mathrm{d}A\)，\(u-V\) 的分布 Laplace 算子为零。

分布调和函数的光滑性可以这样恢复：其磨光 \(h_\varepsilon\) 调和；在一个内部圆盘上再用固定径向核 \(\eta\)，平均值性质给 \(h_\varepsilon=h_\varepsilon*\eta\)。局部 \(L^1\) 极限使右端连同每个阶导数一致收敛，得到光滑调和代表 \(h=h*\eta\)。这是 \term[weyl-yinli]{Weyl 引理}[Weyl's lemma]在平面上的证明。

势决定非调和部分，调和函数决定不被测度看见的自由度。在一般 Radon 密度情形，对数势可能在薄集上为负无穷，必须另行选择次调和代表；这里不把局部可积函数的任意零测修改说成同一个点态次调和函数。


正分布的零阶估计还说明其局部有限性。取 \(K\Subset D\)，所有 \(0\le\varphi\le1\)、支集在 \(K\) 的测试函数均满足 \(T(\varphi)\le T(\chi)\)。表示测度在 \(K\) 的质量因此不超过同一有限常数。不同紧集选择不同 \(\chi\)；不能从局部有限性推出全域测度有限。

\ProseParagraphBreak

例如 \(u(z)=|z|^2\) 的 Riesz 测度在每个有界紧集有限，却在全平面质量无穷。\(u=\log|z|\) 的测度总质量一，但是函数在无穷远不有界。测度有限性与函数有界性是不同性质，后面 Green 势的边界控制不能只根据其中一项猜出。

\ProseParagraphBreak

正分布一定是测度，任意分布却不一定如此。点质量的导数 \(\partial_x\delta_0\) 对测试函数的作用是 \(-\partial_x\varphi(0)\)。可取支集缩小、最大模不超过一而原点导数任意大的光滑函数，所以它不满足紧集上的零阶最大范数界，不能由局部有限测度表示。它也不正：非负测试函数在原点的一阶导数可有不同符号，只要原点不是极值点。正性在测度识别中承担的正是把所有阶连续性降为零阶控制的作用。

\ProseParagraphBreak

对于无零点全纯函数，\(\mu_{\log|f|}=0\)，但函数的增长可以很大。例如 \(e^{z^3}\) 的对数模为 \(\operatorname{Re}z^3\)，调和且测度零。因此 Laplace 测度只记录势的源，并不记录全部调和自由度。要确定一个函数，还须给边界数据、增长条件或某种归一化；这也是下一章边值问题与再下一章 Green 最小性各自出现的原因。


\ProseParagraphBreak

对一般有限正测度也可以局部构造对数势。设 \(\mu\) 支集于紧集 \(K\)，定义 \(V(z)=\int_K\log|z-w|\,\mathrm{d}\mu(w)\)，允许负无穷。对任意有界集 \(L\)，
\[
\sup_{w\in K}\int_L|\log|z-w||\,\mathrm{d}A(z)<\infty,
\]
因为单位距离以内的负部由 \(2\pi\int_0^1r|\log r|\,\mathrm{d}r\) 控制，外面的正部由 \(K,L\) 的有界性控制。Tonelli 因而给 \(V\in L^1_{\mathrm{loc}}\)，尤其它不能在整个开集恒为负无穷。分布换序及对数基本解得到 \(\Delta V=2\pi\mu\)。

\ProseParagraphBreak

这个积分本身提供了适当的点态代表。对局部 \(z\) 变量，核有共同有限上界；减去该上界后用 Fatou 引理，得到 \(\limsup V(z_j)\le V(z)\)，故上半连续。每个核的圆周平均不小于中心值，对正测度积分后仍保持该不等式；若负部的换序有疑问，可先以光滑对数核逼近，再用下降极限处理。因此 \(V\) 次调和，而不只是一个 Laplace 正的分布等价类。

\ProseParagraphBreak

给定次调和 \(u\)，在两个同心内部圆盘 \(B'\Subset B\Subset D\) 上取 \(\mu=\mu_u|_{\overline B}\)，上述势满足 \(\Delta(u-V)=0\) 于 \(B\)，Weyl 引理给调和函数 \(h\)，于是 \(u=V+h\) 几乎处处。要恢复处处等式，可利用次调和函数的面积平均在半径趋零时回到自身：平均不等式给下界，上半连续性给上界，所以平均极限等于点值，包括值为负无穷的情形。双方几乎处处相等使所有面积平均相等，最终得到 \(u=V+h\) 处处成立。这是局部 Riesz 分解；把全部测度放入势时，远处边界与增长还须另外控制，本证明只承诺内部局部范围。


\begin{chaptersummary}
平均值性质把调和方程转为局部比较。Poisson 核的正性给出 Harnack 控制，核的可微性给出内部导数估计；这些估计又使局部有界的调和函数族具有紧致性。使用正调和函数的表示或 Harnack 链时，正性承担着控制全部函数值的作用，不能仅以某一点的有界性代替。

\ProseParagraphBreak
调和共轭在小圆盘上总能构造，全局存在却要检验闭曲线周期。周期全为零时，共轭可取成单值函数；周期只要求落在 \(2\pi\mathbb Z\) 时，指数已经单值，因而得到全局的全纯对数模。环域中的对数项展示了这两个条件的差别。

\ProseParagraphBreak
次调和函数以平均不等式保留对数奇性，其 Riesz 测度是 \((2\pi)^{-1}\Delta u\)。正分布的零阶估计保证这一对象确实是局部有限测度，分布判据则从几乎处处等价类中选出唯一次调和代表。对数势具有指定的 Riesz 测度，Weyl 引理把剩余的零 Laplace 分布恢复为调和函数，从而得到局部分解。测度决定非调和部分，边界条件、周期或增长限制还要控制调和部分；后续边值问题将进一步讨论这些条件如何配合。
\end{chaptersummary}

\BookExercises

% 八题均为自拟或对原自拟题的实质改编；新增训练分别涉及非径向周期、径向凸性、单侧增长、曲线源、圆环势、边界唯一性、高阶极值和势的收敛。解答只使用本章及此前工具，未附会第三方题源。
\Needspace{6\baselineskip}
\begin{exercise}[由圆周均值识别环域中的障碍]\label{b3:ex:15-draft-01}
设 \(0<\rho<R\)，实函数 \(u\) 在环域 \(A=\{\rho<|z|<R\}\) 调和，记
\[
 m(r)\coloneqq\frac1{2\pi}\int_0^{2\pi}u(re^{\mathrm it})\,\mathrm dt.
\]
证明对任意 \(\rho<r_1<r_2<R\)，商
\[
 c\coloneqq\frac{m(r_2)-m(r_1)}{\log(r_2/r_1)}
\]
与所选半径无关，并且存在 \(A\) 上的单值全纯函数 \(F\)，使 \(u=c\log|z|+\operatorname{Re}F\)。据此判别何时存在无零点全纯函数 \(f\) 满足 \(\log|f|=u\)，写出所有这样的 \(f\)；给定整数 \(n\ge2\)，再判别这样的 \(f\) 何时具有 \(A\) 上的全纯 \(n\) 次根。
\end{exercise}
\begin{solution}
在内部圆周上对平均求导，并用极坐标中的调和方程，得到
\[
 (rm'(r))'
 =\frac1{2\pi}\int_0^{2\pi}(ru_{rr}+u_r)\,\mathrm dt
 =-\frac1{2\pi r}\int_0^{2\pi}u_{tt}\,\mathrm dt=0.
\]
所以 \(m(r)=c\log r+d\)，所给商确实与半径无关。沿正向圆周积分共轭形式 \(\omega=-u_y\,\mathrm dx+u_x\,\mathrm dy\)，有
\[
 \int_{|z|=r}\omega
 =\int_0^{2\pi}r u_r(re^{\mathrm it})\,\mathrm dt=2\pi c.
\]
环域的闭曲线周期由这一圆周周期生成。函数 \(u-c\log|z|\) 的共轭形式在生成圆周上的积分为零，因而全部周期为零。由全局共轭判据，它是某个单值全纯函数 \(F\) 的实部。这也证明分解中的实系数 \(c\) 唯一；\(F\) 只差纯虚常数。

若 \(u=\log|f|\)，对数模的周期条件要求 \(2\pi c\in2\pi\mathbb Z\)，即 \(c\in\mathbb Z\)。反过来，当 \(c\) 为整数时，\(z^c e^F\) 在环域上全纯且无零点，对数模正是 \(u\)。任何另一解与它的商模恒为1，故为模一常数。因此全部解为
\[
 f(z)=e^{\mathrm i\beta}z^c e^{F(z)},\quad \beta\in\mathbb R.
\]
若 \(f=g^n\)，则 \(g\) 也无零点，且 \(\log|g|=u/n\)。将已经证明的判据用于 \(u/n\)，得到必要条件 \(c/n\in\mathbb Z\)。此条件也充分，因为
\[
 g(z)=e^{\mathrm i\beta/n}z^{c/n}e^{F(z)/n}
\]
给出一个全纯根，其余根只差一个 \(n\) 次单位根。判别只取决于圆周均值的对数斜率，并不要求原函数径向对称。
\end{solution}

\Needspace{6\baselineskip}
\begin{exercise}[径向次调和性与凸函数的斜率]\label{b3:ex:15-draft-02}
设 \(I=(\log\rho,\log R)\)，其中 \(0<\rho<R\)，且 \(\psi:I\to\mathbb R\) 连续。在环域 \(A=\{\rho<|z|<R\}\) 上定义 \(u(z)=\psi(\log|z|)\)。证明 \(u\) 次调和当且仅当 \(\psi\) 凸。凸时令 \(\nu=\mathrm d\psi'_+\) 为其右导数确定的 Stieltjes 测度，证明对 \(\varphi\in C_c(A)\)，
\[
 \int_A\varphi\,\mathrm d\mu_u
 =\frac1{2\pi}\int_I\int_0^{2\pi}
       \varphi(e^{t+\mathrm i\theta})\,\mathrm d\theta\,\mathrm d\nu(t).
\]
求 \(\mu_u(\{|z|=r\})\) 及 \(\mu_u(\{r<|z|<s\})\)，其中 \(\rho<r<s<R\)，并说明导数跳跃如何产生圆周上的质量。
\end{exercise}
\begin{solution}
对实测试函数 \(\varphi\in C_c^\infty(A)\)，记
\[
 \Phi(t,\theta)=\varphi(e^{t+\mathrm i\theta}),\quad
 \overline\Phi(t)=\frac1{2\pi}\int_0^{2\pi}\Phi(t,\theta)\,\mathrm d\theta.
\]
令 \(r=e^t\) 作极坐标换元。面积因子与 Laplace 算子的径向因子相消，角向二阶导数积分为零，于是
\[
 \langle\Delta u,\varphi\rangle
 =\int_I\int_0^{2\pi}\psi(t)
       (\Phi_{tt}+\Phi_{\theta\theta})\,\mathrm d\theta\,\mathrm dt
 =2\pi\int_I\psi(t)\overline\Phi''(t)\,\mathrm dt.
\]
若 \(u\) 次调和，取 \(\varphi(z)=\chi(\log|z|)\)，其中 \(0\le\chi\in C_c^\infty(I)\)，便得 \(\int_I\psi\chi''\ge0\)。这表示一维分布 \(\psi''\) 非负。在任意内部区间上磨光 \(\psi\)，所得光滑函数的二阶导数非负，所以为凸函数；连续性使磨光局部一致收敛到 \(\psi\)，凸不等式随极限保留，故 \(\psi\) 凸。

反过来，有限凸函数局部绝对连续，右导数 \(\psi'_+\) 单调且局部有界，其 Stieltjes 测度 \(\nu\) 为正测度。对光滑紧支集 \(\chi\)，两次分部积分分别按绝对连续函数和 Stieltjes 积分理解，给出
\[
 \int_I\psi\chi''\,\mathrm dt
 =-\int_I\psi'_+\chi'\,\mathrm dt
 =\int_I\chi\,\mathrm d\nu.
\]
代回前式，得到 \(\Delta u\ge0\) 及题中的测度公式，先对光滑测试函数成立，再由测度唯一性延拓到连续紧支集函数。由于 \(u\) 连续，局部磨光回到其原有点值；分布判据因此说明这个 \(u\) 本身次调和。

Stieltjes 测度在一点的质量是右、左极限之差，故
\[
 \mu_u(\{|z|=r\})
 =\psi'_+(\log r)-\psi'_-(\log r),
\]
而开环域的质量为
\[
 \mu_u(\{r<|z|<s\})
 =\psi'_-(\log s)-\psi'_+(\log r).
\]
前一质量按角度均匀分布在圆周上。例如斜率在 \(t_0\) 向上跳 \(a>0\)，就产生总质量 \(a\) 的均匀圆周测度，圆周半径为 \(e^{t_0}\)。若 \(\psi\in C^2\)，则相应的面积密度为 \(\psi''(\log|z|)/(2\pi|z|^2)\)。这两个描述允许光滑密度与圆周质量同时出现。
\end{solution}

\Needspace{6\baselineskip}
\begin{exercise}[单侧增长也限制调和函数的次数]\label{b3:ex:15-draft-03}
设实调和函数 \(u\) 定义在整个复平面，且存在 \(C>0\)、实数 \(\alpha\ge0\)，使
\[
 u(z)\le C(1+|z|)^\alpha.
\]
证明 \(u\) 是次数至多 \(\lfloor\alpha\rfloor\) 的调和多项式；这里没有假设 \(-u\) 满足同样的上界。进一步，若整数 \(N\ge1\) 满足
\[
 \max_{|z|=R}u(z)=o(R^N)\quad(R\to\infty),
\]
证明次数至多 \(N-1\)，并给出说明两种次数界都不能改进的例子。
\end{exercise}
\begin{solution}
先把单侧上界转为较小圆盘上的双侧上界。设 \(M_R\ge0\) 且 \(u\le M_R\) 于 \(B(0,2R)\)。正调和函数 \(M_R-u\) 的 Harnack 不等式给出，对 \(|z|\le R\)，
\[
 M_R-u(z)\le3(M_R-u(0)).
\]
结合 \(u(z)\le M_R\)，得到
\[
 \sup_{|z|\le R}|u(z)|\le2M_R+3|u(0)|.
\]
第一种假设下可取 \(M_R=C(1+2R)^\alpha\)。固定 \(a\in\mathbb C\)，令整数 \(k=\lfloor\alpha\rfloor+1\)。当 \(R\) 充分大时，\(a\) 附近半径与 \(R\) 成固定比例的圆盘包含在 \(B(0,R)\) 内，内部导数估计因此给出，对每个 \(|\beta|=k\)，
\[
 |\partial^\beta u(a)|
 \le C_\beta R^{-k}\bigl(2C(1+2R)^\alpha+3|u(0)|\bigr).
\]
因 \(k>\alpha\)，令 \(R\to\infty\) 得所有 \(k\) 阶导数为零。Taylor 公式于是表明 \(u\) 是次数至多 \(k-1\) 的多项式；它原来满足的调和方程仍成立。

第二种假设下，取
\[
 M_R=\max\left(0,\max_{|z|=2R}u(z)\right)=o(R^N).
\]
最大值原理使它也控制整个 \(B(0,2R)\)。同一 Harnack 估计给 \(\sup_{|z|\le R}|u(z)|=o(R^N)\)，其中常数项因 \(N\ge1\) 可并入小 \(o\)。这次对所有 \(N\) 阶导数使用内部估计，极限为零，故次数至多 \(N-1\)。

对第一种界，\(u(z)=\operatorname{Re}z^{\lfloor\alpha\rfloor}\) 达到所许次数；\(\lfloor\alpha\rfloor=0\) 时取非零常数。对第二种界，\(\operatorname{Re}z^{N-1}\) 满足小 \(o\) 条件且次数恰为 \(N-1\)。单侧条件能够控制次数，关键在于每个大圆盘上都可以先减去一个共同上界，再使用正性。
\end{solution}

\Needspace{6\baselineskip}
\begin{exercise}[双曲线上的 Riesz 测度]\label{b3:ex:15-draft-04}
在复平面上设
\[
 h(z)=\operatorname{Re}(z^2)-1,\quad u(z)=\max(h(z),0),\quad
 \Gamma=\{x^2-y^2=1\}.
\]
直接用测试函数证明
\[
 \mathrm d\mu_u=\frac{|z|}{\pi}\,\mathrm ds_\Gamma,
\]
其中 \(\mathrm ds_\Gamma\) 为 \(\Gamma\) 上的弧长测度。进一步求 \(\mu_u(B(0,R))\) 对所有 \(R>0\) 的显式表达式，并说明 \(\Gamma\) 的两个顶点处有无点质量。
\end{exercise}
\begin{solution}
调和函数 \(h\) 与零函数的最大值次调和。记 \(E=\{h>0\}\)；在其边界 \(\Gamma\) 上，\(h=0\) 且 \(\nabla h=(2x,-2y)\ne0\)。对 \(\varphi\in C_c^\infty(\mathbb C)\)，在 \(E\) 与足够大的圆盘之交上使用 Green 公式。外圆上的测试函数及其导数均为零，故
\[
 \langle\Delta u,\varphi\rangle
 =\int_E h\Delta\varphi\,\mathrm dA
 =\int_\Gamma(h\partial_n\varphi-\varphi\partial_n h)\,\mathrm ds.
\]
\(E\) 的外法向朝向 \(h\) 减小的一侧，所以 \(\partial_n h=-|\nabla h|=-2|z|\)。由 \(h=0\) 得
\[
 \langle\Delta u,\varphi\rangle
 =\int_\Gamma2|z|\varphi\,\mathrm ds,
\]
除以 \(2\pi\) 即得所求测度。

当 \(R\le1\)，开圆盘不与 \(\Gamma\) 相交，质量为零。当 \(R>1\)，两支曲线可分别参数化为
\[
 x=\pm\sqrt{1+y^2},\quad |y|<Y,\quad
 Y=\sqrt{\frac{R^2-1}{2}}.
\]
沿任一支有
\[
 |z|\,\mathrm ds
 =\frac{1+2y^2}{\sqrt{1+y^2}}\,\mathrm dy
 =\mathrm d\bigl(y\sqrt{1+y^2}\bigr).
\]
把两支的积分相加，得到
\[
 \mu_u(B(0,R))=
 \begin{cases}
 0,&0<R\le1,\\[2pt]
 \dfrac2\pi\sqrt{R^4-1},&R>1.
 \end{cases}
\]
两个顶点 \((\pm1,0)\) 都是光滑曲线上的普通点，密度在其附近连续；弧长测度对单点为零，所以没有额外点质量。这里的非调和部分集中在整个曲线界面，不能只从函数不光滑的几个可见位置猜测点质量。
\end{solution}

\Needspace{6\baselineskip}
\begin{exercise}[均匀圆环的对数势]\label{b3:ex:15-draft-05}
设 \(0<a<b\)，令
\[
 \mathrm d\mu(w)=\frac{\mathbf1_{\{a<|w|<b\}}}{\pi(b^2-a^2)}\,\mathrm dA(w),\quad
 V(z)=\int_{\mathbb C}\log|z-w|\,\mathrm d\mu(w).
\]
求 \(V\) 在 \(|z|\le a\)、\(a<|z|<b\)、\(|z|\ge b\) 三个区域的显式表达式。证明势在中空圆盘内为常数，且在两条边界圆周处一阶导数连续；据此直接核对 \(\mu_V=\mu\)，包括边界圆周上是否出现质量。
\end{exercise}
\begin{solution}
记 \(r=|z|\)、\(D=b^2-a^2\)。旋转不变性使 \(V\) 只依赖 \(r\)。当 \(r\ne s\) 时，提出较大的半径，再对无零点函数的全纯对数实部使用平均值性质，可得
\[
 \frac1{2\pi}\int_0^{2\pi}\log|r-se^{\mathrm it}|\,\mathrm dt
 =\log\max(r,s).
\]
例如 \(r<s\) 时，余项为 \(\log|1-(r/s)e^{-\mathrm it}|\)，平均为零。对数核的局部可积性允许极坐标换序；参数 \(s=r\) 这单个半径不影响面积积分。因此
\[
 V(r)=\frac2D\int_a^b s\log\max(r,s)\,\mathrm ds.
\]
利用 \(\int s\log s\,\mathrm ds=(s^2/2)\log s-s^2/4\)，分别计算得到
\[
 V(r)=
 \begin{cases}
 \dfrac{b^2\log b-a^2\log a}{D}-\dfrac12,&0\le r\le a,\\[6pt]
 \dfrac{b^2\log b-a^2\log r-(b^2-r^2)/2}{D},&a<r<b,\\[6pt]
 \log r,&r\ge b.
 \end{cases}
\]
第一段显示中空圆盘内的势为常数，第三段则等于总质量为1的原点点质量在外部产生的势；在中间的有质量区域，二者都不能代替实际的势。

三段在端点取值一致，径向导数分别为
\[
 V'(r)=
 \begin{cases}
 0,&0<r<a,\\[2pt]
 \dfrac{r^2-a^2}{Dr},&a<r<b,\\[6pt]
 \dfrac1r,&r>b.
 \end{cases}
\]
在 \(a\) 处中间导数趋于0，在 \(b\) 处趋于 \(1/b\)，故 \(V\) 在整个平面为 \(C^1\) 函数。在开圆环内，\(V''+V'/r=2/D\)，而在其余两区为零。分区使用 Green 公式时，函数值和法向一阶导数的边界项都相消，所以分布 Laplace 算子没有附加的圆周项。因此
\[
 \frac1{2\pi}\Delta V
 =\frac{\mathbf1_{\{a<|z|<b\}}}{\pi D}\,\mathrm dA=\mu.
\]
\end{solution}

\Needspace{6\baselineskip}
\begin{exercise}[几乎处处边值何时能够确定调和函数]\label{b3:ex:15-draft-06}
设实函数 \(u\) 在 \(\mathbb D\) 调和，且对几乎每个 \(t\in[0,2\pi)\)，有 \(u(re^{\mathrm it})\to0\) 当 \(r\uparrow1\)。若存在 \(1<p<\infty\)，使
\[
 \sup_{0<r<1}\frac1{2\pi}\int_0^{2\pi}|u(re^{\mathrm it})|^p\,\mathrm dt<\infty,
\]
证明 \(u\equiv0\)。再给出满足同样径向边值条件、所有圆周 \(L^1\) 范数均等于1的非零正调和函数，说明把 \(p>1\) 改成 \(p=1\) 后结论不成立。
\end{exercise}
\begin{solution}
以 \(\mathrm dm=\mathrm dt/(2\pi)\) 表示归一化圆周测度，令 \(\|u(re^{\mathrm it})\|_{L^p(m)}\le M\)。固定 \(\varepsilon>0\)，设
\[
 E_r=\{t:|u(re^{\mathrm it})|>\varepsilon\}.
\]
几乎处处收敛使 \(m(E_r)\to0\)。在 \(E_r\) 上用 Hölder 不等式，在其补集上用 \(\varepsilon\) 控制，得到
\[
 \int|u(re^{\mathrm it})|\,\mathrm dm
 \le\varepsilon+M\,m(E_r)^{1-1/p}.
\]
先令 \(r\uparrow1\)，再令 \(\varepsilon\downarrow0\)，证明圆周 \(L^1\) 范数趋于零。正的指数 \(1-1/p\) 在这一步不可省略。

固定 \(z\in\mathbb D\)。对 \(|z|<r<1\) 使用半径 \(r\) 的 Poisson 表示：
\[
 u(z)=\int_0^{2\pi}
       \frac{r^2-|z|^2}{|re^{\mathrm it}-z|^2}
       u(re^{\mathrm it})\,\frac{\mathrm dt}{2\pi}.
\]
当 \(r\) 充分接近1时，核由 \((r+|z|)/(r-|z|)\) 一致控制。右端的绝对值因圆周 \(L^1\) 范数趋零而趋零，故 \(u(z)=0\)。任意性给 \(u\equiv0\)。

反例为边界点1的 Poisson 核
\[
 v(z)=\operatorname{Re}\frac{1+z}{1-z}
     =\frac{1-|z|^2}{|1-z|^2}.
\]
它正调和，\(v(0)=1\)，且对 \(t\not\equiv0\pmod{2\pi}\)，有 \(v(re^{\mathrm it})\to0\)。由正性和均值性质，\(\int|v(re^{\mathrm it})|\,\mathrm dm=1\) 对每个 \(r\) 成立。圆周质量越来越集中在角度0附近，正是这种集中使有界的 \(L^1\) 范数无法推出 \(L^1\) 收敛。
\end{solution}

\Needspace{6\baselineskip}
\begin{exercise}[正调和函数的高阶方向导数]\label{b3:ex:15-draft-07}
设 \(u\ge0\) 在 \(\mathbb D\) 调和，\(u(0)=1\)，整数 \(n\ge2\)。对 \(\theta\in\mathbb R\)，定义方向导数
\[
 \mathrm D_\theta^n u(0)
 \coloneqq\left.\frac{\mathrm d^n}{\mathrm ds^n}u(se^{\mathrm i\theta})\right|_{s=0}.
\]
证明 \(|\mathrm D_\theta^n u(0)|\le2n!\)，并用边界表示测度的支集分别刻画正、负两种取等情形。若再有 \(u(z)=1+O(|z|^n)\) 当 \(z\to0\)，证明正向取等 \(\mathrm D_\theta^n u(0)=2n!\) 唯一确定
\[
 u(z)=\operatorname{Re}\frac{1+e^{-\mathrm in\theta}z^n}{1-e^{-\mathrm in\theta}z^n}.
\]
\end{exercise}
\begin{solution}
令 \(\nu\) 为 \(u\) 的边界表示测度，它是单位圆周上的概率测度。Poisson 核在每个内部紧圆盘上一致展开为
\[
 \frac{1-|z|^2}{|\zeta-z|^2}
 =1+2\operatorname{Re}\sum_{k=1}^\infty z^k\overline\zeta^{\,k},
 \quad |\zeta|=1.
\]
积分并沿给定方向求 \(n\) 阶导数，得到
\[
 \mathrm D_\theta^n u(0)
 =2n!\int_{\partial\mathbb D}
       \operatorname{Re}(e^{\mathrm in\theta}\overline\zeta^{\,n})\,\mathrm d\nu(\zeta).
\]
被积函数介于 \(-1\) 与1之间，所求界立即成立。

积分等于1，当且仅当 \(\nu\) 集中在被积函数等于1的地方；这是因为非负连续函数 \(1-\operatorname{Re}(e^{\mathrm in\theta}\overline\zeta^{\,n})\) 的积分必须为零。因此正向取等当且仅当
\[
 \operatorname{supp}\nu\subseteq
 \{\zeta\in\partial\mathbb D:\zeta^n=e^{\mathrm in\theta}\}.
\]
同理，负向取等当且仅当支集包含在 \(\{\zeta:\zeta^n=-e^{\mathrm in\theta}\}\) 中。两种集合都各有 \(n\) 个点，其上的任意概率权重都给出取等函数，所以常数确实达到，且一般并不只有一个取等函数。

在正向取等情形，写 \(\omega=e^{2\pi\mathrm i/n}\)，则
\[
 \nu=\sum_{j=0}^{n-1}a_j\delta_{e^{\mathrm i\theta}\omega^j},
 \quad a_j\ge0,\quad\sum_j a_j=1.
\]
条件 \(u(z)-1=O(|z|^n)\) 使前 \(n-1\) 个复矩 \(\int\overline\zeta^{\,k}\,\mathrm d\nu\) 全为零：若其中首个非零矩的阶数为 \(k<n\)，Poisson 展开在某个方向就有非零的 \(k\) 阶首项，与该增长条件矛盾。于是
\[
 \sum_{j=0}^{n-1}a_j\omega^{-jk}=0\quad(1\le k<n).
\]
连同 \(k=0\) 时和为1，对这些等式作有限 Fourier 反演，或直接用单位根的正交关系，可得 \(a_j=1/n\)。均匀权重使全部非 \(n\) 的倍数阶矩消失，而第 \(\ell n\) 阶矩为 \(e^{-\mathrm i\ell n\theta}\)。因此
\[
 u(z)=1+2\operatorname{Re}\sum_{\ell=1}^\infty
                 (e^{-\mathrm in\theta}z^n)^\ell
     =\operatorname{Re}\frac{1+e^{-\mathrm in\theta}z^n}{1-e^{-\mathrm in\theta}z^n}.
\]
它确实正调和并满足全部要求，唯一性亦由权重的确定得到。
\end{solution}

\Needspace{6\baselineskip}
\begin{exercise}[测度弱收敛与对数势的两种收敛]\label{b3:ex:15-draft-08}
设 \(\mu_j,\mu\) 为支集都包含在同一紧集 \(K\subset\mathbb C\) 中的有限正测度，且对每个 \(f\in C(K)\)，有 \(\int f\,\mathrm d\mu_j\to\int f\,\mathrm d\mu\)。记其对数势为
\[
 V_j(z)=\int_K\log|z-w|\,\mathrm d\mu_j(w),\quad
 V(z)=\int_K\log|z-w|\,\mathrm d\mu(w).
\]
证明 \(V_j\to V\) 于 \(L^1_{\mathrm{loc}}(\mathbb C)\)。再构造一列概率测度，使上述弱收敛成立，但在一个固定点 \(z_0\) 处，\(V_j(z_0)\) 均有限且不趋于 \(V(z_0)\)。
\end{exercise}
\begin{solution}
对常函数1应用弱收敛，得到总质量有共同有限上界 \(C\)，并可把 \(\mu(K)\) 也包括在这个上界内。固定紧圆盘 \(L\)，对 \(M>0\) 定义下截断核
\[
 k_M(z,w)=\max(\log|z-w|,-M),
\]
在 \(z=w\) 时取值 \(-M\)。它在 \(L\times K\) 上连续有界。记相应的截断势为 \(V_{j,M}\)、\(V_M\)；对每个固定 \(z\)，弱收敛给 \(V_{j,M}(z)\to V_M(z)\)，共同的质量上界又提供关于 \(z\in L\) 的一致常数界。支配收敛定理因此给
\[
 \|V_{j,M}-V_M\|_{L^1(L)}\longrightarrow0
 \quad\text{对每个固定的 }M.
\]

截断误差可以独立于测度的集中位置来控制。Tonelli 定理和整个小圆盘上的估计给出
\begin{align*}
 \|V_{j,M}-V_j\|_{L^1(L)}
 &\le C\int_{|s|<e^{-M}}(-M-\log|s|)\,\mathrm dA(s)\\
 &=2\pi C\int_0^{e^{-M}}r(-M-\log r)\,\mathrm dr
 =\frac{\pi C}{2}e^{-2M}.
\end{align*}
同一估计适用于 \(V_M-V\)。这些势可能在个别点为 \(-\infty\)，但局部可积性保证此处的 \(L^1\) 计算有意义。由三角不等式，先令 \(j\to\infty\)，再令 \(M\to\infty\)，得到 \(\|V_j-V\|_{L^1(L)}\to0\)。紧圆盘任意，故有局部 \(L^1\) 收敛。

点态结论的反例可保持总质量始终为1。对 \(j\ge2\)，令
\[
 \mu_j=\left(1-\frac1j\right)\delta_1
              +\frac1j\delta_{e^{-j}},\quad \mu=\delta_1.
\]
所有支集均包含在实线段 \([0,1]\) 中；任意连续测试函数有界，而第二项的质量趋零，所以 \(\mu_j\) 弱收敛到 \(\mu\)。在固定观察点 \(z_0=0\)，却有
\[
 V_j(0)=\left(1-\frac1j\right)\log1
                  +\frac1j\log(e^{-j})=-1,\quad V(0)=0.
\]
每个 \(V_j(0)\) 都有限，但不收敛到极限势的点值。越来越靠近观察点的少量质量可以被对数奇性放大；核的积分尾界仍足以保证局部 \(L^1\) 收敛，却不控制这个固定点的值。
\end{solution}

% !TeX root = ../../Book3.tex
\chapter{Dirichlet 问题与 Perron 方法}
\label{b3:ch:16}
\BookChapterNumber{b3:ch:16}
% 本章摘要（不计入编号）
\begin{chapterabstract}
圆盘 Poisson 积分可以直接从边值计算内部函数。一般区域的解则由 Perron 次函数族构造，是否达到每个边界值取决于障碍函数。本章完成包络的调和性证明，并在正则域上把解表示为边界调和测度的积分。
\end{chapterabstract}

% 本章目标（不计入编号）
\begin{learninggoals}
能计算圆盤 Poisson 延拓；能说明 Perron 族的调和提升为何保留次函数条件；能构造外锥障碍并检验边界极限；能在给定边界延拓条件下使用调和测度的共形不变性。
\end{learninggoals}

在单位圆周给定 \(\cos(3t)\)，内部解是 \(r^3\cos(3\theta)\)。如果把圆盘中心挖去，再独立给这个孤立边界点赋值，解却未必接受该数据。一个边值问题因而包含两项责任：构造内部调和函数，证明它达到边界。前者由平均值、比较原理和 Harnack 收敛解决，后者由障碍解决。

\ProseParagraphBreak

本章使用\chapref{b3:ch:15}的次调和比较与 Harnack 原理。圆盘核的计算可与 Axler、Bourdon、Ramey 的《Harmonic Function Theory》\cite{Axler2001Harmonic}的本地背景材料比较；一般域的势论继续阅读 Ransford 的《Potential Theory in the Complex Plane》\cite{Ransford1995Potential}。

% !TeX root = ../../Book3.tex
\section{Dirichlet 问题}
\label{b3:sec:1601}

给定边界温度，怎样确定内部平衡状态？第十五章已经用 Poisson 积分解决了圆盘上的连续边值问题。本节从这个解出发，研究边界误差怎样传入内部、不同频率怎样衰减，以及边值不连续时应当采用什么极限。上半平面的公式还会说明，无界区域的唯一性需要另加增长限制。

\subsection{边值问题}
\label{b3:subsec:160101}

设 \(D\subseteq\mathbb C\) 为有界区域，\(\varphi\in C(\partial D;\mathbb R)\)。寻找一个 \(u\in C(\overline D;\mathbb R)\)，使
\[
 \Delta u=0\quad\text{于 }D,
 \quad
 u=\varphi\quad\text{于 }\partial D,
\]
称为给定数据 \(\varphi\) 的\term[dirichlet-wenti]{Dirichlet 问题}[Dirichlet problem]。这里方程要求 \(u|_D\in C^2(D)\)，边界条件则按闭包上的连续性理解；满足这些要求的函数称为经典解。域内没有热源对应 Laplace 方程，边界温度则选择众多调和函数中的一个。

\ProseParagraphBreak
唯一性与存在性承担不同的任务。若 \(u,v\) 都是同一数据的经典解，则 \(u-v\) 在有界域 \(D\) 调和，连续到 \(\overline D\)，边界值为零。对差及其负号应用最大值原理，得到 \(u=v\)。这个论证不要求边界光滑，却没有构造任何解。去心圆盘上在孤立边界点另指定一个温度，就可能得不到经典解；下一节将给出具体证明。

\ProseParagraphBreak
比较原理同时给出数据扰动的控制。设 \(\varphi,\psi\in C(\partial D;\mathbb R)\) 分别已有经典解 \(u,v\)，并令
\[
 d=\sup_{\partial D}|\varphi-\psi|.
\]
函数 \(u-v-d\) 与 \(v-u-d\) 在边界非正，因此
\[
 \sup_D|u-v|\le d.
\]
若 \(\varphi\le\psi\)，对 \(u-v\) 直接比较还得 \(u\le v\)。所以一旦存在性成立，边界数据的最大误差不会在内部放大。特别地，若一列可解的连续边值一致收敛，则其解在整个闭包上一致 Cauchy；极限的内部调和性可由圆周平均等式传给极限，再用\cref{b3:thm:15-mean-value}的连续函数逆命题得到。

\ProseParagraphBreak
边界上的同一个公式，不能不加检验地用作内部函数。例如单位圆周上的 \(|z|^2\) 恒等于一，故其调和延拓是常数一；函数 \(|z|^2=x^2+y^2\) 本身满足 \(\Delta|z|^2=4\)，并不是这份边值的解。边界条件只规定圆周上的值，域内公式还必须满足方程。

\ProseParagraphBreak
实值连续数据是经典问题的基本范围。复值数据可以分别延拓实部、虚部；有界但不连续的数据仍可作 Poisson 积分，不过此时必须逐点说明边界极限，不能预先要求解属于 \(C(\overline D)\)。

\subsection{圆盘 Poisson 核}
\label{b3:subsec:160102}

圆盘中的构造要同时满足两件事：对内部点保持调和性，对接近的边界点集中权重。设 \(R>0\)，\(|z|<R\)。第十五章引入的 Poisson 核为
\[
 P_R(z,t)=\frac{R^2-|z|^2}{|Re^{\mathrm i t}-z|^2}
 =\operatorname{Re}\frac{Re^{\mathrm i t}+z}{Re^{\mathrm i t}-z}.
\]
右端是关于 \(z\) 的全纯函数的实部，因而调和；左端直接显示它为正。分子中的 \(R^2-|z|^2\) 还给出正确的归一化：
\[
 P_R(z,t)=1+2\operatorname{Re}\sum_{n\ge1}
       \left(\frac zR\right)^n e^{-\mathrm i nt},
 \quad
 \frac1{2\pi}\int_{-\pi}^{\pi}P_R(z,t)\,\mathrm dt=1.
\]
几何级数在每个较小闭盘上关于 \(t\) 一致收敛，所以可逐项积分；所有非零频率的积分均为零。保持常数与集中质量是核逼近边值时各自需要的性质。

\ProseParagraphBreak
把 \(z=re^{\mathrm i\theta}\)、\(\rho=r/R\) 代入，可把半径尺度从角变量中分离出来。记
\[
 p_\rho(s)\coloneqq
 \frac{1-\rho^2}{1-2\rho\cos s+\rho^2},
 \quad 0\le\rho<1.
\]
则 \(P_R(re^{\mathrm i\theta},t)=p_\rho(t-\theta)\)，而 \(p_\rho\) 为偶的 \(2\pi\) 周期函数。接近边界时，控制它的尺度是 \(\varepsilon=1-\rho\)，即到边界的径向距离除以 \(R\)。

\Needspace{6\baselineskip}
\begin{lemma}[Poisson 核的集中估计]\label{b3:lem:16-poisson-localization}
设 \(1/2\le\rho<1\)，\(\varepsilon=1-\rho\)，并令
\(p_\rho(t)=(1-\rho^2)/(1-2\rho\cos t+\rho^2)\)。对 \(|t|\le\pi\)，有
\[
 \frac{\varepsilon}{\varepsilon^2+t^2}
 \le p_\rho(t)
 \le\frac{\pi^2\varepsilon}{\varepsilon^2+t^2}.
\]
特别地，对 \(0<\delta\le\pi\)，
\[
 \frac1{2\pi}\int_{\delta\le|t|\le\pi}p_\rho(t)\,\mathrm dt
 \le\frac{\pi\varepsilon}{\delta}.
\]
\end{lemma}
\begin{proof}
分母可写为
\[
 1-2\rho\cos t+\rho^2
 =\varepsilon^2+4\rho\sin^2(t/2).
\]
在 \(|t|\le\pi\) 上，\(|t|/\pi\le|\sin(t/2)|\le|t|/2\)，故
\[
 \frac2{\pi^2}(\varepsilon^2+t^2)
 \le\varepsilon^2+4\rho\sin^2(t/2)
 \le\varepsilon^2+t^2.
\]
分子为 \(\varepsilon(1+\rho)\)，介于 \(\varepsilon\) 与 \(2\varepsilon\) 之间，得到双边估计。利用偶性和 \(t\ge\delta\) 时 \(\varepsilon^2+t^2\ge t^2\)，进一步有
\[
 \frac1{2\pi}\int_{\delta\le|t|\le\pi}p_\rho(t)\,\mathrm dt
 \le\pi\varepsilon\int_\delta^\infty\frac{\mathrm dt}{t^2}
 =\frac{\pi\varepsilon}{\delta}.
\]
\end{proof}

在中心 \(\rho=0\) 时，核恒等于一。接近边界时，它在 \(t=0\) 的高度为 \((1+\rho)/(1-\rho)\)，在 \(|t|\) 与 \(1-\rho\) 同阶的角区间中仍有相同量级。高度增大、宽度缩小，而总积分保持 \(2\pi\)；任何固定远弧的质量则由上面的估计趋于零。这个尺度关系将同时给出边界收敛速度和跳跃点的极限。

\subsection{Poisson 积分}
\label{b3:subsec:160103}

用正的归一化核对边界数据加权，便得到圆盘解。下面保留一个可直接引用的 Dirichlet 问题版本；圆盘上的存在性与边界连续性的证明已经在上一章完成。

\Needspace{6\baselineskip}
\begin{theorem}[圆盘 Dirichlet 问题]\label{b3:thm:16-poisson}
设 \(R>0\)，\(\varphi:\partial B(0,R)\to\mathbb R\) 连续，并令
\[
 P_R(z,t)=\frac{R^2-|z|^2}{|Re^{\mathrm i t}-z|^2},
 \quad |z|<R.
\]
则存在唯一 \(u\in C(\overline{B(0,R)};\mathbb R)\)，在 \(B(0,R)\) 调和且在边界等于 \(\varphi\)。它在内部由
\[
 u(z)=\frac1{2\pi}\int_0^{2\pi}
 P_R(z,t)\varphi(Re^{\mathrm i t})\,\mathrm dt
\]
给出。
\end{theorem}
\begin{proof}
\cref{b3:thm:15-poisson-continuous}适用于任意正半径圆盘上的实值连续圆周数据，正是这里的条件。该定理用内部紧集上的核导数界证明积分调和，用近弧上的连续性和远弧上的质量消失证明任意方向的边界连续性，再用有界域最大值原理证明唯一性。因此把其 Poisson 积分用于当前 \(\varphi\)，便得到所述经典解。
\end{proof}

记这一延拓为 \(\mathcal P_R\varphi\)。积分直接显示它对数据线性、保持常数并保持大小关系，且
\[
 |(\mathcal P_R\varphi)(z)|
 \le\|\varphi\|_\infty,
 \quad
 (\mathcal P_R\varphi)(0)
 =\frac1{2\pi}\int_0^{2\pi}\varphi(Re^{\mathrm i t})\,\mathrm dt.
\]
中心只看到边界平均，其余角向信息随着离开中心逐渐出现。为精确描述这种变化，先把边界数据与各个角频率配对。

\Needspace{6\baselineskip}
\begin{proposition}[角频率的延拓与衰减]\label{b3:prop:16-poisson-frequencies}
设 \(R>0\)，\(\varphi:\partial B(0,R)\to\mathbb C\) 为有界 Borel 可测函数，\(u\) 由其 Poisson 积分定义。令
\[
 c_n=\frac1{2\pi}\int_0^{2\pi}
       \varphi(Re^{\mathrm i t})e^{-\mathrm i nt}\,\mathrm dt,
 \quad n\in\mathbb Z.
\]
则对 \(0\le r<R\)，
\[
 u(re^{\mathrm i\theta})
 =\sum_{n\in\mathbb Z}c_n
       \left(\frac rR\right)^{|n|}e^{\mathrm i n\theta}.
\]
该级数在每个内部闭盘上绝对一致收敛，且可按实变量 \(x,y\) 逐项求任意阶导数。其和的实部与虚部均调和。若 \(0<q<1\)、\(N\ge0\) 为整数，则截断到 \(|n|\le N\) 后，在 \(|z|\le qR\) 上的误差不超过
\[
 \frac{2\|\varphi\|_\infty q^{N+1}}{1-q}.
\]
\end{proposition}
\begin{proof}
有界性给出 \(|c_n|\le\|\varphi\|_\infty\)。核的几何级数展开在 \(|z|\le qR\) 上关于 \(t\) 绝对一致收敛，乘以有界可测的边界数据后仍可交换求和与积分。于是
\[
 u(z)=c_0+\sum_{n\ge1}
 \left[c_n\left(\frac zR\right)^n+
       c_{-n}\left(\frac{\overline z}R\right)^n\right].
\]
这与角变量表达相同。在 \(|z|\le qR\) 上，取 \(k\) 阶实偏导数后，每个 \(n\ge k\) 的项可由常数乘以 \(n^kq^{n-k}\) 控制；这个数列可求和。有限个低次项也可直接求导。因此各阶导数级数一致收敛，每项的实部、虚部都调和，故和也调和。最后对 \(|n|>N\) 的项取绝对值并求几何级数，得到所述余项界。
\end{proof}

系数 \(c_n\) 称为边界数据的\term[fourier-xishu]{Fourier 系数}[Fourier coefficient]。上面的推导只用核的几何级数，不要求边界函数自身的 Fourier 级数在任何边界点收敛。因子 \((r/R)^{|n|}\) 已使内部展开绝对收敛；当 \(r/R\) 固定小于一时，较高频率受到更强的压制。

\ProseParagraphBreak
对 \(n\ge1\)，数据 \(\cos(nt)\)、\(\sin(nt)\) 分别延拓为
\[
 \left(\frac rR\right)^n\cos(n\theta),
 \quad
 \left(\frac rR\right)^n\sin(n\theta).
\]
正频率 \(e^{\mathrm i nt}\) 对应 \((z/R)^n\)，负频率对应 \((\overline z/R)^{|n|}\)。实数据满足 \(c_{-n}=\overline{c_n}\)，两部分配对成为实函数；调和延拓并不要求负频率消失。

\ProseParagraphBreak
例如圆周上的数据 \(\varphi(z)=(\operatorname{Re}z)^2\) 可写成
\[
 \varphi(Re^{\mathrm i t})
 =R^2\cos^2t=\frac{R^2}{2}(1+\cos2t).
\]
它的调和延拓因而为
\[
 u(x+\mathrm i y)=\frac{R^2+x^2-y^2}{2}.
\]
在圆周上由 \(x^2+y^2=R^2\) 得 \(u=x^2\)，内部则有 \(\Delta u=0\)。中心值 \(R^2/2\) 是边界平均，不是把边界公式 \(x^2\) 直接代入中心所得的零。

\subsection{Poisson 积分的边界行为与调和测度}
\label{b3:subsec:160104}

连续边值已经保证任意方向的边界连续性。沿同心圆趋近时，核的尺度估计还能把数据的连续模转成统一误差，说明光滑程度怎样影响逼近速度。

\Needspace{6\baselineskip}
\begin{proposition}[径向边界误差]\label{b3:prop:16-poisson-boundary-rate}
设 \(R>0\)，\(\varphi:\partial B(0,R)\to\mathbb R\) 连续，\(u=\mathcal P_R\varphi\)。定义角变量的连续模
\[
 \omega_\varphi(\delta)=
 \sup_{\substack{\theta\in\mathbb R\\|t|\le\delta}}
 |\varphi(Re^{\mathrm i(\theta+t)})
       -\varphi(Re^{\mathrm i\theta})|,
 \quad 0<\delta\le\pi.
\]
若 \(R/2\le r<R\)，\(\varepsilon=1-r/R\)，则有绝对常数 \(C\)，使
\[
 \sup_\theta|u(re^{\mathrm i\theta})-\varphi(Re^{\mathrm i\theta})|
 \le\omega_\varphi(\delta)
       +C\|\varphi\|_\infty\frac{\varepsilon}{\delta}.
\]
因此左端随 \(r\uparrow R\) 趋于零。若进一步有
\[
 |\varphi(Re^{\mathrm i(\theta+t)})
       -\varphi(Re^{\mathrm i\theta})|
 \le H|t|^\alpha
 \quad(\theta\in\mathbb R,\ |t|\le\pi),
\]
其中 \(H\ge0\)，则当 \(0<\alpha<1\) 时，上述左端至多为
\(C_\alpha H\varepsilon^\alpha\)；当 \(\alpha=1\) 时，至多为
\(CH\varepsilon\log(e/\varepsilon)\)。常数与 \(R,r,\varphi,H\) 无关，\(C_\alpha\) 可依赖 \(\alpha\)。
\end{proposition}
\begin{proof}
写 \(\rho=r/R\)。核的归一化与周期性给出
\[
 u(re^{\mathrm i\theta})-\varphi(Re^{\mathrm i\theta})
 =\frac1{2\pi}\int_{-\pi}^{\pi}p_\rho(t)
 \bigl[\varphi(Re^{\mathrm i(\theta+t)})
       -\varphi(Re^{\mathrm i\theta})\bigr]\,\mathrm dt.
\]
在 \(|t|<\delta\) 上，边值差不超过 \(\omega_\varphi(\delta)\)，核质量至多一；其余部分用边值差至多 \(2\|\varphi\|_\infty\) 及\cref{b3:lem:16-poisson-localization}的远弧估计，得到第一式。连续函数在紧圆周上一致连续，故先使 \(\omega_\varphi(\delta)\) 小，再固定 \(\delta\) 令 \(\varepsilon\downarrow0\)，即得一致收敛。

有 Hölder 条件时，不再先固定一条近弧，而在整个积分中保留差值大小。核的上界给
\[
 |u(re^{\mathrm i\theta})-\varphi(Re^{\mathrm i\theta})|
 \le CH\int_0^\pi
        \frac{\varepsilon t^\alpha}{\varepsilon^2+t^2}\,\mathrm dt.
\]
若 \(0<\alpha<1\)，代入 \(t=\varepsilon s\)，右侧至多为
\[
 CH\varepsilon^\alpha
       \int_0^\infty\frac{s^\alpha}{1+s^2}\,\mathrm ds.
\]
该积分在零附近由 \(s^\alpha\) 控制，在无穷远由 \(s^{\alpha-2}\) 控制，两端均可积。若 \(\alpha=1\)，直接计算得到
\[
 \int_0^\pi\frac{\varepsilon t}{\varepsilon^2+t^2}\,\mathrm dt
 =\frac{\varepsilon}{2}
     \log\left(1+\frac{\pi^2}{\varepsilon^2}\right)
 \le C\varepsilon\log(e/\varepsilon),
\]
于是得到端点估计。
\end{proof}

Lipschitz 情形的对数因子一般不能删去。取单位圆周数据
\(\varphi(e^{\mathrm i t})=|t|\)，其中 \(-\pi\le t\le\pi\)，并作周期延拓。它是到角度零的圆周距离，满足角变量的 Lipschitz 条件。在点 \(1\) 的径向线上，令 \(r=1-\varepsilon\)、\(0<\varepsilon\le1/2\)，则
\[
 (\mathcal P_1\varphi)(r)-\varphi(1)
 =\frac1\pi\int_0^\pi p_r(t)t\,\mathrm dt
 \ge\frac{\varepsilon}{2\pi}
       \int_\varepsilon^1\frac{\mathrm dt}{t}
 =\frac{\varepsilon}{2\pi}\log(1/\varepsilon).
\]
这里使用了核的下界及 \(t\ge\varepsilon\) 时
\(\varepsilon^2+t^2\le2t^2\)。因此对所有 Lipschitz 边值声称误差为 \(O(1-r)\) 并不成立。

\ProseParagraphBreak
固定 \(R>0\) 与 \(z\in B(0,R)\)，把核当作边界上的密度，定义概率测度
\[
 \mathrm d\omega^z(Re^{\mathrm i t})
 =\frac1{2\pi}P_R(z,t)\,\mathrm dt.
\]
它称为圆盘的\term[diaohe-cedu]{调和测度}[harmonic measure]，满足
\[
 (\mathcal P_R\varphi)(z)=\int_{\partial B(0,R)}\varphi\,\mathrm d\omega^z.
\]
中心处它等于归一化弧长；非中心处则对较近的边界赋予较大权重。这里的测度由显式密度定义，稍后的\cref{b3:sec:1604}将在一般正则域上从解算子构造它。

\ProseParagraphBreak
对有界 Borel 数据，同一积分仍有意义，并由\cref{b3:prop:16-poisson-frequencies}给出内部调和函数。改变有限个边界点的取值不会改变积分，因为每个点的弧长测度为零。在任何一个数据连续的边界点 \(\xi\)，积分仍任意方向趋于 \(\varphi(\xi)\)：近 \(\xi\) 的弧上只需该点的连续性，远弧上的边值差由统一有界性控制，而核质量趋于零。全边界连续性的丢失，并不妨碍在其余连续点取到正确边值。

\ProseParagraphBreak
跳跃点处的情况取决于接近方向。下面直接用核的缩放计算这种差别，不需要预先知道某条圆弧的调和测度公式。

\Needspace{6\baselineskip}
\begin{proposition}[跳跃边值的方向极限]\label{b3:prop:16-poisson-jump}
设 \(\varphi:\partial B(0,1)\to\mathbb R\) 为有界 Borel 可测函数，\(u=\mathcal P_1\varphi\)。设在某个 \(\theta_0\in\mathbb R\) 处存在单侧极限
\[
 A_-=\lim_{t\uparrow0}\varphi(e^{\mathrm i(\theta_0+t)}),
 \quad
 A_+=\lim_{t\downarrow0}\varphi(e^{\mathrm i(\theta_0+t)}).
\]
对任意固定 \(\lambda\in\mathbb R\)，令
\(z_\varepsilon=(1-\varepsilon)e^{\mathrm i(\theta_0+\lambda\varepsilon)}\)，其中 \(0<\varepsilon<1/2\)。则
\[
 \lim_{\varepsilon\downarrow0}u(z_\varepsilon)
 =A_-\left(\frac12-\frac{\arctan\lambda}{\pi}\right)
  +A_+\left(\frac12+\frac{\arctan\lambda}{\pi}\right).
\]
特别地，径向极限为 \((A_-+A_+)/2\)。这些极限与
\(\varphi(e^{\mathrm i\theta_0})\) 本身的取值无关。
\end{proposition}
\begin{proof}
旋转后可设 \(\theta_0=0\)。取 \(0<\delta<\pi/2\)，把积分分成 \(|t|<\delta\) 与其余部分。当 \(\varepsilon\) 足够小时，核的中心角 \(\lambda\varepsilon\) 与零的角距离小于 \(\delta/2\)。远弧到中心的圆周角距离至少为 \(\delta/2\)，所以远弧贡献由集中估计趋于零。

在近弧作换元 \(t=\varepsilon s\)。对每个固定实数 \(s\)，有
\[
 \frac{\varepsilon}{2\pi}
 p_{1-\varepsilon}\bigl(\varepsilon(s-\lambda)\bigr)
 \longrightarrow\frac1{\pi\bigl(1+(s-\lambda)^2\bigr)}.
\]
这是把分子写成 \(\varepsilon(2-\varepsilon)\)，分母写成
\(\varepsilon^2+2(1-\varepsilon)(1-\cos(\varepsilon(s-\lambda)))\)
后，使用 \((1-\cos h)/h^2\to1/2\) 得到的。近弧中
\(|\varepsilon(s-\lambda)|<\pi\)，核上界又给出
\[
 \frac{\varepsilon}{2\pi}
 p_{1-\varepsilon}\bigl(\varepsilon(s-\lambda)\bigr)
 \le\frac{C}{1+(s-\lambda)^2}.
\]
将近弧积分的被积函数在 \(|s|\ge\delta/\varepsilon\) 处补零，右侧乘以 \(\|\varphi\|_\infty\) 就是整条实轴上的共同可积控制。对 \(s<0\)，数据趋于 \(A_-\)；对 \(s>0\)，趋于 \(A_+\)。支配收敛因此给出
\[
 \lim_{\varepsilon\downarrow0}u(z_\varepsilon)
 =\frac{A_-}{\pi}\int_{-\infty}^{0}
       \frac{\mathrm ds}{1+(s-\lambda)^2}
  +\frac{A_+}{\pi}\int_{0}^{\infty}
       \frac{\mathrm ds}{1+(s-\lambda)^2}.
\]
两项积分分别为 \(\pi/2-\arctan\lambda\) 与
\(\pi/2+\arctan\lambda\)，得到结论。单个角度的函数值不影响积分。
\end{proof}

当 \(A_-\ne A_+\) 时，改变 \(\lambda\) 就得到不同极限，因而解不能连续延拓到这个边界点。径向得到左右平均，是因为此时核关于该点对称；它并没有恢复一份连续边值。圆弧指示函数的端点属于这种情形，具体圆弧和实轴区间的权重将在\cref{b3:sec:1604}中计算。

\ProseParagraphBreak
虽然有限个点的指定数值不影响 Poisson 积分，但从“两个有界调和函数在其他边界点同值”推出内部相同，仍须一个比较论证。直接使用跳跃的指示数据会在弧端点留下未控制的极限；采用连续的短弧上界，可以在每个边界点应用已证最大值原理。

\Needspace{6\baselineskip}
\begin{lemma}[有限边界例外集的比较]\label{b3:lem:16-finite-boundary-exceptions}
设 \(v:B(0,1)\to\mathbb R\) 调和且有上界，\(E\subset\partial B(0,1)\) 为有限集。若
\[
 \limsup_{B(0,1)\ni z\to\xi}v(z)\le0
 \quad(\xi\in\partial B(0,1)\setminus E),
\]
则 \(v\le0\) 于圆盘。特别地，若 \(v\) 有界且在除 \(E\) 外的每个边界点均趋于零，则 \(v=0\)。
\end{lemma}
\begin{proof}
取 \(M\ge0\) 使 \(v\le M\)。给定 \(\eta>0\)，用有限条开弧覆盖 \(E\)，使总角长小于 \(\eta\)。在每条弧中取一个连续函数，在相应例外点的小邻弧上等于一，在大弧外等于零，并在中间线性过渡。取这些函数的最大值，得到连续圆周函数 \(\chi_\eta\)，满足
\[
 0\le\chi_\eta\le1,\quad
 \chi_\eta=1\text{ 于 }E,\quad
 \int_0^{2\pi}\chi_\eta(e^{\mathrm i t})\,\mathrm dt<\eta.
\]
当 \(E\) 为空时直接取 \(\chi_\eta=0\)。令
\(h_\eta=\mathcal P_1\chi_\eta\)，则 \(h_\eta\ge0\)，且连续取边值 \(\chi_\eta\)。

在 \(E\) 以外，\(\limsup(v-Mh_\eta)\le-M\chi_\eta\le0\)；在 \(E\) 上，由 \(v\le M\) 与 \(h_\eta\to1\)，同样得到上极限不超过零。因此调和函数 \(v-Mh_\eta\) 在有界圆盘的每个边界点均满足上极限条件，\cref{b3:cor:15-boundary-limsup}给出
\(v\le Mh_\eta\)。固定 \(z\) 后，核满足
\[
 0\le h_\eta(z)
 \le\frac{1+|z|}{1-|z|}\,\frac{\eta}{2\pi}.
\]
令 \(\eta\downarrow0\) 得 \(v(z)\le0\)。最后若 \(v\) 有界且边界极限为零，对 \(v\) 与 \(-v\) 分别应用前一结论，得到 \(v=0\)。
\end{proof}

上半平面给出这一引理的一个用途：将它映到圆盘后，无穷远只对应一个边界点。实轴上的积分核可以先直接构造。对 \(z=x+\mathrm i y\)、\(y>0\)，函数
\[
 \frac1\pi\operatorname{Im}\frac1{t-z}
 =\frac1\pi\frac{y}{(x-t)^2+y^2}
\]
关于 \(z\) 调和且为正。代入 \(t=x+ys\)，它对 \(t\) 的积分成为
\(\pi^{-1}\int_{\mathbb R}(1+s^2)^{-1}\,\mathrm ds=1\)，所以也能用来延拓有界数据。

\Needspace{6\baselineskip}
\begin{theorem}[上半平面的 Poisson 积分]\label{b3:thm:16-half-plane-poisson}
设 \(\psi:\mathbb R\to\mathbb R\) 为有界 Borel 可测函数，定义
\[
 U(x,y)=\frac1\pi\int_{\mathbb R}
       \frac{y\,\psi(t)}{(x-t)^2+y^2}\,\mathrm dt,
 \quad x\in\mathbb R,\ y>0.
\]
则 \(U\) 在上半平面调和，且 \(|U|\le\|\psi\|_\infty\)。若 \(\psi\) 在 \(x_0\in\mathbb R\) 连续，则
\[
 U(x,y)\longrightarrow\psi(x_0)
 \quad\text{当 }(x,y)\to(x_0,0),\ y>0.
\]
若 \(\psi\) 在整条实轴连续，则 \(U\) 是唯一有界、在上半平面调和并连续取实轴边值 \(\psi\) 的函数。
\end{theorem}
\begin{proof}
核为正且积分为一，因此积分绝对收敛，并给出 \(|U|\le\|\psi\|_\infty\)。固定上半平面中的紧集 \(K\)。其上的 \(x,y\) 有界，且 \(y\) 有正下界。核及其至多二阶的 \(x,y\) 偏导数，在 \(t\) 的有界区间上统一有界，在 \(|t|\) 足够大时由 \(C_K/(1+t^2)\) 控制；这一点可由核的有理表达直接求导得到。乘以有界的 \(\psi\) 后仍有共同可积控制，故可在积分号下求二阶导数。核本身是关于 \(z\) 的全纯函数的虚部，Laplace 算子为零，因而 \(\Delta U=0\)。

令 \(M=\|\psi\|_\infty\)。固定连续点 \(x_0\) 与 \(\gamma>0\)，取 \(\delta>0\)，使 \(|t-x_0|<\delta\) 时
\(|\psi(t)-\psi(x_0)|<\gamma\)。若 \(|x-x_0|<\delta/2\)，则远部满足
\(|t-x|\ge|t-x_0|/2\)，所以
\[
 \frac1\pi\int_{|t-x_0|\ge\delta}
       \frac{y}{(x-t)^2+y^2}\,\mathrm dt
 \le\frac{4y}{\pi}\int_{|s|\ge\delta}\frac{\mathrm ds}{s^2}
 =\frac{8y}{\pi\delta}.
\]
将 \(U(x,y)-\psi(x_0)\) 按这两部分拆开，近部用核的总质量一，远部用差值至多 \(2M\)，得到
\[
 |U(x,y)-\psi(x_0)|
 \le\gamma+\frac{16My}{\pi\delta}.
\]
先令 \((x,y)\to(x_0,0)\)，再令 \(\gamma\downarrow0\)，便得到任意方向的边界极限。

现在假设 \(\psi\) 在整条实轴连续，并取另一有界调和解 \(V\)。差 \(W=V-U\) 有界且连续取零实轴边值。映射
\[
 F(\zeta)=\mathrm i\,\frac{1+\zeta}{1-\zeta},
 \quad |\zeta|<1,
\]
将单位圆盘双全纯地映到上半平面，因为
\[
 \operatorname{Im}F(\zeta)=\frac{1-|\zeta|^2}{|1-\zeta|^2}>0,
 \quad
 F^{-1}(z)=\frac{z-\mathrm i}{z+\mathrm i}.
\]
除圆周点 \(1\) 外，\(F\) 连续取有限实轴值。链式法则与 Cauchy--Riemann 方程给出
\(\Delta(W\circ F)=|F'|^2\cdot(\Delta W)\circ F=0\)，故 \(W\circ F\) 是有界调和函数，并在除 \(1\) 外的所有圆周点趋于零。\cref{b3:lem:16-finite-boundary-exceptions}遂给 \(W\circ F=0\)，从而 \(V=U\)。证明没有要求在无穷远预先存在一个边值。

\end{proof}

有界性在唯一性证明中控制了映到圆周点 \(1\) 的无穷远部分。若放弃这项条件，零函数与 \(U(x,y)=y\) 都连续取零实轴边值，却在内部不同。同一证明还允许在实轴上忽略有限个边界点：若两个有界调和函数在其余每个实点的差都趋于零，拉回后的例外集只是这些点的像再加上 \(1\)，仍为有限集，故两函数相等。

\ProseParagraphBreak
实轴非紧，各连续点的边界收敛本身不保证全实轴上的一致收敛。

\ProseParagraphBreak
一些积分值可以通过唯一性计算，而不必直接积分。数据
\(\psi(t)=1/(1+t^2)\) 的有界调和延拓是
\[
 U(x,y)=\frac{1+y}{x^2+(1+y)^2}
       =-\operatorname{Im}\frac1{x+\mathrm i y+\mathrm i}.
\]
右端的极点位于下半平面，在 \(y\ge0\) 上介于零与一之间，并连续取指定边值，故由唯一性就是 Poisson 积分。对数据 \(\cos(\kappa t)\)，其中 \(\kappa\in\mathbb R\)，函数
\[
 U(x,y)=e^{-|\kappa|y}\cos(\kappa x)
\]
同样直接满足调和方程、有界性和边界条件。频率越高，向上半平面内部的衰减越快；这与圆盘中 \((r/R)^{|n|}\) 的衰减是两个可明确计算的模型。

% !TeX root = ../../Book3.tex
\section{Perron 方法}
\label{b3:sec:1602}

一般区域没有可直接写出的核。Perron 方法从所有低于边界数据的次调和函数出发；圆盘上的调和替换迫使它们的上包络满足平均等式。

\subsection{次函数}
\label{b3:subsec:160201}

称连续次调和函数 \(v:D\to\mathbb R\) 为边值 \(\varphi\) 的\term[perron-cixanshu]{Perron 次函数}[Perron subfunction]，若对每个 \(\xi\in\partial D\) 有 \(\limsup_{z\to\xi}v(z)\le\varphi(\xi)\)。本章只采用连续的次函数；这个族已足以完成连续边值的构造。

常数 \(m=\min\varphi\) 是次函数。每个次函数由次调和最大值原理满足 \(v\le M=\max\varphi\)，所以取上确界不会在内部发散。连续性允许在任何内部闭圆盘上作 Poisson 替换。

\subsection{Perron 族}
\label{b3:subsec:160202}

把全部次函数组成的族记为 \(\mathcal P_\varphi\)，上包络记为 \(H_\varphi(z)=\sup_{v\in\mathcal P_\varphi}v(z)\)。有限最大值仍在族中。

若 \(B\Subset D\)，把 \(v\) 在 \(B\) 内换成同边值的 Poisson 解，外部保持 \(v\)，记为 \(v^B\)。比较原理给 \(v^B\ge v\)。该函数在圆周连续，且仍次调和；下文通过先比较原次函数、再比较替换部分给出跨圆周的完整核查。替换只影响内部紧集，所以原边界限制不变。这个封闭性是上包络证明的核心。

\subsection{Perron 解}
\label{b3:subsec:160203}

\begin{theorem}[Perron]\label{b3:thm:16-perron}
有界区域 \(D\) 与连续边值 \(\varphi\) 的包络 \(H_\varphi\) 在 \(D\) 调和，并满足 \(m\le H_\varphi\le M\)。
\end{theorem}
\begin{proof}
固定内部圆盘 \(B\) 及 \(a\in B\)。选 \(v_n\) 使 \(v_n(a)\to H_\varphi(a)\)，逐次取有限最大值，得到递增序列，再在 \(B\) 调和替换。所得 \(h_n\) 递增、有界，在 \(B\) 局部一致收敛到调和函数 \(h\)，且 \(h(a)=H_\varphi(a)\)、\(h\le H_\varphi\)。
任取另一项 \(v\in\mathcal P_\varphi\)，将 \(\max(v,h_n)\)（在全域使用对应替换后的函数）再提升，得到调和函数 \(k_n\ge h_n,v\)。它们同样递增有界，极限 \(k\) 调和。中心满足 \(h(a)\le k(a)\le H_\varphi(a)=h(a)\)，所以非负调和函数 \(k-h\) 在内部为零，由强最小值原理恒零。于是 \(v\le h\) 于 \(B\)，取全族上确界得 \(H_\varphi\le h\)。因此包络在每个内部圆盘调和。
\end{proof}
证明没有先断言任意上包络连续；连续性由局部调和识别产生。这一步避免了把上确界与上半连续正则化混同。

\subsection{唯一性}
\label{b3:subsec:160204}

若一个经典解 \(u\) 已存在，则它本身是次函数，每个次函数又由比较原理不超过 \(u\)，故 \(H_\varphi=u\)。因此 Perron 构造不产生第二种不同的经典解。

没有边界正则性时，\(H_\varphi\) 仍是内部调和函数，却可能忽略某些指定的边界值。去心圆盘上，外圆取零、孤立边界点原点取一，得到 Perron 解恒零：有界调和函数的孤立奇点可去，延拓后外圆边值迫使零函数。该例将在下一节同时说明障碍函数为何必要。


\ProseParagraphBreak

调和替换在圆周上的粘合可以作如下严格核查。设 \(w=v^B\)，并在测试圆盘 \(Q\Subset D\) 取调和函数 \(h\)，满足 \(w\le h\) 于 \(\partial Q\)。由于 \(v\le w\)，先由 \(v\) 的次调和比较得到 \(v\le h\) 于 \(Q\)。在 \(Q\cap B\) 中，\(w-h\) 调和；其边界分成 \(\partial Q\cap\overline B\) 与 \(Q\cap\partial B\)，前者由假设非正，后者因 \(w=v\le h\) 非正。最大值原理便给 \(w\le h\) 于交集。其余部分 \(w=v\le h\)，所以整个 \(Q\) 比较成立。这里先比较原次函数，再比较替换部分，避免在交界处循环使用尚未证明的次调和性。

\ProseParagraphBreak

Perron 证明中选择序列只需要在一个点逼近包络，而不要求全域同时选出最优函数。具体可令 \(v_n(a)>H_\varphi(a)-1/n\)，将 \(m,v_1,\ldots,v_n\) 的最大值记为 \(s_n\)，再替换得 \(w_n=s_n^B\)。序列 \(w_n\) 在全域递增，且属于同一 Perron 族；其圆盘限制 \(h_n\) 调和。由于 \(h_n(a)\ge s_n(a)\) 且 \(h_n\le H_\varphi\)，中心极限正是包络。对任意竞争者 \(v\)，使用 \(\max(v,w_n)\) 而非只定义在 \(B\) 中的 \(\max(v,h_n)\)，就得到合法的全域次函数。这些域的区别保证第二轮替换仍保留边界约束。

\ProseParagraphBreak

取上确界并不通常与取极限交换。Perron 证明采用的事实是：对每个固定竞争者 \(v\)，新调和极限 \(k\) 在圆盘上支配它，同时 \(k(a)=h(a)\)；非负调和函数的强最小值原理强迫 \(k=h\)。于是一个最初只在中心被优化的序列，最终在整个圆盘支配全部竞争者。这个步骤使用线性的调和方程和正性，才产生包络的调和性；只看有限最大值的次调和封闭性还不足以完成证明。

\ProseParagraphBreak

去心圆盘例子的包络恒零也能在不预先调用可去奇点定理的情况下证明。对任意次函数 \(v\)，外圆边界上限零，内边界点上限至多一，故 \(v\le1\)。给定 \(\varepsilon>0\)，在足够小的内圆 \(|z|=r\) 上 \(v\le1+\varepsilon\)，环域上的调和函数
\[
q_r(z)=(1+\varepsilon)\frac{\log|z|}{\log r}
\]
外圆值零、内圆值 \(1+\varepsilon\)。比较给 \(v(z)\le q_r(z)\)；固定 \(z\ne0\) 让 \(r\downarrow0\)，得到 \(v(z)\le0\)。常数零又是合法次函数，所以包络恰为零。孤立边值一在内部没有影响，这一计算完全在当前已证明的环域比较范围内进行。

\ProseParagraphBreak

边界数据的单调性不需要域正则：若 \(\varphi\le\psi\)，则 \(\mathcal P_\varphi\subseteq\mathcal P_\psi\)，故 \(H_\varphi\le H_\psi\)。加常数与乘正数也分别对应次函数族的平移与缩放，所以 \(H_{\varphi+c}=H_\varphi+c\)、\(H_{a\varphi}=aH_\varphi\) 对 \(a>0\) 成立。两者推出
\[
\|H_\varphi-H_\psi\|_{L^\infty(D)}\le\|\varphi-\psi\|_{L^\infty(\partial D)}.
\]
一般加法的线性不能只由族包含轻率推出；本章在正则域用经典解的唯一性证明完整线性，已经足够用于测度表示。

\ProseParagraphBreak

圆环提供检验一般存在理论的另一组显式解。设 \(r_0<|z|<R\)，令 \(x=\log(|z|/r_0)\)、\(L=\log(R/r_0)\)，则 \(0<x<L\)。极坐标下 Laplace 方程乘以 \(|z|^2\) 后成为 \(u_{xx}+u_{\theta\theta}=0\)。对内、外边界各一个 Fourier 模式，解分别为
\[
 \frac{\sinh(n(L-x))}{\sinh(nL)}e^{\mathrm{i}n\theta},\qquad
 \frac{\sinh(nx)}{\sinh(nL)}e^{\mathrm{i}n\theta}\quad(n\ge1).
\]
它们在相应边界取一倍模式，在另一边界取零；对负频率取共轭。零频率则为 \(1-x/L\) 与 \(x/L\)。每一式都可直接求导检验，唯一性保证它与 Perron 解相同。
\ProseParagraphBreak
若内、外边界为三角多项式，有限叠加便得到解。对一般连续周期数据，可由非负的 Fej\'er 核逼近：
\[
 F_N(t)=\frac1N\left|\sum_{j=0}^{N-1}e^{\mathrm{i}jt}\right|^2,
 \qquad \frac1{2\pi}\int_{-\pi}^{\pi}F_N(t)\,\mathrm{d}t=1.
\]
在 \(|t|\ge\delta\) 上，有限几何和给 \(F_N(t)\le C_\delta/N\)；在小弧上用数据的一致连续性。因此卷积 \(F_N*\varphi\) 一致趋于 \(\varphi\)，每次卷积都是三角多项式。对这些数据的圆环解，最大值原理使内部一致差不超过两条边界数据的最大差，故在闭圆环上一致收敛。极限在内部调和，因为任意小圆盘上的平均值性质可通过一致极限；边值也正确。这给出圆环上存在性的第二个完整构造，并显示对数项恰来自角向零频率。

% !TeX root = ../../Book3.tex
\section{边界正则性}
\label{b3:sec:1603}

Perron 包络已经在内部调和，但这还没有说明它是否取到每一个指定的边界值。本节把问题集中到一个边界点：若能找到在这个点趋零、远离这个点却保持正下界的超调和函数，就能用它同时控制所有连续边值的误差。这样的函数称为障碍；它将局部几何信息转化为边界比较。

\subsection{障碍函数}
\label{b3:subsec:160301}

在边界点 \(\xi\) 附近，我们希望一个正函数只在趋向 \(\xi\) 时变小。仅要求它在域内逐点为正并不够：沿另一个边界点趋近时，它仍可能趋于零。定义中的环带下界排除这种情形，也为稍后的截断和拼接留出余量。

\Needspace{6\baselineskip}
\begin{definition}\label{b3:def:16-barrier}
设 \(D\subseteq\mathbb C\) 为有界区域，\(\xi\in\partial D\)。若存在 \(r_0>0\) 及连续函数
\[
 b:D\cap B(\xi,r_0)\longrightarrow(0,\infty),
\]
使 \(b\) 在该开集的每个连通分支超调和，且满足
\[
 \lim_{z\to\xi,\ z\in D}b(z)=0,
 \quad
 \inf_{\substack{z\in D\\\delta\le|z-\xi|\le r'}}b(z)>0
 \quad(0<\delta<r'<r_0),
\]
则称 \(b\) 为 \(\xi\) 处的连续正超调和\term[zhangai-hanshu]{障碍函数}[barrier]，简称局部障碍。环带与 \(D\) 的交集为空时，下界条件不另加限制。
\end{definition}
这里的超调和是指 \(-b\) 次调和。函数 \(b\) 不必连续延拓到整个局部边界；在一个可能不闭的集合 \(D\cap\{|z-\xi|=s\}\) 上逐点为正，也不能单凭连续性推出统一正下界。因而定义直接要求内部环带上的统一下界，而不以人工圆周的紧性代替它。

\Needspace{6\baselineskip}
\begin{lemma}[局部障碍的全域化]\label{b3:lem:16-global-barrier}
设 \(D\subseteq\mathbb C\) 为有界区域，\(\xi\in\partial D\)，且 \(\xi\) 有连续正超调和局部障碍 \(b\)。则存在有界、连续、正的超调和函数 \(B\) 于 \(D\)，在 \(\xi\) 的某个域内邻域等于 \(b\)，并满足
\[
 \lim_{z\to\xi,\ z\in D}B(z)=0,
 \quad
 \inf_{D\setminus B(\xi,\delta)}B>0
 \quad(\delta>0).
\]
最后一式只要求于非空的集合。
\end{lemma}
\begin{proof}
设 \(b\) 定义在 \(D\cap B(\xi,r_0)\)。取 \(0<s<r_0/2\)。由环带条件，可选 \(c>0\)，使
\[
 b(z)\ge2c
 \quad\left(z\in D,\ s\le|z-\xi|\le2s\right).
\]
若该环带中没有域内点，任取 \(c>0\) 即可。定义
\[
 B(z)=
 \begin{cases}
 \min\{b(z),c\},&|z-\xi|<3s/2,\\
 c,&|z-\xi|\ge3s/2.
 \end{cases}
\]
在整个域内环带 \(s<|z-\xi|<2s\) 上，两种表达都等于 \(c\)。因此拼接处已有一个共同的常值邻域，而不是仅在人工圆周上取值相同。函数 \(B\) 连续且处处为正。有限个超调和函数的最小值仍超调和，这是\cref{b3:prop:15-subharmonic-operations}对其负号取最大值的结论；再用\cref{b3:prop:15-local-comparison}的局部性质，得到全域超调和性。

当 \(z\to\xi\) 时 \(b(z)\to0\)，所以在充分小的邻域有 \(b<c\)，从而 \(B=b\)。对远处的下界，若 \(\delta\ge s\)，则 \(D\setminus B(\xi,\delta)\) 上 \(B=c\)。若 \(0<\delta<s\)，原障碍在 \(\delta\le|z-\xi|\le s\) 上有正下界；将它与 \(c\) 取较小值，便控制了整个 \(D\setminus B(\xi,\delta)\)。
\end{proof}
圆盘给出一个可以直接核验的全域障碍。设 \(D=B(a,R)\)、\(|\xi-a|=R\)，令
\[
 b(z)=1-\operatorname{Re}\frac{(z-a)\overline{(\xi-a)}}{R^2}.
\]
它是调和函数，而且
\[
 2R^2b(z)=R^2-|z-a|^2+|z-\xi|^2\ge|z-\xi|^2.
\]
在圆盘内部第一项严格为正，因此 \(b>0\)；趋于 \(\xi\) 时 \(b\to0\)，而 \(|z-\xi|\ge\delta\) 时 \(b(z)\ge\delta^2/(2R^2)\)。这个恒等式把“远离指定边界点”需要的统一下界直接算了出来。

\subsection{正则边界点}
\label{b3:subsec:160302}

边界正则性首先是解对数据的性质，不依赖我们是否已经找到某种几何障碍。由\cref{b3:thm:16-perron}，有界区域上每份连续实边值都有一个内部调和的 Perron 包络；定义只问这个包络在一个指定边界点的极限。

\Needspace{6\baselineskip}
\begin{definition}\label{b3:def:16-regular-point}
设 \(D\subseteq\mathbb C\) 为有界区域，\(\xi\in\partial D\)。对每个 \(\varphi\in C(\partial D;\mathbb R)\)，令 \(H_\varphi\) 为由连续次调和 Perron 次函数组成的上包络。若
\[
 \lim_{z\to\xi,\ z\in D}H_\varphi(z)=\varphi(\xi)
 \quad\left(\varphi\in C(\partial D;\mathbb R)\right),
\]
则称 \(\xi\) 为 \(D\) 的\term[zhengze-bianjiedian]{正则边界点}[regular boundary point]。若 \(\partial D\) 的每个点都正则，称 \(D\) 为\term[zhengze-yu]{正则域}[regular domain]。
\end{definition}
量词“每个连续边值”不能省略。常数数据在任何边界点都会取到正确极限，不能用一份常数数据检验正则性。下面的等价命题既给充分条件，也说明局部障碍确实描述了这一性质。

\Needspace{6\baselineskip}
\begin{proposition}[障碍与边界正则性]\label{b3:prop:16-barrier-regular}
设 \(D\subseteq\mathbb C\) 为有界区域，\(\xi\in\partial D\)。则 \(\xi\) 为正则边界点，当且仅当它具有连续正超调和局部障碍。若 \(\xi\) 正则，还可以取一个在整个 \(D\) 调和的正障碍，且它在每个 \(D\setminus B(\xi,\delta)\) 上有统一正下界。
\end{proposition}
\begin{proof}
先设局部障碍存在，由\cref{b3:lem:16-global-barrier}取全域障碍 \(B\)。固定 \(\varphi\in C(\partial D;\mathbb R)\) 和 \(\varepsilon>0\)。选 \(\delta>0\)，使
\[
 |\varphi(\eta)-\varphi(\xi)|<\varepsilon
 \quad\left(\eta\in\partial D,\ |\eta-\xi|<\delta\right).
\]
设 \(M=\max_{\partial D}|\varphi-\varphi(\xi)|\)。全域障碍在 \(D\setminus B(\xi,\delta/2)\) 上有正下界 \(c_\delta\)，取 \(A\ge0\) 使 \(Ac_\delta\ge M\)。若该集合为空，所有边界点都在上述邻域中，取 \(A=0\) 即可。

函数
\[
 L(z)=\varphi(\xi)-\varepsilon-AB(z),
 \quad U(z)=\varphi(\xi)+\varepsilon+AB(z)
\]
分别连续次调和、连续超调和。在距离 \(\xi\) 小于 \(\delta\) 的边界点，正性和数据连续性给 \(\limsup L\le\varphi\)、\(\liminf U\ge\varphi\)。在其余边界点，域内的趋近点最终满足 \(|z-\xi|\ge\delta/2\)，统一下界 \(c_\delta\) 及 \(Ac_\delta\ge M\) 给出相同结论。因此 \(L\) 属于 Perron 次函数族。对任意该族中的 \(v\)，函数 \(v-U\) 次调和，边界上极限非正；\cref{b3:prop:15-subharmonic-maximum}在有界域上给 \(v\le U\)。取上确界得到
\[
 \varphi(\xi)-\varepsilon-AB(z)
 \le H_\varphi(z)\le
 \varphi(\xi)+\varepsilon+AB(z).
\]
先令 \(z\to\xi\)，利用 \(B(z)\to0\)，再令 \(\varepsilon\downarrow0\)，便得正则性。

反过来，只假定 \(\xi\) 正则。取连续边值 \(\psi(\eta)=|\eta-\xi|^2\)，并令 \(B=H_\psi\)。函数 \(v(z)=|z-\xi|^2\) 连续次调和，且在每个边界点取到 \(\psi\)，所以它本身就是合法次函数。因此
\[
 0<|z-\xi|^2\le B(z)\le\max_{\partial D}|\eta-\xi|^2
 \quad(z\in D).
\]
Perron 定理保证 \(B\) 调和；\(\xi\) 的正则性又给 \(B(z)\to\psi(\xi)=0\)。而 \(|z-\xi|\ge\delta\) 时 \(B(z)\ge\delta^2\)，故 \(B\) 是所需全域正调和障碍。这个构造没有使用其他边界点的正则性。
\end{proof}
因此只需检验一份与指定点相关的数据：若边值 \(|\eta-\xi|^2\) 的 Perron 包络在 \(\xi\) 趋于零，上面的构造就给出障碍，继而迫使所有连续数据在 \(\xi\) 取到正确极限。正则性也因而是局部的：局部障碍在相同的域内边界几何上可以直接使用，全域部分由截断补齐。

\ProseParagraphBreak

在正则域上，\(H_\varphi\) 在每个边界点取连续数据的值，连同内部的调和性便给 \(H_\varphi\in C(\overline D)\)。经典解的唯一性来自最大值原理。因此，Perron 法的内部存在性与障碍的边界到达性质合在一起，完成了正则域上连续 Dirichlet 问题的存在唯一性。

\subsection{外锥条件}
\label{b3:subsec:160303}

外锥属于补集。具体地，若 \(\mathbb C\setminus D\) 在 \(\xi\) 附近包含一个以 \(\xi\) 为顶点、角宽为正的实心扇形，域的局部部分便落在余下角宽 \(\beta<2\pi\) 的扇形中。\(\beta\) 表示容纳域的角宽；它不是补集外锥的角宽。选定旋转角 \(\theta_0\) 后，把域内点写成
\[
 z=\xi+r e^{\mathrm{i}(\theta_0+\theta)},
 \quad 0<r<r_0,\quad |\theta|<\beta/2.
\]
因为 \(\beta<2\pi\)，这里可以选取一个连续辐角分支。域只须包含在这个扇形中，不必等于扇形，也不要求原边界是直线。

\ProseParagraphBreak

极坐标下的调和方程提示取 \(r^\alpha\cos(\alpha\theta)\)：径向二阶项和角向二阶项恰好相消。但障碍还要求整个允许角范围内保持正下界，所以 \(\alpha\) 的选择不仅由调和方程决定。

\Needspace{6\baselineskip}
\begin{proposition}[外锥条件与正则性]\label{b3:prop:16-exterior-cone-regular}
设 \(D\subseteq\mathbb C\) 为有界区域，\(\xi\in\partial D\)。若存在 \(r_0>0\)、\(\theta_0\in\mathbb R\) 及 \(0<\beta<2\pi\)，使
\[
 D\cap B(\xi,r_0)
 \subseteq\{\,\xi+r e^{\mathrm{i}(\theta_0+\theta)}
       \mid 0<r<r_0,\ |\theta|<\beta/2\,\},
\]
则 \(\xi\) 正则。对任意 \(0<\alpha<\pi/\beta\)，函数
\[
 b(z)=r^\alpha\cos(\alpha\theta)
\]
都是 \(\xi\) 处的连续正调和局部障碍。
\end{proposition}
\begin{proof}
在上述辐角分支上，\(b\) 是 \(\bigl(e^{-\mathrm{i}\theta_0}(z-\xi)\bigr)^\alpha\) 的实部，故调和。令 \(c_\alpha=\cos(\alpha\beta/2)\)。严格不等式 \(\alpha\beta/2<\pi/2\) 给 \(c_\alpha>0\)，并且
\[
 c_\alpha r^\alpha\le b(z)\le r^\alpha.
\]
上界给 \(z\to\xi\) 时 \(b(z)\to0\)，下界给每个 \(\delta\le r\le r'<r_0\) 上的统一正下界 \(c_\alpha\delta^\alpha\)。因此 \(b\) 满足局部障碍的全部条件，由\cref{b3:prop:16-barrier-regular}得到正则性。
\end{proof}
\begin{figure}[htbp]
\centering
\begin{tikzpicture}[>=stealth,font=\small,scale=1.12]
 \fill[gray!30] (0,0)--(120:2) arc (120:240:2)--cycle;
 \draw[dashed] (0,0) circle (2);
 \draw (120:2)--(0,0)--(240:2);
 \draw[dotted,->] (0,0)--(2.35,0) node[right] {\(\theta=0\)};
 \draw[->] (-120:0.65) arc (-120:120:0.65);
 \node[right] at (0.68,0.22) {\(\beta\)};
 \draw[thick] (-0.38,1.6)..controls (-0.4,0.7) and (-0.1,0.35)..(0,0)
   ..controls (0.3,-0.65) and (0.25,-1.3)..(0.8,-1.8);
 \node at (1.18,0.8) {\(D\)};
 \node[align=center] at (-1.3,0) {补集\\外锥};
 \fill (0,0) circle (0.025) node[below left] {\(\xi\)};
 \node[above right] at (45:2) {人工圆周 \(r=r_0\)};
\end{tikzpicture}
\caption[外锥与正障碍的扇形]{阴影锥包含于补集，域在局部只包含于余下角宽 \(\beta<2\pi\) 的扇形，不要求等于该扇形。粗曲线仅示意可能的域边界。选择 \(\alpha\beta/2<\pi/2\) 后，整个允许角范围内都有 \(\cos(\alpha\theta)\ge\cos(\alpha\beta/2)>0\)，故人工圆周上障碍具有统一正下界。}
\label{b3:fig:16-exterior-cone-barrier}
\end{figure}

\cref{b3:fig:16-exterior-cone-barrier}中的粗曲线示意原域边界，虚线圆周则是局部化人为加入的边界。余弦的统一正下界在整个环带有效，因而可以使用\cref{b3:lem:16-global-barrier}；不必假定人工圆周与域的交集是紧的。若取临界值 \(\alpha=\pi/\beta\)，余弦在两条侧边趋于零，环带上的正下界就可能消失。这个临界函数仍然调和，却未必符合当前障碍的定义。

\ProseParagraphBreak

例如局部 Lipschitz 边界可旋转成图像 \(y=g(x)\)，其中 \(g(0)=0\)、\(|g(x)-g(x')|\le L|x-x'|\)，域局部位于 \(y>g(x)\)。由于 \(g(x)\ge-L|x|\)，下方的锥 \(y<-(L+1)|x|\) 包含于补集。于是 Lipschitz 边界点满足外锥条件；光滑边界及角宽小于 \(2\pi\) 的多边形角点也属于这一范围。再入角对应 \(\beta>\pi\)，并不妨碍正则性。

\Needspace{6\baselineskip}
\begin{proposition}[外锥下的局部 Hölder 估计]\label{b3:prop:16-corner-holder}
设 \(D\subseteq\mathbb C\) 为有界区域，\(\xi\in\partial D\)。假设某个 \(r_0>0\) 下，\(D\cap B(\xi,r_0)\) 包含于以 \(\xi\) 为顶点、角宽 \(0<\beta<2\pi\) 的扇形。设 \(\varphi\in C(\partial D;\mathbb R)\)，\(H_\varphi\) 为连续次函数 Perron 包络。固定 \(0<\alpha<\pi/\beta\)。若存在 \(H\ge0\) 和 \(0<\gamma\le1\)，使
\[
 |\varphi(\eta)-\varphi(\xi)|\le H|\eta-\xi|^\gamma
 \quad\left(\eta\in\partial D\cap B(\xi,r_0)\right),
\]
则存在 \(r_1>0\)、\(C\ge0\)，使
\[
 |H_\varphi(z)-\varphi(\xi)|
 \le C|z-\xi|^{\min(\alpha,\gamma)}
 \quad\left(z\in D\cap B(\xi,r_1)\right).
\]
常数可依赖局部扇形、所选指数、半径及边值数据。
\end{proposition}
\begin{proof}
取 \(b=r^\alpha\cos(\alpha\theta)\)，令
\[
 q=\min(1,\gamma/\alpha),\quad p=\alpha q=\min(\alpha,\gamma),
 \quad c_\alpha=\cos(\alpha\beta/2)>0.
\]
在域内部 \(b>0\)、\(\Delta b=0\)，链式法则给
\[
 \Delta(b^q)=q(q-1)b^{q-2}|\nabla b|^2\le0,
 \quad c_\alpha^q r^p\le b^q\le r^p.
\]
故 \(b^q\) 是连续正超调和局部障碍。由\cref{b3:lem:16-global-barrier}取其全域化 \(B\)，并选 \(0<\rho<r_0\)，使 \(B=b^q\) 于 \(D\cap B(\xi,\rho)\)。

若边界点 \(\eta\) 满足 \(0<|\eta-\xi|<\rho/2\)，则
\[
 \liminf_{z\to\eta,\ z\in D}B(z)
 \ge c_\alpha^q|\eta-\xi|^p,
 \quad
 |\varphi(\eta)-\varphi(\xi)|
 \le H(\rho/2)^{\gamma-p}|\eta-\xi|^p.
\]
因此可用同一个倍数 \(A B\) 支配这部分边界误差。在其余边界点，域内趋近点最终离 \(\xi\) 至少 \(\rho/4\)，而 \(B\) 在这个范围有统一正下界；增大 \(A\) 即可支配其余有界误差。在 \(\xi\) 本身，\(B\to0\)，误差也为零。

于是 \(\varphi(\xi)-AB\) 是 Perron 次函数，\(\varphi(\xi)+AB\) 是支配所有次函数的超调和上界。和障碍正则性证明中的比较相同，得到
\[
 |H_\varphi(z)-\varphi(\xi)|\le AB(z)\le A|z-\xi|^p
 \quad\left(z\in D\cap B(\xi,\rho)\right).
\]
取 \(C=A\)、\(r_1=\rho\) 即得结论。
\end{proof}
当边值指数 \(\gamma\) 小于障碍指数 \(\alpha\) 时，取 \(q=\gamma/\alpha<1\) 会使障碍在顶点附近下降得慢一些，足以控制较粗糙的边值误差；凹幂产生的超调和性保留了比较方向。仅有某个障碍存在，并不保证它具有这样的幂次上下界，所以定性的边界到达不能自动加强为固定指数的 Hölder 估计。

\ProseParagraphBreak

再入角确实可能降低解的边界光滑性。在截断扇形 \(0<r<1\)、\(|\theta|<\beta/2\) 上，取 \(\pi<\beta<2\pi\)，则
\[
 u(r,\theta)=r^{\pi/\beta}\cos(\pi\theta/\beta)
\]
调和、连续到闭域，并在两条侧边及顶点取零。在中轴上，\(u(r,0)/r=r^{\pi/\beta-1}\to\infty\)，所以它在顶点不是 Lipschitz 函数，尽管靠近顶点的边界数据恒为零。正则性保证连续到达，并未保证一阶边界正则性。

\subsection{障碍算例与 Wiener 判据概览}
\label{b3:subsec:160304}

外接圆提供另一种直接可算的障碍，但须限制接触集。若同一个外圆沿一段边界接触，下面的对数会同时在这些点趋于零，不能单独识别指定点。

\Needspace{6\baselineskip}
\begin{proposition}[唯一接触的外球障碍]\label{b3:prop:16-exterior-ball-barrier}
设 \(D\subseteq\mathbb C\) 为有界区域，\(\xi\in\partial D\)。若存在 \(a\in\mathbb C\)、\(R>0\)、\(r_0>0\)，使
\[
 \overline{B(a,R)}\subseteq\mathbb C\setminus D,
 \quad |\xi-a|=R,
 \quad
 \overline D\cap\overline{B(\xi,r_0)}\cap\overline{B(a,R)}=\{\xi\},
\]
则 \(b(z)=\log(|z-a|/R)\) 在 \(D\cap B(\xi,r_0)\) 是连续正调和障碍，因而 \(\xi\) 正则。
\end{proposition}
\begin{proof}
域内有 \(|z-a|>R\)，故 \(b>0\)，且它在不含 \(a\) 的该开集调和。趋于 \(\xi\) 时，由 \(|\xi-a|=R\) 得 \(b\to0\)。固定 \(0<\delta<r'<r_0\)，考虑紧集
\[
 F=\overline D\cap\{\,z\mid\delta\le|z-\xi|\le r'\,\}.
\]
若 \(F\) 非空，唯一接触条件使 \(F\) 与 \(\overline{B(a,R)}\) 不交。连续函数 \(|z-a|-R\) 因而在 \(F\) 上有严格正的最小值，所以 \(\log(|z-a|/R)\) 在相应域内环带上有统一正下界。由障碍判据得结论。
\end{proof}
对于 \(C^2\) 图像边界，可把接触点移到原点、切线转为水平线，写 \(D\) 局部在 \(y>g(x)\)，其中 \(g(0)=g'(0)=0\)、\(|g(x)|\le Cx^2\)。取下方圆心 \((0,-a)\) 和足够小的半径 \(a>0\)，使整个圆盘留在图像邻域且 \(1/(2a)>C\)。该圆的上半边满足
\[
 -a+\sqrt{a^2-x^2}\le-\frac{x^2}{2a}<-Cx^2\le g(x)
 \quad(0<|x|<a).
\]
圆盘位于域外，并只在原点接触。这说明光滑边界为何常能使用外球障碍；只有 Lipschitz 正则性时，则直接使用外锥更方便。

\ProseParagraphBreak

圆环的两条边界圆均正则。\subsecref{b3:subsec:160204}中的显式解表明，内圈取一、外圈取零时
\[
 u_r(z)=\frac{\log(R/|z|)}{\log(R/r)}
 \quad(r<|z|<R).
\]
固定非零内点后，\(r\downarrow0\) 时该值按 \(1/\log(R/r)\) 的量级趋于零。若把内圈缩成孤立点并在那里指定值一，\hyperref[b3:exmp:16-punctured-disk]{上一节的穿孔圆盘算例}已经用圆环比较证明 Perron 包络恒零。该点不正则，因而也不存在本节定义的障碍。

\ProseParagraphBreak

外锥和孤点的差别还可以用能量度量。下面只建立所需的相对容量及两个直接计算，为理解更一般的边界判据提供参照。

\Needspace{6\baselineskip}
\begin{definition}\label{b3:def:16-relative-capacity}
设 \(U\subseteq\mathbb C\) 为有界开集，\(K\Subset U\) 为紧集。定义 \(K\) 相对于 \(U\) 的\term[xiangdui-rongliang]{相对容量}[relative capacity]为
\[
 \operatorname{cap}(K,U)\coloneqq
 \inf\left\{\,\int_U|\nabla v|^2\,\mathrm{d}A
 \ \middle|\ v\in C_c^\infty(U;\mathbb R),\quad
 v\ge1\text{ 于 }K\text{ 的某邻域}\,\right\}.
\]
\end{definition}
\symidx{\(\operatorname{cap}(K,U)\)：紧集 \(K\) 相对于开集 \(U\) 的容量}

紧集增大时，试函数限制增加，容量不减；外域增大时，可用试函数增加，容量不增。二维能量还具有尺度不变性：换元 \(z=\xi+s\zeta\) 使梯度的平方乘 \(s^{-2}\)，面积元乘 \(s^2\)，两者抵消。这解释了为什么平面上的比较必须保留内外尺度之比。

\Needspace{6\baselineskip}
\begin{proposition}[圆盘的相对容量]\label{b3:prop:16-annulus-capacity}
若 \(0<r<R\)，则
\[
 \operatorname{cap}\bigl(\overline{B(0,r)},B(0,R)\bigr)
 =\frac{2\pi}{\log(R/r)},
 \quad
 \operatorname{cap}\bigl(\{0\},B(0,R)\bigr)=0.
\]
\end{proposition}
\begin{proof}
对任意合法试函数 \(v\)，有 \(v(re^{\mathrm{i}\theta})\ge1\)，而靠近 \(R\) 的圆周值为零。沿每条半径积分，再用带权 Cauchy--Schwarz 不等式，得到
\[
 \begin{aligned}
 1&\le\left|\int_r^R\partial_s v(se^{\mathrm{i}\theta})\,\mathrm{d}s\right|,\\
 1&\le\left(\int_r^R|\partial_s v(se^{\mathrm{i}\theta})|^2s\,\mathrm{d}s\right)
           \log(R/r).
 \end{aligned}
\]
对 \(\theta\) 积分，并用 \(|\nabla v|^2\ge|\partial_s v|^2\)，得到能量下界 \(2\pi/\log(R/r)\)。

为达到这个下界的极限，先取 \(r<r'<R'<R\)，令
\[
 w(z)=
 \begin{cases}
 1,&|z|\le r',\\
 \dfrac{\log(R'/|z|)}{\log(R'/r')},&r'<|z|<R',\\
 0,&|z|\ge R'.
 \end{cases}
\]
这个函数连续、分段光滑；沿两条连接圆周积分分部时函数值项相消，所以其弱梯度就是各片的通常梯度。直接计算得
\[
 \int_{\mathbb C}|\nabla w|^2\,\mathrm{d}A
 =\frac{2\pi}{\log^2(R'/r')}\int_{r'}^{R'}\frac{\mathrm{d}s}{s}
 =\frac{2\pi}{\log(R'/r')}.
\]
取非负光滑单位核 \(\eta_\varepsilon\)，支集半径为 \(\varepsilon\)，其中
\(\varepsilon<\tfrac12\min(r'-r,R-R')\)。磨光 \(w_\varepsilon=w*\eta_\varepsilon\) 在 \(|z|<r'-\varepsilon\) 等于一，支集包含于 \(\overline{B(0,R'+\varepsilon)}\Subset B(0,R)\)，所以属于容量的试函数类。又由弱梯度卷积公式和 Jensen 不等式，
\[
 |\nabla w_\varepsilon|^2
 \le |\nabla w|^2*\eta_\varepsilon,
 \quad
 \int|\nabla w_\varepsilon|^2\,\mathrm{d}A
 \le\int|\nabla w|^2\,\mathrm{d}A.
\]
令 \(r'\downarrow r\)、\(R'\uparrow R\)，便得相反的容量上界。

单点的情形也有明确试函数。取 \(0<\delta<R/8\)，在上述构造中令 \(r'=2\delta\)、\(R'=R/2\)，再用支集半径小于 \(\delta/2\) 的非负单位核磨光。所得函数在 \(|z|<\delta\) 等于一，支集包含于 \(B(0,3R/4)\)，能量不超过
\[
 \frac{2\pi}{\log(R/(4\delta))}\longrightarrow0
 \quad(\delta\downarrow0).
\]
容量非负，故单点的相对容量为零。
\end{proof}
圆环解 \(u_r\) 正是上述截断对数在过渡环带内的调和部分；其能量的分母与小孔影响的分母相同。容量为零表示可以用能量任意小的光滑函数覆盖这个点，并不表示它与域的欧氏距离较大。

\ProseParagraphBreak

外锥则在每个足够小的尺度留下统一的容量下界。设补集在 \(\xi\) 附近包含一个固定角宽的锥。由缩放，可取只依赖角宽的 \(0<\kappa<1/4\)，使每个充分小的 \(s>0\) 都有一个闭圆盘 \(\overline{B(a_s,\kappa s)}\) 位于该锥内，且 \(|a_s-\xi|=s/2\)。于是该圆盘包含于
\(K_s=(\mathbb C\setminus D)\cap\overline{B(\xi,s)}\)，而 \(B(\xi,2s)\subset B(a_s,3s)\)。由容量的两种单调性及圆盘公式，
\[
 \operatorname{cap}\bigl(K_s,B(\xi,2s)\bigr)
 \ge\operatorname{cap}\bigl(\overline{B(a_s,\kappa s)},B(a_s,3s)\bigr)
 =\frac{2\pi}{\log(3/\kappa)}>0.
\]
下界与 \(s\) 无关。外锥在不断缩小的尺度上保留了能量影响，孤点则没有这种性质。

\ProseParagraphBreak

\term[wiener-panju]{Wiener 判据}[Wiener criterion]把一般边界点的正则性与补集在逐层尺度上的容量条件联系起来。圆盘和外锥提供了这种尺度比较的可算模型。推广到任意闭补集时，需要用容量构造势、累积不同尺度的影响，并建立相应条件的必要性。Ransford 的《Potential Theory in the Complex Plane》\cite{Ransford1995Potential}可作为继续学习的资料。

% !TeX root = ../../Book3.tex
\section{调和测度}
\label{b3:sec:1604}

正则域上的解算子保持线性与正性。测度表示把整个解算子压缩成一族边界概率测度，共形映射再将这些测度从一个区域搬到另一个区域。

\subsection{解的表示}
\label{b3:subsec:160401}

设 \(D\) 是有界正则域。固定 \(z\in D\)，映射 \(\varphi\mapsto H_\varphi(z)\) 为 \(C(\partial D)\) 上的正线性泛函：两解相加仍给相加边值，唯一性使解算子线性；非负边值由最大值原理给非负解。常数一的解为一。紧空间上的正泛函表示遂给唯一概率测度 \(\omega_D^z\)，满足
\[
H_\varphi(z)=\int_{\partial D}\varphi\,\mathrm{d}\omega_D^z.
\]
这里使用第一卷定理11.4.1的 Riesz--Markov 表示。对于不同 \(z\)，这些测度通常互不相同；Harnack 估计使其对非负连续测试函数的积分在内部紧集上相互控制。

\subsection{共形不变性}
\label{b3:subsec:160402}

\begin{proposition}\label{b3:prop:16-harmonic-measure-conformal}
若 \(F:D\to D'\) 共形双射，且延拓为闭包间同胚，两域均有界正则，则
\(F_\#\omega_D^z=\omega_{D'}^{F(z)}\)。
\end{proposition}
\begin{proof}
对 \(\varphi\in C(\partial D')\)，\(H_\varphi\circ F\) 调和，因为 \(\Delta(u\circ F)=|F'|^2\cdot(\Delta u)\circ F\)。它连续取边值 \(\varphi\circ F\)，故由唯一性等于 \(H_{\varphi\circ F}\)。两边在 \(z\) 的值给两测度对全部连续测试函数积分相等，表示的唯一性得结论。
\end{proof}
若映射只在内部存在，边界像可能需要素端而不是通常拓扑边界。此时不能无说明地写相同的推前公式。

\subsection{概率解释的预告}
\label{b3:subsec:160403}

概率论中，调和测度可解释为平面 Brownian 运动首次退出区域时的位置分布。这个解释需要先构造过程，再证明退出时间、停时和边界到达的性质；它不是概率测度表示定理本身的内容。

圆盘中心的退出分布由旋转对称性应当均匀，与 Poisson 核的中心值一致。非中心处则偏向附近边界弧。本章使用的是已构造的解与测度，不以这一直观概率解释代替 Perron 证明。


\ProseParagraphBreak

解表示蕴含测度之间的定量关系。取 \(z,w\) 属于同一内部紧集，用有限 Harnack 链得到 \(C^{-1}H_\varphi(w)\le H_\varphi(z)\le C H_\varphi(w)\) 对所有非负连续 \(\varphi\) 成立。正测度的序关系因此给
\[
C^{-1}\omega_D^w(E)\le\omega_D^z(E)\le C\omega_D^w(E)
\]
对每个 Borel 边界集合 \(E\) 成立：先由连续函数积分不等式得到测度支配，再用正则性从紧集、开集推广。尤其是否具有零调和测度与内部观察点无关，虽然具体权重随点而变。这个结论只比较同一区域内的点；跨域比较需要域的关系或共形映射。

\ProseParagraphBreak

对边界弧 \(I=[\alpha,\beta]\) 的圆盘权重，中心值为 \((\beta-\alpha)/(2\pi)\)。非中心处可积分核，或用 Möbius 变换 \(T_z(\zeta)=(\zeta-z)/(1-\overline z\zeta)\) 将 \(z\) 搬到原点，得到 \(\omega^z(I)=|T_z(I)|/(2\pi)\)，这里弧长在像的单位圆上测量。原来的近弧会被展开，远弧会被压缩；这把 Poisson 核分母 \(|e^{\mathrm{i}t}-z|^2\) 中的距离效应转成圆周变换的导数。必须取相应连续像弧，不能只减两个主辐角而忽略跨越分支切线的情形。

\ProseParagraphBreak

例如上半平面的区间 \([a,b]\) 在 \(x+\mathrm{i}y\) 处的权重是
\[
\omega^{x+\mathrm{i}y}([a,b])=\frac1\pi\left(\arctan\frac{b-x}{y}-\arctan\frac{a-x}{y}\right).
\]
取 \(y\downarrow0\)，当 \(x\in(a,b)\) 时极限一，区间外极限零，端点径向极限二分之一。这个算式既检验归一化，也显示端点与弧内部的边界行为不同。若同时改变 \(x\) 和 \(y\) 以斜向接近端点，比例 \((a-x)/y\) 可趋于不同实数，因而极限也相应变化。

\ProseParagraphBreak

调和测度一般不等于普通边界长度；即使圆盘内也只有中心处是均匀的。对有尖角的域，共形变换在边界的伸缩会加强某些位置的权重，正则域这一性质本身没有承诺与弧长的密度存在。下一章光滑边界的 Green 函数法向导数将给具体密度公式，其额外边界光滑性正是微分边界核所需的条件。

% 本章小结（不计入编号）
\begin{chaptersummary}
Poisson 核以正性和质量集中同时解决圆盘存在性与边界极限。Perron 上包络通过内部调和提升成为调和函数，边界障碍再把它识别为经典解。正则域上的唯一性使解算子保持线性，正性使其具有概率测度表示。孤立缺口与外锥分别展示边界几何造成的失败和成功。
\end{chaptersummary}


\BookExercises

% 本题组自拟；各题新增训练动作见题意，未复制第三方原题。
\begin{exercise}\label{b3:ex:16-draft-01}
求单位圆周边值 \(1+2\cos t-3\sin(2t)\) 的延拓及中心值。
\end{exercise}
\begin{solution}
延拓为 \(1+2r\cos\theta-3r^2\sin(2\theta)\)，中心值一。各项为全纯多项式实部或虚部，边界值直接核对，唯一性保证结果。
\end{solution}

\begin{exercise}\label{b3:ex:16-draft-02}
若连续边值 \(\varphi\ge0\) 在一段弧上严格为正，证明其圆盘延拓在每个内部点严格为正。
\end{exercise}
\begin{solution}
连续性给一段正长度小弧及常数 \(c>0\)，使边值至少为 \(c\)。Poisson 核在该弧严格正，积分的这部分严格正；其余部分非负。
\end{solution}

\begin{exercise}\label{b3:ex:16-draft-03}
求单位圆盘中心处半圆弧的调和测度，并解释为何对一般内部点不必等于二分之一。
\end{exercise}
\begin{solution}
中心处为归一化弧长，所以等于二分之一。趋近半圆弧内部的边界点时，相应 Poisson 积分趋一，趋近另一半圆弧内部时趋零，故该权重不是常数。
\end{solution}

\begin{exercise}\label{b3:ex:16-draft-04}
证明有界调和函数在圆盘的孤立缺口可去。
\end{exercise}
\begin{solution}
在半径 \(R\) 圆周取同边值 Poisson 解 \(h\)，差 \(w\) 有界且外圆为零。在 \(\varepsilon<|z|<R\) 用调和函数 \(C\log(R/|z|)/\log(R/\varepsilon)\) 上下控制 \(w\)。固定 \(z\) 后令 \(\varepsilon\downarrow0\)，得 \(w=0\)，故由 \(h\) 延拓。
\end{solution}

\begin{exercise}\label{b3:ex:16-draft-05}
构造上半平面上两种零边值调和解，并写出保证唯一性的一种附加条件。
\end{exercise}
\begin{solution}
零函数与 \(y\) 都连续取零实轴边值。若要求解有界，则差可奇反射到整平面，反射后的有界调和函数由 Harnack 推出的 Liouville 原理为常数，奇性迫使零。局部反射可由 Poisson 唯一性核验。
\end{solution}

\begin{exercise}\label{b3:ex:16-draft-06}
设有界正则域上 \(\varphi\le\psi\)。证明两个解满足相同的序关系，并估计它们的最大范数差。
\end{exercise}
\begin{solution}
差的边值非负，最大值原理得 \(H_\varphi\le H_\psi\)。把边值差夹在 \(\pm\|\varphi-\psi\|_\infty\) 之间并比较常数解，得内部最大范数差不超过边值差的最大范数。
\end{solution}

\begin{exercise}\label{b3:ex:16-draft-07}
在角宽 \(3\pi/2\) 的扇形顶点构造障碍，说明为何 \(r\cos\theta\) 不适用。
\end{exercise}
\begin{solution}
取 \(\alpha=1/2\)，函数 \(r^{1/2}\cos(\theta/2)\) 在 \(|\theta|\le3\pi/4\) 为正，固定半径上有正下界，到顶点为零。\(r\cos\theta\) 在允许角范围一部分为负，不能作为正障碍。
\end{solution}

\begin{exercise}\label{b3:ex:16-draft-08}
证明有界正则域中任意非空相对开边界片 \(I\) 若承载某个非零非负连续边值，则 \(\omega_D^z(I)>0\) 对所有内部点成立。
\end{exercise}
\begin{solution}
取支集在 \(I\) 的非零非负连续函数 \(\varphi\)。其解非负；若某个内部点为零，则强最小值原理迫使恒零，与边界非零矛盾。因此 \(0<H_\varphi(z)\le\|\varphi\|_\infty\omega_D^z(I)\)。
\end{solution}

% !TeX root = ../../Book3.tex
\chapter{Green 函数、Green 域与多连通区域}
\label{b3:ch:17}
\BookChapterNumber{b3:ch:17}
% 本章摘要（不计入编号）
\begin{chapterabstract}
Green 函数把对数点源与零边界约束合在一个正核中。本章由 Perron 解和穷竭构造核，证明最小性、对称性、共形不变性和区域单调性；随后讨论 Green 势、容量与多连通域的周期分解。
\end{chapterabstract}

% 本章目标（不计入编号）
\begin{learninggoals}
能核对圆盘和半平面 Green 核的符号与常数；能利用穷竭和正核上界判断存在性；能区分 Green 域与覆盖意义的双曲型；能用边界调和测度调整有限连通域的共轭周期。
\end{learninggoals}

圆盘中心的函数 \(\log(R/|z|)\) 在边界为零，在原点有一个正无穷点源。把半径增大，内部势就增加；若一直增大到全平面，这个核发散，说明域的大小会影响点源能否有有限的零边界响应。

\ProseParagraphBreak

本章以\chapref{b3:ch:15}的调和比较和\chapref{b3:ch:16}的 Perron 方法为前提。一般势论与容量的进一步阅读为 Ransford 的《Potential Theory in the Complex Plane》\cite{Ransford1995Potential}与 Armitage、Gardiner 的《Classical Potential Theory》\cite{Armitage2001Classical}；仅核书目和主题的部分不承担这里的证明来源。

% !TeX root = ../../Book3.tex
\section{Green 函数}
\label{b3:sec:1701}

在区域内放置一个点源，并让边界保持零势，需要同时安排局部奇性和全边界条件。对数基本解给出点源，Dirichlet 问题则补上一个调和修正。本章采用正 Green 核，奇性为 \(-\log|z-w|\)，因而其 Laplace 算子带负号。先在有界正则域完成这项构造，再说明光滑边界如何用于核的对称性。

\subsection{定义}
\label{b3:subsec:170101}

设 \(D\subseteq\mathbb C\) 为有界区域，\(w\in D\)。若实函数 \(g_D(\cdot,w)\) 在 \(D\setminus\{w\}\) 严格为正且调和，在每个 \(\xi\in\partial D\) 满足
\[
 \lim_{D\ni z\to\xi}g_D(z,w)=0,
\]
并且 \(g_D(z,w)+\log|z-w|\) 可在 \(w\) 附近延拓为调和函数，则称它为以 \(w\) 为极点的经典零边值\term[green-hanshu]{Green 函数}[Green function]。局部奇性中的系数固定为一；只规定趋于正无穷，不能确定点源的大小。

\ProseParagraphBreak
记调和余项为 \(h_w\)，则在整个 \(D\) 内有
\[
 g_D(z,w)=-\log|z-w|+h_w(z),\quad z\ne w.
\]
这里 \(h_w\) 在极点外由等式定义，在极点处由假设延拓；因此它在全域调和。由\cref{b3:prop:15-log-fundamental}，对数函数局部可积且
\(\Delta\log|z-w|=2\pi\delta_w\)，所以
\[
 -\Delta_zg_D(\cdot,w)=2\pi\delta_w
\]
按分布成立。这个等式记录点源，零边值则记录区域的约束，两项都属于核的定义。

\ProseParagraphBreak
为了在边界使两项相消，所求修正必须满足
\(h_w(\xi)=\log|\xi-w|\)。因此使用的是正号的对数边值。因为 \(w\) 在内部，它到边界有正距离，这份边值在紧集 \(\partial D\) 上连续且有界；上一章的连续 Dirichlet 理论正好可以处理它。

\Needspace{6\baselineskip}
\begin{theorem}[有界正则域上的 Green 核]\label{b3:thm:17-green-regular}
设 \(D\subseteq\mathbb C\) 为有界正则域，即每个连续实边值的 Perron 解都连续达到全部边界值。对 \(w\in D\)，令
\[
 \varphi_w(\xi)=\log|\xi-w|,\quad \xi\in\partial D,
 \quad h_w=H_{\varphi_w}.
\]
则
\[
 g_D(z,w)=-\log|z-w|+h_w(z),\quad z\in D\setminus\{w\},
\]
是以 \(w\) 为极点的经典零边值 Green 函数。
\end{theorem}
\begin{proof}
极点到边界的正距离保证 \(\varphi_w\in C(\partial D;\mathbb R)\)。
\cref{b3:thm:16-perron}给出内部调和的 Perron 解，而 \(D\) 的正则性保证
\(h_w\in C(\overline D)\) 且 \(h_w|_{\partial D}=\varphi_w\)。因此所定义的
\(g_D\) 在极点外调和，连续取零边值，其对数奇性也符合定义。

还须证明严格正性。先取 \(\rho>0\) 使
\(\overline{B(w,\rho)}\subset D\)。调和函数 \(h_w\) 在这个闭圆盘上有界，故
\[
 \min_{|z-w|=\varepsilon}g_D(z,w)
 \ge-\log\varepsilon-\sup_{\overline{B(w,\rho)}}|h_w|
 \longrightarrow+\infty
 \quad(\varepsilon\downarrow0).
\]
取足够小的 \(\varepsilon\)，在
\(D\setminus\overline{B(w,\varepsilon)}\) 上对 \(-g_D\) 应用边界最大值原理：原边界上它为零，小圆上它为负，所以 \(g_D\ge0\) 于余域；若余域有多个连通分支，就在每个分支分别使用同一原理。任意固定 \(z\ne w\) 都可放进这样的余域，故全域非负。

区域 \(D\setminus\{w\}\) 仍连通：连接两点的内部折线若经过 \(w\)，可在一个内部小圆盘中以圆弧绕开该点。若 \(g_D\) 在某个非极点位置为零，强最小值原理便使它在整个 \(D\setminus\{w\}\) 恒零，与极点附近趋于正无穷矛盾。因此 \(g_D>0\)。
\end{proof}

构造只用正则性，没有用边界导数。若记
\[
 d_w=\min_{\xi\in\partial D}|\xi-w|,\quad
 L_w=\max_{\xi\in\partial D}|\xi-w|,
\]
则 \(0<d_w\le L_w<\infty\)，最大值原理给出
\[
 \log d_w\le h_w(z)\le\log L_w,\quad z\in D.
\]
因此极点附近有明确的比较
\[
 \log\frac{d_w}{|z-w|}
 \le g_D(z,w)\le\log\frac{L_w}{|z-w|}.
\]
左端离极点较远时可以为负，严格正性仍由上面的域内比较证明。在圆盘
\(B(w,R)\) 中，边值 \(\varphi_w\) 恒等于 \(\log R\)，所以
\(h_w=\log R\)，得到 \(g_{B(w,R)}(z,w)=\log(R/|z-w|)\)。

\subsection{唯一性}
\label{b3:subsec:170102}

同一对数奇性的差不再含点源，因此经典核的唯一性仍归结为调和函数的零边值唯一性。这一步只要求区域有界；正则性是保证存在性的条件。

\Needspace{6\baselineskip}
\begin{proposition}[经典零边值核的唯一性]\label{b3:prop:17-green-uniqueness}
设 \(D\subseteq\mathbb C\) 为有界区域，\(w\in D\)。在
\(D\setminus\{w\}\) 调和、在每个边界点连续取零值，且
\(g(z)+\log|z-w|\) 在 \(w\) 附近调和可去的实函数 \(g\) 至多一个。
这里不必预先假定候选函数为正。
\end{proposition}
\begin{proof}
若 \(g_1,g_2\) 是两个候选，则在极点附近写成
\(g_j=-\log|z-w|+h_j\)。它们的差等于 \(h_1-h_2\)，故可在 \(w\) 处延拓为调和函数。延拓后的差在全域调和，并在每个边界点趋于零。由有界域的边界上极限原理\cref{b3:cor:15-boundary-limsup}，对差与负差分别比较，得到 \(g_1-g_2=0\)。
\end{proof}

经典核还有一个不要求竞争者满足零边值的性质。它将用于把定义扩展到边界粗糙或无界的区域。

\Needspace{6\baselineskip}
\begin{proposition}[经典 Green 核的最小性]\label{b3:prop:17-classical-green-minimal}
设 \(D\subseteq\mathbb C\) 为有界正则域，\(w\in D\)，\(g_D(\cdot,w)\) 为其经典 Green 核。若实函数 \(G\) 在 \(D\setminus\{w\}\) 为正且调和，并且 \(G(z)+\log|z-w|\) 在 \(w\) 附近调和可去，则
\[
 g_D(z,w)\le G(z),\quad z\in D\setminus\{w\}.
\]
对 \(G\) 不要求边界极限存在，也不要求它在靠近边界时有界。
\end{proposition}
\begin{proof}
函数 \(q=g_D-G\) 的奇性相消，故延拓为 \(D\) 上的调和函数。因为
\(G>0\)，对每个 \(\xi\in\partial D\) 都有
\[
 \limsup_{D\ni z\to\xi}q(z)
 \le\lim_{D\ni z\to\xi}g_D(z,w)=0.
\]
区域有界，\cref{b3:cor:15-boundary-limsup}直接给 \(q\le0\)，即得结论。
\end{proof}

最小性比较的是同一对数奇性的全部正竞争函数。若向经典核加一个正的调和函数，奇性与正性仍然保留，但一般不再取零边值，也不再是最小者。对无界域，甚至通常边界上的零值也可能保留：上半平面的函数
\(y\) 为正、调和且在实轴为零，可以加到任何已有的正零边值候选上。因此后面的穷竭构造将以最小性识别一般 Green 核，而不要求它在每个不正则边界点取零。

\ProseParagraphBreak
以下边界积分限定在有限条实解析 Jordan 曲线围成的有界域。这里曲线的实解析参数要求导数处处非零，因而排除了尖点。这类边界局部是 \(C^1\) 图像，域的补集包含一个外锥，由\cref{b3:prop:16-exterior-cone-regular}可知域正则。零边值的反射还会给出法向导数，不能在证明反射之前先假定这些导数存在。

\Needspace{6\baselineskip}
\begin{lemma}[常边值调和函数的局部反射]\label{b3:lem:17-harmonic-reflection}
设 \(D\subseteq\mathbb C\) 为有界区域，其边界由有限条互不相交的正则实解析 Jordan 曲线组成。设 \(I\) 是一段开边界弧，\(V\subseteq\mathbb C\) 是包含 \(I\) 的开集，实函数 \(u\) 在 \(D\cap V\) 调和，并在 \(I\) 上连续取同一常值 \(c\)。则对每个 \(\xi\in I\)，\(u\) 都可调和延拓到包含 \(\xi\) 的一个开邻域。因此 \(u\) 在该弧局部光滑直到边界。
\end{lemma}
\begin{proof}
在 \(\xi\) 附近取实解析参数 \(\gamma(t)\)，使 \(\gamma(0)=\xi\)、
\(\gamma'(0)\ne0\)。其幂级数在一个复圆盘中给出全纯延拓
\(\Gamma(\zeta)\)。由复逆函数定理，缩小圆盘后 \(\Gamma\) 双全纯，且将一段实轴映为边界弧。正则 Jordan 边界在局部把邻域分成两侧；调整参数方向后，可使上半圆盘映到 \(D\) 侧。选 \(a>0\) 足够小，使闭上半圆盘
\(\{|\zeta|\le a,\operatorname{Im}\zeta\ge0\}\) 都落在这张坐标的范围内，且半圆弧除两端外映入 \(D\)。

令 \(v(\zeta)=u(\Gamma(\zeta))-c\)。链式法则给
\[
 \Delta v=|\Gamma'|^2\cdot(\Delta u)\circ\Gamma=0
\]
于上半圆盘。它连续到闭上半圆盘，并在直径上为零。在完整圆周
\(|\zeta|=a\) 上定义数据
\[
 \psi(\zeta)=
 \begin{cases}
 v(\zeta),&\operatorname{Im}\zeta\ge0,\\
 -v(\overline\zeta),&\operatorname{Im}\zeta<0.
 \end{cases}
\]
两端 \(\zeta=\pm a\) 的值为零，因此 \(\psi\) 是连续圆周数据。由
\cref{b3:thm:16-poisson}，其 Poisson 解 \(V\) 在整圆盘调和并连续到圆周。

函数 \(-V(\overline\zeta)\) 具有相同边值 \(\psi\)，圆盘解的唯一性给
\(V(\zeta)=-V(\overline\zeta)\)，所以 \(V\) 在直径上为零。在上半圆盘，
\(V\) 与 \(v\) 都调和，且在半圆弧与直径上取相同值；有界域最大值原理给
\(V=v\)。于是
\[
 c+V\circ\Gamma^{-1}
\]
在 \(\xi\) 的完整邻域中调和，并在域内与原来的 \(u\) 相同。这就是所需延拓。
\end{proof}

对固定极点 \(w\)，边界小邻域可选得不含 \(w\)，故该引理适用于
\(g_D(\cdot,w)\)。边界紧且分量有限，取有限个这样的邻域就得到整个边界上的光滑性及法向导数的连续性。它们足以在删去极点小圆后使用 Green 公式；一般粗糙域的对称性将由穷竭传递，不在粗糙边界上写通常法向积分。

\subsection{对称性}
\label{b3:subsec:170103}

Green 核描述源点与观察点之间的作用。交换二者以后核值不变，需要把两份核放进同一个积分恒等式；两处奇性分别贡献一个函数值。

\Needspace{6\baselineskip}
\begin{proposition}[Green 核的对称性]\label{b3:prop:17-green-symmetry}
设 \(D\subseteq\mathbb C\) 为有界区域，其边界由有限条互不相交的正则实解析 Jordan 曲线组成，\(g_D\) 为对数奇性取 \(-\log|z-w|\) 的经典零边值 Green 核。则对不同的 \(z,w\in D\)，有
\[
 g_D(z,w)=g_D(w,z).
\]
\end{proposition}
\begin{proof}
固定 \(z\ne w\)，在变量 \(\zeta\) 中记
\(u(\zeta)=g_D(\zeta,z)\)、\(v(\zeta)=g_D(\zeta,w)\)。取足够小的
\(\varepsilon>0\)，使以 \(z,w\) 为中心的两个闭圆盘互不相交且均包含于
\(D\)。令
\[
 \Omega_\varepsilon
 =D\setminus\bigl(\overline{B(z,\varepsilon)}
                  \cup\overline{B(w,\varepsilon)}\bigr).
\]
两函数在这个域中调和，且由边界反射光滑到原边界。这里使用第二卷推论1.2.3的 Green 公式：区域有界、边界为 \(C^1\)，函数在闭域的邻域中为 \(C^2\)。原边界实解析，新增小圆也光滑；反射与极点外的光滑性给出所需邻域延拓，故这些条件均满足。对恒等式
\[
 \operatorname{div}(u\nabla v-v\nabla u)=u\Delta v-v\Delta u
\]
使用平面散度公式，得到 Green 第二公式
\[
 0=\int_{\partial\Omega_\varepsilon}
       (u\,\partial_nv-v\,\partial_nu)\,\mathrm ds,
\]
其中 \(n\) 一律指 \(\Omega_\varepsilon\) 的外单位法向。原边界的每个分量上
\(u=v=0\)，故那里没有贡献。

在 \(z\) 附近写 \(u(\zeta)=-\log|\zeta-z|+a(\zeta)\)，其中 \(a\) 光滑调和；函数 \(v\) 在这里也光滑。小圆上的外法向朝向被删圆盘的中心，故
\[
 \begin{aligned}
 u&=-\log\varepsilon+O(1),&
 \partial_nu&=\frac1\varepsilon+O(1),\\
 v&=v(z)+O(\varepsilon),&
 \partial_nv&=O(1).
 \end{aligned}
\]
所有误差均对圆周方向一致。因此
\[
 \int_{\partial B(z,\varepsilon)}
       (u\,\partial_nv-v\,\partial_nu)\,\mathrm ds
 =-2\pi v(z)+O(\varepsilon|\log\varepsilon|)+O(\varepsilon).
\]
第一项 \(u\,\partial_nv\) 乘上圆周长后趋于零；保留下来的是带负号的
\(-v\,\partial_nu\)。

在 \(w\) 的小圆上两函数的角色交换：
\(v=-\log\varepsilon+O(1)\)、\(\partial_nv=1/\varepsilon+O(1)\)，而
\(u=u(w)+O(\varepsilon)\)、\(\partial_nu=O(1)\)。这里得到
\[
 \int_{\partial B(w,\varepsilon)}
       (u\,\partial_nv-v\,\partial_nu)\,\mathrm ds
 =2\pi u(w)+O(\varepsilon|\log\varepsilon|)+O(\varepsilon).
\]
两小圆积分之和为零，令 \(\varepsilon\downarrow0\)，得
\(u(w)=v(z)\)，即所述对称性。
\end{proof}

小圆上的法向不是从其中心向外，而是余域的向外方向，即朝向被删掉的圆盘。这个方向使对数奇项的法向导数为 \(+1/\varepsilon\)；两处主项的相反符号则来自
\(u\,\partial_nv-v\,\partial_nu\) 中两函数的位置。对称性还表明，对固定
\(z\)，函数 \(w\mapsto g_D(z,w)\) 在 \(D\setminus\{z\}\) 调和，因为它等于已按第一个变量构造的 \(g_D(w,z)\)。

\subsection{共形不变性}
\label{b3:subsec:170104}

共形变换保持二维调和方程，但对数奇性会多出一个调和项。经典核的边界条件还需要映射在闭包上的对应；先保留这项条件，可以直接使用唯一性。

\Needspace{6\baselineskip}
\begin{proposition}[经典 Green 核的共形不变性]\label{b3:prop:17-classical-green-conformal}
设 \(D,D'\subseteq\mathbb C\) 为有界正则域，\(F:D\to D'\) 双全纯，且
\(F\) 延拓为 \(\overline D\) 到 \(\overline{D'}\) 的同胚。则对不同的
\(z,w\in D\)，两域的经典 Green 核满足
\[
 g_D(z,w)=g_{D'}(F(z),F(w)).
\]
\end{proposition}
\begin{proof}
固定 \(w\)，令 \(G(z)=g_{D'}(F(z),F(w))\)。它为正，并由调和函数复合公式在
\(D\setminus\{w\}\) 调和。闭包同胚将 \(\partial D\) 对应到
\(\partial D'\)，故 \(G\) 在每个边界点趋于零。

在 \(w\) 的小圆盘上写
\[
 F(z)-F(w)=(z-w)q(z),\quad q(w)=F'(w)\ne0.
\]
缩小圆盘使全纯函数 \(q\) 处处非零，于是 \(\log|q|\) 调和。若
\(g_{D'}(\eta,F(w))=-\log|\eta-F(w)|+h(\eta)\)，则
\[
 G(z)+\log|z-w|=h(F(z))-\log|q(z)|
\]
在 \(w\) 附近调和。因而 \(G\) 具有正确的奇性和零边值，由
\cref{b3:prop:17-green-uniqueness}得到 \(G=g_D(\cdot,w)\)。
\end{proof}

映射在边界上不必有非零导数；本命题只在那里使用同胚和连续性，导数非零只用在内部极点。相似变换也说明为何核不能只看奇性中的距离：若
\(F(z)=az+b\)、\(a\ne0\)，则
\(-\log|F(z)-F(w)|=-\log|z-w|-\log|a|\)，调和修正必须补回这项常数，才能保持整个核值不变。

\ProseParagraphBreak
没有通常边界延拓的内部共形双射不满足本命题的边界假设。一般 Green 域仍有同一不变性，但需先建立最小正核，再把正竞争函数双向拉回进行比较；这条不使用边界延拓的证明将在\cref{b3:sec:1703}中给出。

% !TeX root = ../../Book3.tex
\section{Green 域与抛物型域}
\label{b3:sec:1702}

上一节在有界正则域中把点源与连续零边值配在一起。若区域无界，或者边界含有不正则点，就不能继续把通常的边界值当作唯一的约束。本节先问一个更基本的问题：在整个区域内，能否有一个始终为正、只有指定对数奇性的函数？下一节再从所有这样的函数中构造最小者。

\subsection{Green 域}
\label{b3:subsec:170201}

为了在尚未讨论无穷远行为时准确表述存在性，先把竞争函数类写出来。

\Needspace{6\baselineskip}
\begin{definition}[正基本核与 Green 域]\label{b3:def:17-green-domain}
设 \(D\subseteq\mathbb C\) 为区域，\(w\in D\)。称实函数 \(q:D\setminus\{w\}\to(0,\infty)\) 为极点 \(w\) 的正基本核，如果它在 \(D\setminus\{w\}\) 调和，且 \(q(z)+\log|z-w|\) 在 \(w\) 的某个圆盘邻域中延拓为调和函数。记所有正基本核的集合为 \(\mathcal Q_D(w)\)。若对每个 \(w\in D\) 都有 \(\mathcal Q_D(w)\ne\varnothing\)，称 \(D\) 为\term[green-yu]{Green 域}[Greenian domain]。当 \(\mathcal Q_D(w)\) 中存在逐点不超过其余各项的成员时，称这个最小成员为一般域上的 Green 核，记为 \(g_D(\cdot,w)\)。
\end{definition}

“最小”指与全体竞争函数逐点比较，不是只在某个内点使数值最小。一个类如果有最小成员，该成员必唯一。上一节已经证明，有界正则域的经典 Green 核就是其正基本核类的最小成员。下一节将证明，只要该类非空，最小成员就存在；还将证明一处极点的存在性足以推出全部极点的存在性。因此这里的定义与用有限穷竭极限定义 Green 域等价，但现在不把尚未构造的极限作为论证前提。

\ProseParagraphBreak
有界域不必正则，也一定有正基本核。固定 \(w\in D\)，选择 \(R>0\) 使 \(D\subset B(w,R)\)，则
\[
 q(z)=\log\frac R{|z-w|}
\]
在 \(D\setminus\{w\}\) 严格为正、调和，且 \(q+\log|z-w|=\log R\)。所以每个有界区域都是 Green 域。这个竞争核一般不在 \(\partial D\) 上取零；其用途是给随后构造的最小核一个有限上界。边界是否正则影响连续零边值的达到，不决定有界域是否有 Green 核。

\ProseParagraphBreak
正基本核也可以限制到子域。若 \(D\subset E\)，\(w\in D\)，则 \(\mathcal Q_E(w)\) 的每一项在 \(D\) 上仍有正确奇性和正性。因此 Green 域的每个子域也是 Green 域。这一判断只用竞争函数的存在，尚不需要最小核的区域单调性。

\subsection{Green 函数不存在的情形}
\label{b3:subsec:170202}

全平面的困难出现在无穷远。假设极点 \(w\) 有一个正基本核 \(q\)，则
\[
 h(z)=q(z)+\log|z-w|
\]
延拓为整个平面的调和函数。在以 \(w\) 为中心的圆周上取平均，得到
\[
 \frac1{2\pi}\int_0^{2\pi}q(w+re^{\mathrm it})\,\mathrm dt
 =h(w)-\log r.
\]
当 \(r\) 足够大时右端为负，与 \(q>0\) 矛盾。因此整个复平面没有正基本核，更没有 Green 核。这里不是因为没有可写的基本解：\(-\log|z-w|\) 正是所需基本解，只是任何整平面调和修正都不能使它在全部半径上保持正平均。

\ProseParagraphBreak
删除有限个点是否可以补救这一困难，需要先控制那些点附近可能出现的奇性。对于任意调和函数，孤立奇性可以比对数强得多；正性却排除了这些振荡主部。

\Needspace{6\baselineskip}
\begin{lemma}[Bôcher 孤立奇性分解]\label{b3:lem:17-positive-singularity}
设 \(R>0\)，\(u\ge0\) 在 \(0<|z|<R\) 调和。则存在唯一的 \(c\ge0\) 及整个圆盘 \(B(0,R)\) 上的实调和函数 \(h\)，使
\[
 u(z)=-c\log|z|+h(z)\quad(0<|z|<R).
\]
\end{lemma}
\begin{proof}
对每个整数 \(n\) 定义圆周系数
\[
 a_n(r)=\frac1{2\pi}\int_0^{2\pi}u(re^{\mathrm it})e^{-\mathrm int}\,\mathrm dt,
 \quad 0<r<R.
\]
任意闭圆环内的 \(u\) 光滑，故可以对 \(r\) 求导；将极坐标 Laplace 方程乘以 \(e^{-\mathrm int}\)，并对角变量积分两次分部，得到
\[
 a_n''(r)+\frac1r a_n'(r)-\frac{n^2}{r^2}a_n(r)=0.
\]
当 \(n=0\) 时，解为 \(a_0(r)=A+B\log r\)。圆周平均非负，令 \(r\downarrow0\) 可知 \(B\le0\)。当 \(n\ne0\) 时，换元 \(x=\log r\) 将方程化为常系数方程，因而
\[
 a_n(r)=\alpha_n r^{|n|}+\beta_n r^{-|n|}.
\]
正性给出关键估计
\[
 |a_n(r)|\le\frac1{2\pi}\int_0^{2\pi}u(re^{\mathrm it})\,\mathrm dt
 =a_0(r)=O(1+|\log r|).
\]
两端乘以 \(r^{|n|}\)，再令 \(r\downarrow0\)，便知 \(\beta_n=0\)。负幂的增长比对数快，因而不能在正性所给的圆周积分界下保留。

固定 \(0<S<R\)，令 \(h_S\) 为圆周数据 \(u(Se^{\mathrm it})-B\log S\) 的 Poisson 延拓。零频率系数为 \(A\)，非零系数为 \(\alpha_n S^{|n|}\)。由\cref{b3:prop:16-poisson-frequencies}，在任意 \(0<r<S\) 上，\(h_S(re^{\mathrm it})\) 的系数恰为 \(A\) 及 \(\alpha_n r^{|n|}\)。这些也正是 \(u(re^{\mathrm it})-B\log r\) 的全部系数。

两个连续圆周函数的所有 Fourier 系数相同，就必相等。为提醒这一唯一性的依据，第16章圆环构造已证明连续周期函数的 Fej\'er 平均一致收敛；如果差函数的全部系数为零，它的每个 Fej\'er 平均都为零，因此差函数为零。于是 \(u-B\log|z|=h_S\) 于 \(0<|z|<S\)。不同 \(S\) 的调和延拓在交叠的去心圆盘一致，连续性使原点值也相同，故拼成整个 \(B(0,R)\) 上的调和函数 \(h\)。取 \(c=-B\ge0\) 得到表示。零频率确定 \(B\)，其余函数由差确定，所以表示唯一。
\end{proof}

这个结论也称正调和函数的\term[bocher-gulixing-fenjie]{Bôcher 孤立奇性分解}[Bôcher decomposition of an isolated singularity]，这里证明的是二维版本。若允许 \(u\) 改变符号，\(\operatorname{Re}(1/z)\) 就有负幂奇性；它的圆周平均为零，绝对值积分却很大，证明中的 \(|a_n|\le a_0\) 不再成立。

\Needspace{6\baselineskip}
\begin{proposition}[有限穿孔平面无正 Green 核]\label{b3:prop:17-finite-punctures-parabolic}
设 \(A\subset\mathbb C\) 为有限集，允许为空。则在 \(D=\mathbb C\setminus A\) 上，对任何 \(w\in D\) 都不存在正基本核。
\end{proposition}
\begin{proof}
空集情形已由圆周平均证明。设 \(A=\{a_1,\ldots,a_m\}\)，假设 \(q\in\mathcal Q_D(w)\)。在每个 \(a_j\) 的一个不含其余缺口及 \(w\) 的小圆盘内，上一引理给
\[
 q(z)=-c_j\log|z-a_j|+h_j(z),\quad c_j\ge0,
\]
其中 \(h_j\) 跨过 \(a_j\) 调和。因此
\[
 H(z)=q(z)+\log|z-w|+\sum_{j=1}^m c_j\log|z-a_j|
\]
延拓为整平面调和函数：在每个被删点，恰由相应局部分解消去奇性；在 \(w\) 则由正基本核的定义消去奇性。

取 \(r\) 大于 \(|w|,|a_1|,\ldots,|a_m|\)。对任意 \(|a|<r\)，有
\[
 \frac1{2\pi}\int_0^{2\pi}\log|re^{\mathrm it}-a|\,\mathrm dt=\log r.
\]
这可由 \(\log|re^{\mathrm it}-a|=\log r+\operatorname{Re}\log(1-ae^{-\mathrm it}/r)\) 的一致收敛幂级数直接积分，所有非零频率平均为零。于是
\[
 \frac1{2\pi}\int_0^{2\pi}q(re^{\mathrm it})\,\mathrm dt
 =H(0)-\left(1+\sum_{j=1}^m c_j\right)\log r.
\]
右端最终为负，左端始终为正，矛盾。
\end{proof}

每个缺口允许新增的系数均为非负，它们只能使大圆平均中的负对数斜率更大。因而不能靠有限个孤立缺口补偿无穷远的正性失败。这也是论证中必须分析缺口奇性，而不能只说“孤立点没有影响”的原因。

\subsection{\texorpdfstring{\(\mathbb C\)}{C} 的抛物性}
\label{b3:subsec:170203}

从势论角度，在任何极点都没有正基本核的区域称为\term[paowuxing-yu]{抛物型域}[parabolic domain]。这里的名称针对 Green 核的不存在性，与时间演化方程的抛物分类无关。全平面及有限穿孔平面都是上述直接证明的例子。下一节会说明，正基本核若在一个极点存在，就在所有极点存在；因此平面区域不会出现“某些极点属于 Green 情形、另一些属于抛物情形”的混合情况。

\ProseParagraphBreak
圆盘扩大的例子已经能看出发散机制。以原点为极点的经典核为 \(g_{B(0,R)}(z,0)=\log(R/|z|)\)，这个函数直接满足上一节的全部定义条件。固定 \(z\ne0\)，随 \(R\to\infty\) 核值趋于正无穷。每个有界圆盘都可容纳有限点源响应，但把零边界不断推向无穷远后，在固定观察点得不到有限极限。

\ProseParagraphBreak
维数在这里有实际作用。在三维，正基本解与 \(|x-w|^{-1}\) 成比例，它在整个空间保持正且在无穷远趋零；二维的基本解则是对数。因而“把边界推到无穷远”是否仍保留有限正核，与基本解的尺度有关，不能把平面的圆周平均论证原样移到高维。

\subsection{覆盖类型与 Green 存在性的区别}
\label{b3:subsec:170204}

Riemann 曲面论通常把万有覆盖为圆盘的曲面称为双曲型；有些势论文献也用“双曲型”表示存在正 Green 函数。两种说法描述不同条件，本章统一用“Green 域”指核的存在性。

\ProseParagraphBreak
去心圆盘是有界域，所以是 Green 域；它的万有覆盖也是圆盘。另一方面，\cref{b3:prop:12-planar-hyperbolic}说明，补集至少含两个点的平面域具有圆盘万有覆盖，故 \(\mathbb C\setminus\{0,1\}\) 属于这种覆盖类型；该结论沿用第12章已标明的一般统一化输入。但\cref{b3:prop:17-finite-punctures-parabolic}表明它没有正 Green 核。这个例子同时具有圆盘覆盖与势论抛物性。已知万有覆盖的共形类型，尚不足以判断点源在原域内是否产生有限正响应。

\ProseParagraphBreak
单连通情形则可以直接转移。若 \(D\subsetneq\mathbb C\) 为单连通域，Riemann 映照定理给出共形双射 \(F:D\to\mathbb D\)。固定 \(w\in D\)，取圆盘的经典 Green 核，并令
\[
 q(z)=g_{\mathbb D}(F(z),F(w)).
\]
它在极点外正且调和。因 \(F'(w)\ne0\)，在 \(w\) 附近有
\[
 F(z)-F(w)=(z-w)A(z),\quad A(w)=F'(w)\ne0,
\]
其中 \(A\) 全纯且局部无零。因此复合核加上 \(\log|z-w|\) 后，只新增调和项 \(-\log|A(z)|\)，仍具有所需奇性。它是 \(\mathcal Q_D(w)\) 的成员，故任意真单连通平面域都是 Green 域。这个存在性判断只用内部共形双射，不要求 \(F\) 连续延拓到通常拓扑边界。

\ProseParagraphBreak
若只有覆盖映射而没有双射，原域的一个极点会有多个原像。把原域的核拉回覆盖空间，便在每一个原像处产生对数奇性，不能把这个拉回误认为覆盖空间上只有一个极点的 Green 核。这正是多连通情形比共形双射多出的困难。

\subsection{与万有覆盖的联系}
\label{b3:subsec:170205}

设 \(\pi:\mathbb D\to D\) 为万有覆盖，固定 \(w\in D\) 及一个原像 \(\widetilde w\)。第11章的覆盖理论说明，纤维由覆盖变换群 \(\Gamma\) 的轨道 \(\{\gamma\widetilde w:\gamma\in\Gamma\}\) 给出；该集合离散且可数。于是一个自然的候选表达是
\[
 \sum_{\gamma\in\Gamma}
 g_{\mathbb D}(\widetilde z,\gamma\widetilde w).
\]
每项都非负；有限部分和只处理有限多个点源，而完整纤维要求控制无穷项。即使每个单项都有限，和仍可能发散。离散性只保证内部紧集包含有限个极点，并不控制趋近单位圆周的无穷多个极点的总贡献。

\ProseParagraphBreak
要用这一表达构造原域核，至少还须证明：部分和在极点轨道外局部一致收敛；在每个原像附近只留下一个系数为一的对数奇性；和对覆盖变换不变，从而能够下降。要进一步识别为最小核，还需最小性比较。这些问题各有作用，不能用“按群求和”一句话取代。本章的一般构造采用域内穷竭，因而不以该级数公式为前提。

\ProseParagraphBreak
从这里也可理解为何前述两种分类会分开。圆盘覆盖给出每个单独源点的正核，却没有自动给出整个轨道的求和控制；势论中的有限性正是额外要求。下一节把这种要求改写成更便于检验的条件：在一个固定非极点位置，穷竭核的数值有界。Harnack 链随后会把这个一点条件传播到所有内部紧集。

% !TeX root = ../../Book3.tex
\section{Green 函数的构造}
\label{b3:sec:1703}

先在规则有界域构造核，再逐步放大区域。单调性与 Harnack 原理控制极限，最小性使构造不依赖所选穷竭。

\subsection{穷竭}
\label{b3:subsec:170301}

任意平面区域可取 \(D_j\) 光滑有界连通子域，\(\overline{D_j}\subset D_{j+1}\)，且并为 \(D\)。例如先取到补集距离至少 \(1/j\)、模不超过 \(j\) 的紧部分，添加连接路径，再在内部用有限圆盘和光滑边界逼近。极点包含在首项之后。

若 \(D_j\subset D_{j+1}\)，核之差在极点可去，且在 \(\partial D_j\) 非负；最大值原理给 \(g_{D_j}\le g_{D_{j+1}}\)。这是真正的单调量。不能直接说边界值逐项收敛，因为边界本身在移动。

\subsection{Perron 方法}
\label{b3:subsec:170302}

对于光滑有界 \(D_j\)，用 Perron 法解边值 \(\log|z-w|\)，得调和函数 \(h_j\)，令 \(g_j=-\log|z-w|+h_j\)。小圆上它趋正无穷、外边界为零，最大值原理给严格正性。

\begin{theorem}[Green 核的穷竭构造]\label{b3:thm:17-green-exhaustion}
若 \(g_j(z_0,w)\) 对某个 \(z_0\ne w\) 有界，则 \(g_j\) 局部一致收敛于有限正 Green 核；该核是所有具有同一对数奇性的正调和核中的最小者，因而与穷竭无关。
\end{theorem}
\begin{proof}
极点外用 Harnack 收敛得到调和极限。在极点周围固定一小圆，边界上的极限有界；\(h_j=g_j+\log|z-w|\) 在圆盘内调和，最大值原理把边界界传到内部。Harnack 或 Poisson 表示使 \(h_j\) 在较小圆盘收敛到调和函数，所以奇性保持。
任取竞争正核 \(G\)，\(G-g_j\) 在极点可去，且在 \(\partial D_j\) 非负，因此 \(g_j\le G\)。取极限得最小性。不同穷竭产生的有限核彼此作为竞争核，故相等。若一个穷竭发散而另一个有限，后者又给前者共同上界，矛盾。
\end{proof}
有限核的对称性由光滑有界核的对称性取极限得到。

\subsection{圆盘与半平面}
\label{b3:subsec:170303}

圆盘与上半平面的显式核为
\[
g_{B(0,R)}(z,w)=\log\left|\frac{R^2-\overline wz}{R(z-w)}\right|,
\quad
g_{\mathfrak H}(z,w)=\log\left|\frac{z-\overline w}{z-w}\right|.
\]
分子零点在区域外，故余项调和。边界上分子和分母的模相等，且区域内两模之比大于一。圆盘由唯一性，半平面由共形变换和最小性确认。

极点在圆盘中心时核简化成 \(\log(R/|z|)\)。上半平面公式是负对数点源与反射源之差；反射点必须为 \(\overline w\)，不能给反射源保留同号。

\subsection{比较原理}
\label{b3:subsec:170304}

\begin{proposition}\label{b3:prop:17-domain-monotonicity}
若 Green 域 \(D\subset E\)，则 \(g_D(z,w)\le g_E(z,w)\)。
\end{proposition}
\begin{proof}
\(E\) 的核限制到 \(D\) 是一个正竞争核，具有正确奇性，最小性立即给不等式。
\end{proof}
将边界推远会减弱零势约束，核因而增大。对有界域，这个原则用包围圆盘给有限上界；它也使某些无界域可由半平面或其他已知 Green 域控制。


\ProseParagraphBreak

穷竭极限还保留极点参数的局部稳定性。在两个不相交的相对紧小圆盘 \(U,V\Subset D\) 上，光滑穷竭核在两个变量各自调和且对称。固定源点先用 Harnack 控制观察点，再利用对称性在源点方向使用同一估计，就可在 \(U\times V\) 上得到共同上界及导数界。极限因此在远离对角线处连续，并在每个变量中调和。仅仅知道逐点单调极限有限，还不能未经此控制就交换源点极限或把核当成连续的双变量函数。

\ProseParagraphBreak

在圆盘里，核的正性可由一个恒等式直接核对：
\[
|R^2-\overline wz|^2-R^2|z-w|^2=(R^2-|z|^2)(R^2-|w|^2)>0.
\]
两点均在圆盘内部时右端严格正，故模之比大于一。上半平面的相应恒等式是 \(|z-\overline w|^2-|z-w|^2=4\operatorname{Im}z\operatorname{Im}w>0\)。这些实代数算式避免凭图形猜测反射点的位置和源的符号；也说明两核在光滑边界分别降到零。

\ProseParagraphBreak

若 \(D\subsetneq E\) 且两域连通，区域单调性一般是严格的。在 \(D\) 内两核差调和、非负；如果在一个内部非极点位置为零，强最小值原理使差在 \(D\) 恒零。对光滑域，如果 \(\partial D\) 有一个点位于 \(E\) 内，\(E\) 的核在那里严格正而 \(D\) 的核取零，便产生矛盾。对于粗糙域，删去一个极性孔可能不改变核，所以不能仅由集合严格包含就断言严格不等式。严格性需要边界差异在势论中具有实际影响。

% !TeX root = ../../Book3.tex
\section{Green 势、能量与容量}
\label{b3:sec:1704}

点源的位置由一个参数表示，分布源则由测度表示。把 Green 核按源测度积分，可以同时保留对数奇性、区域几何和零边值约束。本节先说明这种叠加在何种意义下解方程，再从核的边界导数恢复调和测度，最后讨论在一个紧集上怎样分配单位质量才能使相互作用能量最小。

\subsection{Green 势}
\label{b3:subsec:170401}

若有限多个源点 \(w_j\) 的质量为 \(a_j>0\)，叠加势就是 \(\sum_j a_jg_D(z,w_j)\)。将离散质量换成连续或奇异的分布时，有限和自然换成积分。对角上约定 \(g_D(w,w)=+\infty\)，于是原子仍保留点源的无穷奇性。

\Needspace{6\baselineskip}
\begin{definition}\label{b3:def:17-green-potential}
设 \(D\subseteq\mathbb C\) 为 Green 域，\(\mu\) 为 \(D\) 上有限正 Borel 测度。函数
\[
 G_D\mu(z)\coloneqq\int_Dg_D(z,w)\,\mathrm d\mu(w),
 \quad z\in D,
\]
称为 \(\mu\) 的\term[green-shi]{Green 势}[Green potential]，积分按非负扩展积分理解，允许值为 \(+\infty\)。
\end{definition}
\symidx{\(G_D\mu\)：正测度 \(\mu\) 的 Green 势}

核的联合下半连续性保证可测性，所以定义不需要先假定势有限。要把它看成一个分布，还须证明局部可积；下面先处理质量支集留在域内紧集上的情形。

\Needspace{6\baselineskip}
\begin{proposition}[紧支测度的 Green 势]\label{b3:prop:17-green-potential-distribution}
设 \(D\subseteq\mathbb C\) 为 Green 域，有限正 Borel 测度 \(\mu\) 的支集包含于某个 \(K\Subset D\)。则 \(G_D\mu\in L^1_{\mathrm{loc}}(D)\)，并满足分布等式
\begin{equation}\label{b3:eq:17-green-source}
 -\Delta G_D\mu=2\pi\mu.
\end{equation}
具体而言，对每个 \(\varphi\in C_c^\infty(D)\)，有
\[
 \int_DG_D\mu(z)\,(-\Delta\varphi(z))\,\mathrm dA(z)
 =2\pi\int_D\varphi(w)\,\mathrm d\mu(w).
\]
\end{proposition}
\begin{proof}
固定观察点紧集 \(L\Subset D\)。由\cref{b3:prop:17-kernel-regularity}，对数余项在内部紧集上有共同界；远离对角线的核又连续。因此存在 \(C\ge0\)，使
\[
 0\le g_D(z,w)\le C+\log^+\frac1{|z-w|}
 \quad(z\in L,\ w\in K),
\]
其中 \(\log^+t=\max\{\log t,0\}\)。这个控制中的常数与源点 \(w\) 无关。由于
\[
 \int_L\log^+\frac1{|z-w|}\,\mathrm dA(z)
 \le2\pi\int_0^1-s\log s\,\mathrm ds=\frac\pi2,
\]
Tonelli 定理给出
\[
 \int_LG_D\mu\,\mathrm dA
 \le\mu(K)\left(C\operatorname{area}(L)+\frac\pi2\right)<\infty.
\]

对测试函数 \(\varphi\)，取 \(L=\operatorname{supp}\varphi\)，上述估计乘以 \(\|\Delta\varphi\|_\infty\) 便给出换序所需的绝对可积性。每个源点的核满足 \(-\Delta_zg_D(z,w)=2\pi\delta_w\)：这是\cref{b3:prop:15-log-fundamental}与核的调和余项给出的分布等式。因此
\[
 \begin{aligned}
 \int_DG_D\mu(z)(-\Delta\varphi(z))\,\mathrm dA(z)
 &=\int_K\left(\int_Dg_D(z,w)(-\Delta\varphi(z))\,\mathrm dA(z)\right)\mathrm d\mu(w)\\
 &=2\pi\int_K\varphi(w)\,\mathrm d\mu(w).
 \end{aligned}
\]
\end{proof}

局部可积只保证势几乎处处有限。例如 \(G_D\delta_w=g_D(\cdot,w)\) 在 \(w\) 处为正无穷，仍然满足上面的分布方程。对支集不紧的有限测度，本节保留势的扩展积分定义；有关局部可积性或无穷远收敛的进一步结论需要另外的控制。

\ProseParagraphBreak

光滑密度会把对数奇性平均掉。此时还可把分布解恢复成经典解，并证明它连续取零边值。

\Needspace{6\baselineskip}
\begin{proposition}[光滑内部源的零边值解]\label{b3:prop:17-green-poisson-source}
设 \(D\subseteq\mathbb C\) 为有界区域，其边界由有限条互不相交的正则实解析 Jordan 曲线组成，\(f\in C_c^\infty(D;\mathbb R)\)。则
\[
 u_f(z)\coloneqq\frac1{2\pi}\int_Dg_D(z,w)f(w)\,\mathrm dA(w)
\]
在 \(D\) 内光滑，连续延拓到 \(\overline D\)，并且是
\[
 -\Delta u=f\quad\text{于 }D,
 \quad u=0\quad\text{于 }\partial D
\]
在 \(C^2(D)\cap C(\overline D)\) 中的唯一解。函数 \(u_f\) 还在每段边界附近具有光滑调和延拓。
\end{proposition}
\begin{proof}
把 \(f\) 在域外延为零，记 \(K=\operatorname{supp}f\)。核的对数控制保证积分在每个点绝对收敛。写
\[
 g_D(z,w)=-\log|z-w|+h(z,w),
\]
其中 \(h(\cdot,w)\) 在填补 \(z=w\) 的可去点后于全域调和。于是 \(u_f=N_f+H_f\)，其中
\[
 \begin{aligned}
 N_f(z)&=-\frac1{2\pi}\int_{\mathbb C}\log|z-w|\cdot f(w)\,\mathrm dA(w),\\
 H_f(z)&=\frac1{2\pi}\int_Kh(z,w)f(w)\,\mathrm dA(w).
 \end{aligned}
\]
对 \(N_f\) 作换元 \(t=z-w\)，可将各阶导数放到 \(f(z-t)\) 上；当 \(z\) 留在固定紧集时，积分只涉及共同有界的 \(t\) 集合，而 \(\log|t|\) 局部可积。因此 \(N_f\) 光滑，且对数基本解给 \(-\Delta N_f=f\)。对 \(H_f\)，\cref{b3:prop:17-kernel-regularity}的共同余项界与\cref{b3:prop:15-interior-derivatives}给出关于 \(z\) 的各阶导数的共同局部界，故可在积分号下求导，得到 \(H_f\) 光滑且调和。这证明了内部方程。

还须控制 \(z\) 接近边界时的源点积分。选定 \(w_0\in D\)，用内部有限条路径及小圆盘把 \(K\) 与 \(w_0\) 连在一个相对紧的连通邻域中。在足够薄的边界邻域内，观察点 \(z\) 不落入这些圆盘。函数 \(w\mapsto g_D(z,w)\) 在这组圆盘上正且调和，所以\cref{b3:prop:15-harnack-chain}的有限链比较给出
\[
 \sup_{w\in K}g_D(z,w)\le C_Kg_D(z,w_0).
\]
所用圆盘固定，常数不随 \(z\) 改变。右端在 \(z\to\partial D\) 时趋于零，因而
\[
 |u_f(z)|\le\frac{C_K}{2\pi}\|f\|_{L^1(D)}g_D(z,w_0)\longrightarrow0.
\]
这给出连续零边值。源 \(f\) 的支集与边界有正距离，故 \(u_f\) 在边界内侧已经调和；\cref{b3:sec:1701}的解析边界拉直与奇反射使它光滑延拓到每段边界的另一侧。最后，两解之差调和并连续取零边值，最大值原理给出唯一性。
\end{proof}

式\eqref{b3:eq:17-green-source}解释了这里的系数：Green 势本身产生 \(2\pi f\)，方程 \(-\Delta u=f\) 的解须除以 \(2\pi\)。后面的边界表示采用同一个归一化。

\subsection{调和测度}
\label{b3:subsec:170402}

上一章从边界解算子的正性得到调和测度，没有要求边界存在法向量。对于解析边界，Green 核的反射正则性允许把这个抽象测度写成弧长密度。先对具有边界导数的函数证明积分公式，再以一致逼近处理一般连续边值。

\Needspace{6\baselineskip}
\begin{proposition}[Green 第二表示公式]\label{b3:prop:17-green-representation}
设 \(D\subseteq\mathbb C\) 为有界区域，其边界由有限条互不相交的正则实解析 Jordan 曲线组成；实函数 \(u\) 在 \(\overline D\) 的一个开邻域上属于 \(C^2\)。则对每个 \(z\in D\)，
\begin{equation}\label{b3:eq:17-green-representation}
 u(z)=-\frac1{2\pi}\int_{\partial D}u(\xi)\partial_{n_\xi}g_D(\xi,z)\,\mathrm ds(\xi)
 -\frac1{2\pi}\int_Dg_D(\xi,z)\Delta u(\xi)\,\mathrm dA(\xi),
\end{equation}
其中 \(n_\xi\) 是区域的外单位法向，\(\mathrm ds\) 为未归一化弧长测度。
\end{proposition}
\begin{proof}
固定 \(z\)，记 \(g(\xi)=g_D(\xi,z)\)，并从 \(D\) 删去闭圆盘 \(\overline{B(z,\varepsilon)}\)。核在外边界附近的光滑性由\cref{b3:sec:1701}的反射论证保证；在删圆后的闭域上，\(u\) 与 \(g\) 都有 Green 第二公式所需的正则性。由于 \(\Delta g=0\)，有
\[
 -\int_{D\setminus\overline{B(z,\varepsilon)}}g\Delta u\,\mathrm dA
 =\int_{\partial D}u\partial_ng\,\mathrm ds
 +\int_{|\xi-z|=\varepsilon}
       (u\partial_ng-g\partial_nu)\,\mathrm ds.
\]
外边界上 \(g=0\)。内圆的外法向是删圆余域的外法向，指向圆心，所以
\[
 g=-\log\varepsilon+O(1),
 \quad \partial_ng=\varepsilon^{-1}+O(1),
 \quad u(\xi)=u(z)+O(\varepsilon).
\]
内圆第一项的积分趋于 \(2\pi u(z)\)，第二项的绝对值为 \(O(\varepsilon|\log\varepsilon|)\)，趋于零。面积积分也收敛：极点附近的对数可积，远离极点的闭域部分上核有界。令 \(\varepsilon\downarrow0\) 并移项即得公式。
\end{proof}

内圆的法向决定公式中的符号。取 \(u\equiv1\)，面积项消失，边界项须给出一；取零边值函数，边界项消失，面积项则与上一小节的源表示相合。

\Needspace{6\baselineskip}
\begin{proposition}[解析边界上的调和测度密度]\label{b3:prop:17-harmonic-measure-density}
设 \(D\subseteq\mathbb C\) 为有界区域，其边界由有限条互不相交的正则实解析 Jordan 曲线组成。对 \(z\in D\)，定义
\[
 q_z(\xi)\coloneqq-\frac1{2\pi}\partial_{n_\xi}g_D(\xi,z)
 \quad(\xi\in\partial D).
\]
则 \(q_z\) 连续、非负，且 \(\int_{\partial D}q_z\,\mathrm ds=1\)。该域的调和测度满足
\begin{equation}\label{b3:eq:17-harmonic-measure-density}
 \mathrm d\omega_D^z(\xi)=q_z(\xi)\,\mathrm ds(\xi).
\end{equation}
因此，对任意实值连续边界数据 \(\varphi\)，其经典调和解为
\[
 H_\varphi(z)=\int_{\partial D}\varphi(\xi)q_z(\xi)\,\mathrm ds(\xi).
\]
\end{proposition}
\begin{proof}
反射使固定源点的核在边界附近光滑，故 \(q_z\) 连续。沿边界的内法向进入区域时核非负，而边界值为零，因而
\[
 \partial_{n_\xi}g_D(\xi,z)
 =\lim_{t\downarrow0}\frac{g_D(\xi,z)-g_D(\xi-tn_\xi,z)}{t}\le0.
\]
对常数函数使用\cref{b3:prop:17-green-representation}，便得 \(q_z\) 的积分为一。

每个解析边界点都满足局部外锥条件，所以上一章的正则性结果保证连续边值问题可解。为把表示公式用于 \(H_\varphi\)，对每条边界分量选一个实解析的正则周期参数 \(\gamma_j(t)\)，周期取 \(2\pi\)。连续函数 \(\varphi\circ\gamma_j\) 可由实三角多项式 \(p_{j,m}\) 一致逼近；这可用\cref{b3:subsec:160204}中已经证明的 Fej\'er 平均取得。由有限条分量组成的边值 \(\varphi_m(\gamma_j(t))=p_{j,m}(t)\) 因而在整个 \(\partial D\) 一致趋于 \(\varphi\)。记其调和解为 \(u_m\)。

这些特殊边值确实给出边界光滑性。固定一段解析弧，参数 \(\gamma_j\) 的局部全纯逆把它拉直为实轴。三角多项式 \(p_{j,m}(t)\) 是某个全纯函数 \(P_{j,m}(t)\) 在实轴上的实部：将其中的实系数正弦、余弦直接延拓到复参数即可。因此
\[
 v(\xi)=\operatorname{Re}P_{j,m}(\gamma_j^{-1}(\xi))
\]
在弧的邻域调和，并取到 \(u_m\) 的边值。差 \(u_m-v\) 连续取零弧值，奇反射使它调和延拓。这样 \(u_m\) 在每个边界点附近都有光滑延拓；有限覆盖与光滑分片也给出闭域邻域上的 \(C^2\) 延拓。对 \(u_m\) 使用 Green 表示，面积项因调和性消失，得到
\[
 u_m(z)=\int_{\partial D}\varphi_m(\xi)q_z(\xi)\,\mathrm ds(\xi).
\]

最大值原理给
\[
 \|u_m-H_\varphi\|_{C(\overline D)}
 \le\|\varphi_m-\varphi\|_{C(\partial D)}\longrightarrow0.
\]
另一方面，\(q_z\,\mathrm ds\) 是概率测度，故边界积分的误差也不超过同一个一致范数。取极限得到连续数据的表示，再由\cref{b3:thm:16-harmonic-measure}中表示测度的唯一性识别出 \(\omega_D^z\)。
\end{proof}

一般连续边值的调和解不由此获得边界 \(C^1\) 正则性；证明只在三角多项式逼近的解上使用边界导数，再通过最大值原理取极限。导数落在固定内点的 Green 核上，正是这种表示适用于连续数据的原因。

\ProseParagraphBreak

在圆盘 \(D=B(0,R)\) 中，对显式 Green 核沿 \(\xi=Re^{\mathrm it}\) 的径向求导，得到
\[
 q_z(Re^{\mathrm it})
 =\frac{R^2-|z|^2}{2\pi R\,|Re^{\mathrm it}-z|^2}.
\]
乘以弧长元 \(R\,\mathrm dt\)，便恢复圆盘的 Poisson 权重。这里的 \(R\) 因子来自弧长与角参数的区别。

若同时给定 \(f\in C_c^\infty(D;\mathbb R)\) 和 \(\varphi\in C(\partial D;\mathbb R)\)，则唯一的经典解由两项叠加给出：
\[
 u(z)=\int_{\partial D}\varphi(\xi)q_z(\xi)\,\mathrm ds(\xi)
 +\frac1{2\pi}\int_Dg_D(z,w)f(w)\,\mathrm dA(w).
\]
它满足 \(-\Delta u=f\) 和 \(u|_{\partial D}=\varphi\)。边界项分配边界数据，面积项叠加内部源。作为常数核验，圆盘上的函数 \(u(z)=(R^2-|z|^2)/4\) 满足 \(-\Delta u=1\) 且边值为零。它在闭域邻域为 \(C^2\)，所以直接由\cref{b3:prop:17-green-representation}得到
\[
 \int_{B(0,R)}g_D(z,w)\,\mathrm dA(w)
 =\frac\pi2(R^2-|z|^2).
\]
这次源 \(1\) 不具有内部紧支集，使用的是已经适用的 Green 第二表示。中心点还可直接计算 \(2\pi\int_0^R s\log(R/s)\,\mathrm ds=\pi R^2/2\)。

\subsection{平衡观点}
\label{b3:subsec:170403}

现在固定一个紧集 \(K\Subset D\)，允许单位质量在 \(K\) 上重新分配。势在每个源点的值再按源测度平均，便得到总相互作用能量。对角上的无穷值使过度集中付出无穷能量；极小问题因而会区分一个集合是否能容纳有限能量的质量分布。

\Needspace{6\baselineskip}
\begin{definition}\label{b3:def:17-green-energy}
设 \(D\subseteq\mathbb C\) 为 Green 域，\(K\Subset D\) 为非空紧集，\(\mathcal P(K)\) 为支持在 \(K\) 上的 Borel 概率测度全体。对任意支持在 \(K\) 上的有限正 Borel 测度 \(\mu\)，定义其\term[green-nengliang]{Green 能量}[Green energy]，并在概率测度上取下确界：
\[
 I_D(\mu)\coloneqq\iint_{K\times K}g_D(z,w)\,\mathrm d\mu(z)\,\mathrm d\mu(w),
 \quad V_D(K)\coloneqq\inf_{\mu\in\mathcal P(K)}I_D(\mu).
\]
积分允许为正无穷。若 \(V_D(K)<\infty\)，且 \(\mu_K\in\mathcal P(K)\) 满足 \(I_D(\mu_K)=V_D(K)\)，则称 \(\mu_K\) 为 \(K\) 相对于 \(D\) 的\term[pingheng-cedu]{平衡测度}[equilibrium measure]。
\end{definition}
\symidx{\(I_D(\mu)\)：测度 \(\mu\) 的 Green 能量}
\symidx{\(V_D(K)\)：紧集 \(K\) 上概率测度的最小 Green 能量}

\Needspace{6\baselineskip}
\begin{theorem}[平衡测度的存在与变分条件]\label{b3:thm:17-equilibrium-existence}
设 \(D\subseteq\mathbb C\) 为 Green 域，\(K\Subset D\) 为非空紧集。如果 \(K\) 上存在有限 Green 能量的概率测度，则平衡测度 \(\mu_K\) 存在。对每个有限能量概率测度 \(\nu\in\mathcal P(K)\)，有
\begin{equation}\label{b3:eq:17-equilibrium-variation}
 \int_KG_D\mu_K\,\mathrm d\nu\ge I_D(\mu_K)=V_D(K),
\end{equation}
左端允许为正无穷。
\end{theorem}
\begin{proof}
先说明概率测度为何能取到弱收敛子列。紧度量空间 \(K\) 上的 \(C(K;\mathbb R)\) 有可数稠密子集：对每个正整数 \(m\)，选一个有限 \(1/m\)-网 \(\{x_i\}\)，以
\[
 \psi_i(x)=\frac{(2/m-|x-x_i|)_+}{\sum_j(2/m-|x-x_j|)_+}
\]
得到连续分片。分母处处为正，\(\sum_i\psi_i=1\)，且 \(\psi_i(x)>0\) 时 \(|x-x_i|<2/m\)。函数 \(\sum_i f(x_i)\psi_i\) 因一致连续性一致趋于 \(f\)；再将有限个系数换成有理数，所有 \(m\) 的这些组合便给出可数稠密族 \(\{f_\ell\}\)。

给定 \(\mu_n\in\mathcal P(K)\)，各数列 \(\int f_\ell\,\mathrm d\mu_n\) 有界。逐个抽子列并取对角线后，它们都收敛。概率质量为一使一致逼近误差在积分下不增大，故对每个 \(f\in C(K;\mathbb R)\)，积分也收敛。其极限 \(\Lambda(f)\) 是正线性泛函，且 \(\Lambda(1)=1\)。第一卷定理11.4.1的 Riesz--Markov 正泛函表示适用于紧 Hausdorff 空间 \(K\)，给出一个概率测度 \(\mu\)，使 \(\Lambda(f)=\int f\,\mathrm d\mu\)。于是选出的子列弱收敛于 \(\mu\)。

接着处理核的奇性。令 \(g_M=\min\{g_D,M\}\)，其中 \(M>0\)，对角值取 \(M\)。由\cref{b3:prop:17-kernel-regularity}，核在对角外连续，而内部紧集上的共同余项界给出
\[
 g_D(z,w)\ge-\log|z-w|-C_K.
\]
因此靠近对角线时 \(g_M\) 恒等于 \(M\)，从而 \(g_M\in C(K\times K)\)。

若 \(\mu_n\) 弱收敛于 \(\mu\)，则它们的积测度也对连续测试函数收敛。为看清这里没有遗漏双变量测试，给定 \(F\in C(K\times K)\)，用上面的细分片构造
\[
 F_m(z,w)=\sum_{i,j}F(x_i,x_j)\psi_i(z)\psi_j(w).
\]
一致连续性给 \(\|F_m-F\|_\infty\to0\)，每项的积分则分解为两个弱收敛积分的乘积。先固定 \(m\) 令 \(n\to\infty\)，再令 \(m\to\infty\)，便得到积测度的收敛。特别地，
\[
 \iint g_M\,\mathrm d\mu\,\mathrm d\mu
 =\lim_n\iint g_M\,\mathrm d\mu_n\,\mathrm d\mu_n
 \le\liminf_n I_D(\mu_n).
\]
再让 \(M\uparrow\infty\)，单调收敛给 \(I_D(\mu)\le\liminf_n I_D(\mu_n)\)。对一列能量趋于 \(V_D(K)\) 的概率测度取上述子列，就得到达到有限极小值的 \(\mu_K\)。

最后记 \(V=I_D(\mu_K)\)、\(J=\int G_D\mu_K\,\mathrm d\nu\)。若 \(J=\infty\)，所需不等式已经成立。若 \(J<\infty\)，对 \(0<t<1\) 取合法概率测度 \(\mu_t=(1-t)\mu_K+t\nu\)。正核的 Tonelli 换序与对称性给
\[
 I_D(\mu_t)=(1-t)^2V+2t(1-t)J+t^2I_D(\nu)\ge V.
\]
减去 \(V\)、除以 \(2t\)，再令 \(t\downarrow0\)，得到 \(J\ge V\)。各量在这一步均有限，所以没有对无穷能量作形式求导。
\end{proof}

变分条件比较的是对竞争测度的平均势。它不直接断言势在紧集的每一点都相等；特别是单点质量可能具有无穷能量，不能任取 \(\nu=\delta_z\) 再把有限能量变分公式当作逐点等式。一般集合上的精细平衡势结论还涉及零容量例外集，本节只在下面的圆盘例子中求出处处可核验的势。

\subsection{容量的联系}
\label{b3:subsec:170404}

能量越小，单位质量越容易分散。因此用最小能量的倒数衡量紧集在给定外域中的容纳能力。这个量保留外域：相同的内部集合，离零势边界的距离不同，其 Green 能量也不同。

\Needspace{6\baselineskip}
\begin{definition}\label{b3:def:17-green-capacity}
设 \(D\subseteq\mathbb C\) 为 Green 域，\(K\Subset D\) 为非空紧集。其\term[green-rongliang]{Green 容量}[Green capacity]定义为
\[
 \operatorname{Cap}_D(K)\coloneqq\frac1{V_D(K)},
 \quad \frac1{+\infty}=0.
\]
\end{definition}
\symidx{\(\operatorname{Cap}_D(K)\)：紧集 \(K\) 相对于 Green 域 \(D\) 的 Green 容量}

这个倒数不会因分母为零而失去意义。正连续函数 \(\min\{g_D,1\}\) 在 \(K\times K\) 上有严格正的最小值 \(a_K\)，故所有概率测度的能量至少为 \(a_K\)。若 \(K\) 只是有限点集，每个概率测度都有原子，对角奇性使其能量为无穷，因而容量为零。集合增大使可选测度增多，所以容量不减；外域增大使核增大，最小能量不减，容量便不增。这两个方向与上一章的相对 Dirichlet 容量一致。

\ProseParagraphBreak

下面完整计算同心闭圆盘的容量。先建立所需的能量非负性；它不能仅从核逐点为正推出，因为参与比较的是符号测度。

\Needspace{6\baselineskip}
\begin{lemma}[圆盘上的符号 Green 能量]\label{b3:lem:17-disk-energy-positive}
设 \(D=B(0,R)\)、\(R>0\)，\(\eta\) 是支持在某个内部紧集上的有限实符号 Borel 测度，并且
\[
 \iint g_D(z,w)\,\mathrm d|\eta|(z)\,\mathrm d|\eta|(w)<\infty.
\]
则绝对收敛的符号能量满足
\[
 I_D(\eta)\coloneqq\iint g_D(z,w)\,\mathrm d\eta(z)\,\mathrm d\eta(w)\ge0.
\]
\end{lemma}
\begin{proof}
先设 \(\mathrm d\eta=f\,\mathrm dA\)，其中 \(f\in C_c^\infty(D;\mathbb R)\)。令 \(V=G_D\eta\)。\cref{b3:prop:17-green-poisson-source}给 \(-\Delta V=2\pi f\)、零边值以及闭域上的光滑性。分部积分的边界项因 \(V=0\) 消失，于是
\[
 I_D(\eta)=\int_DVf\,\mathrm dA
 =\frac1{2\pi}\int_DV(-\Delta V)\,\mathrm dA
 =\frac1{2\pi}\int_D|\nabla V|^2\,\mathrm dA\ge0.
\]

一般情形下，取非负径向函数 \(\rho\in C_c^\infty(B(0,1))\)，使 \(\int\rho\,\mathrm dA=1\)，并令 \(\rho_\varepsilon(a)=\varepsilon^{-2}\rho(a/\varepsilon)\)。把 \(\eta\) 延为平面上的测度，记其磨光为 \(\eta_\varepsilon\)。当 \(\varepsilon\) 充分小时，磨光支集仍在同一个内部紧集上，因而已知 \(I_D(\eta_\varepsilon)\ge0\)。换序把这个能量写成
\[
 I_D(\eta_\varepsilon)=\iint g_\varepsilon(z,w)\,\mathrm d\eta(z)\,\mathrm d\eta(w),
\]
其中
\[
 g_\varepsilon(z,w)
 =\iint g_D(z+a,w+b)\rho_\varepsilon(a)\rho_\varepsilon(b)
            \,\mathrm dA(a)\,\mathrm dA(b).
\]
Green 核在每个变量中超调和，包括对角上的正无穷值。径向非负磨光是圆周平均的非负加权平均，因此超平均不等式可逐变量使用，得到
\[
 0\le g_\varepsilon(z,w)\le g_D(z,w).
\]
所有平移点都留在 \(D\) 内，故这里不涉及把核延到域外。

由对角外的联合连续性，\(g_\varepsilon(z,w)\to g_D(z,w)\) 当 \(z\ne w\)。有限能量假设还排除了 \(|\eta|\) 的原子：若一点质量为 \(a>0\)，对角贡献就使能量至少为 \(a^2M\)，对任意 \(M>0\) 均成立。于是对角线在 \(|\eta|\otimes|\eta|\) 下零测，因为其测度等于 \(\int|\eta|(\{z\})\,\mathrm d|\eta|(z)=0\)。共同控制 \(g_D\) 对这个积测度可积，支配收敛便给
\[
 I_D(\eta)=\lim_{\varepsilon\downarrow0}I_D(\eta_\varepsilon)\ge0.
\]
同一可积控制也使前面的换序对符号测度合法。
\end{proof}

\Needspace{6\baselineskip}
\begin{proposition}[同心闭圆盘的平衡测度与容量]\label{b3:prop:17-disk-capacity}
设 \(0<r<R\)，\(D=B(0,R)\)、\(K=\overline{B(0,r)}\)。令 \(\sigma_r\) 为圆周 \(|w|=r\) 上的均匀概率测度，即 \(\mathrm d\sigma_r(re^{\mathrm it})=\mathrm dt/(2\pi)\)。则
\begin{equation}\label{b3:eq:17-circle-equilibrium-potential}
 G_D\sigma_r(z)=
 \begin{cases}
  \log(R/r),&|z|\le r,\\
  \log(R/|z|),&r<|z|<R.
 \end{cases}
\end{equation}
测度 \(\sigma_r\) 是 \(K\) 的平衡测度，并且
\[
 V_D(K)=\log(R/r),
 \quad \operatorname{Cap}_D(K)=\frac1{\log(R/r)}.
\]
\end{proposition}
\begin{proof}
圆盘核可写成
\[
 g_D(z,w)=\log R+\log\left|1-\frac{\overline wz}{R^2}\right|-\log|z-w|.
\]
若 \(|a|<1\)，由 \(\log(1-ae^{\mathrm it})\) 的幂级数取实部并平均，得到
\[
 \frac1{2\pi}\int_0^{2\pi}\log|1-ae^{\mathrm it}|\,\mathrm dt=0.
\]
这里级数绝对一致收敛，逐项积分有依据。分子项总满足 \(|a|=r|z|/R^2<1\)，故其平均为零。分母项在 \(|z|<r\) 时提出 \(r\)，在 \(|z|>r\) 时提出 \(|z|\)，同一公式分别给出平均 \(\log r\) 与 \(\log|z|\)。

在 \(|z|=r\) 上还须处理圆周积分的一个对数奇点。对 \(1/2\le s<1\)，有
\[
 |1-se^{\mathrm it}|^2=(1-s)^2+4s\sin^2(t/2)
 \ge2\sin^2(t/2),
 \quad |1-se^{\mathrm it}|\le2.
\]
因此其对数绝对值被 \(C+|\log|\sin(t/2)||\) 控制，后者在圆周可积。令 \(s\uparrow1\)，支配收敛将平均零公式传到 \(|a|=1\)。由此圆周上的分母平均也等于 \(\log r\)，证明了包括界面点在内的势公式。记 \(c=\log(R/r)\)，于是 \(I_D(\sigma_r)=c\)。

取任意 \(\nu\in\mathcal P(K)\)。若其能量无穷，下界 \(I_D(\nu)\ge c\) 已成立；以下设其能量有限。由于 \(G_D\sigma_r=c\) 在整个 \(K\) 成立，交叉能量为
\[
 \iint g_D(z,w)\,\mathrm d\nu(z)\,\mathrm d\sigma_r(w)=c.
\]
令 \(\eta=\nu-\sigma_r\)。不等式 \(|\eta|\le\nu+\sigma_r\) 与正核给
\[
 I_D(|\eta|)\le I_D(\nu)+2c+I_D(\sigma_r)<\infty.
\]
因此\cref{b3:lem:17-disk-energy-positive}适用，并给出
\[
 0\le I_D(\eta)
 =I_D(\nu)-2c+I_D(\sigma_r)=I_D(\nu)-c.
\]
下界在 \(\sigma_r\) 处达到，所以它是平衡测度，最小能量及容量公式随之成立。
\end{proof}

同一势也说明本节与上一章的归一化差异。令 \(v=G_D\sigma_r/\log(R/r)\)，则 \(v\) 在内圆盘恒为一，在圆环中为 \(\log(R/|z|)/\log(R/r)\)，且在外圆周为零。其 Dirichlet 能量为
\[
 \int_D|\nabla v|^2\,\mathrm dA
 =\frac{2\pi}{\log^2(R/r)}\int_r^R\frac{\mathrm ds}{s}
 =\frac{2\pi}{\log(R/r)}.
\]
这个分段函数本身不是光滑紧支试函数，\cref{b3:prop:16-annulus-capacity}已用截断和磨光把它逼近为合法竞争函数，并证明了匹配的下界。因此在这个例子中，
\[
 \operatorname{cap}(K,D)=2\pi\operatorname{Cap}_D(K).
\]
因子 \(2\pi\) 同时出现在点源方程、势的能量恒等式和两种容量之间，均由 \(-\Delta(-\log|z|)=2\pi\delta_0\) 决定。一般紧集上两种容量的对应还需要进一步的势论，本节的计算已经说明为什么不能只因名称相近就混用数值。

% !TeX root = ../../Book3.tex
\section{多连通区域}
\label{b3:sec:1705}

多连通域的调和函数仍可由边值确定，但共轭的周期多了有限个独立方向。边界分量的调和测度提供了调整这些周期的自然函数。

\subsection{周期}
\label{b3:subsec:170501}

调和函数 \(u\) 的周期为 \(\int_\gamma(-u_y\,\mathrm{d}x+u_x\,\mathrm{d}y)\)。闭形式使周期只依赖曲线同调类。若区域由 \(m+1\) 条互不相交实解析 Jordan 曲线围成，则绕 \(m\) 个内孔的曲线构成独立周期方向；任一闭曲线的周期为这些基本周期的整数线性组合。

\subsection{调和测度}
\label{b3:subsec:170502}

令 \(\omega_j\) 为第 \(j\) 个内边界取一、其余边界取零的调和解。其边界数据在各分量上连续，各边界弧可拉直并取局部线性障碍，所以 Perron 理论保证存在；减去该分量上的常数后作调和反射，还保证解在边界附近光滑。它们满足 \(0<\omega_j<1\)，连同外边界权重 \(\omega_0\) 有 \(\sum_{j=0}^m\omega_j=1\)。因此常数边界数据 \(c_j\) 的解为 \(\sum c_j\omega_j\)。

环域 \(r<|z|<R\) 的内边界权重是 \(\log(R/|z|)/\log(R/r)\)，其共轭周期为 \(-2\pi/\log(R/r)\)。这给出可计算的周期调整系数。

\subsection{调和共轭的障碍}
\label{b3:subsec:170503}

沿正向边界，\(\omega\) 的积分等于调和函数的外法向通量。令 \(A_{ij}=\int_D\nabla\omega_i\cdot\nabla\omega_j\,\mathrm{d}A\)。Green 第一公式把它等于第 \(i\) 个内边界上的 \(\omega_j\) 通量。

矩阵 \(A\) 对称正定：若 \(\sum c_j\omega_j\) 的梯度能量为零，则该组合恒常；外边界值为零使常数零，再看各内边界值知全部 \(c_j=0\)。因此基本周期可以由 \(\omega_j\) 的线性组合任意调整，拓扑障碍有明确的有限维形式。

\subsection{典范表示}
\label{b3:subsec:170504}

\begin{proposition}\label{b3:prop:17-period-decomposition}
在上述有限连通实解析边界有界域中，每个连续到边界的调和函数可写为
\[
u=\operatorname{Re}F+\sum_{j=1}^m c_j\omega_j,
\]
其中 \(F\) 单值全纯，\(c_j\) 唯一；\(F\) 只差纯虚常数。
\end{proposition}
\begin{proof}
用正定通量矩阵解线性方程，选择 \(c_j\) 消除 \(u\) 在全部基本曲线上的周期。剩余调和函数所有闭曲线周期为零，由\cref{b3:thm:15-global-conjugate}得到单值共轭和 \(F\)。若两组系数都可行，其差的周期向量为零，矩阵可逆使差为零；共轭唯一性给最后结论。
\end{proof}


\ProseParagraphBreak

基本周期的方向须统一。沿区域正向边界行进时，区域在左侧，外边界逆时针、内边界顺时针。若切向为 \((x',y')\)，外法向为 \((y',-x')\)，则
\[
(-u_y\,\mathrm{d}x+u_x\,\mathrm{d}y)(\gamma')=\nabla u\cdot n\,\mathrm{d}s.
\]
所以通量矩阵直接给的是顺时针内边界周期；若基本闭曲线统一选逆时针绕孔，则矩阵前须加负号。环域内边界权重的逆时针周期是 \(-2\pi/\log(R/r)\)，而其沿区域正向内边界通量为 \(2\pi/\log(R/r)\)，正好与正定能量相符。

\ProseParagraphBreak

有限维拓扑描述可以通过割线实现。从每个内边界取一条互不相交的简单弧连接到外边界，沿这些弧割开区域后得到单连通域。在割开的域中闭形式有原函数。沿第 \(j\) 条割线两侧的原函数差为绕该孔的周期，且因导数相同而为常数；所有这些差均零时，原函数才能连续拼回原区域。任意闭曲线跨越各条割线的有向次数为整数，其积分等于对应周期的整数和。这同时解释基本周期足够、为何有 \(m\) 个独立方向，以及周期清零后共轭的单值性。

\ProseParagraphBreak

典范表示中的 \(c_j\) 不是边界数据的简单平均。它由整个调和函数的周期决定。环域中取 \(u(z)=\log|z|\)，其周期为 \(2\pi\)，内边界权重周期为 \(-2\pi/\log(R/r)\)，故 \(c=-\log(R/r)\)，余项是常数 \(\log R\)。取 \(u(z)=\operatorname{Re}z\)，周期为零，所以 \(c=0\)，尽管两个边界上的函数值通常不相同。边界权重分解只适合各分量上常数的边值，而周期分解适合一般调和函数；两种线性分解的系数含义不同。

\ProseParagraphBreak

当边界孔数增加，矩阵的大小也增加；如果有无穷多个孔，当前有限维正定论证并未给出一个自动收敛的无穷级数表示。系数、能量和周期的求和需要额外估计。即使孔数有限，若边界粗糙到通常法向导数不存在，也应改用内部周期和弱能量方法，不能继续逐点积分边界法向导数。本节的实解析边界有界范围正是为了使拓扑方向和通量方向能够用同一个有限矩阵联系起来。

% 本章小结（不计入编号）
\begin{chaptersummary}
有界规则域的 Green 核由边值解构造，一般域则由单调穷竭识别其最小正核。正核不存在的平面反例说明无穷远约束不可省略，也说明覆盖共形类型不等于 Green 存在性。核叠加产生势，能量极小化产生容量。有限连通域的通量矩阵把共轭障碍化成可解的线性方程。
\end{chaptersummary}


\BookExercises

% 本题组自拟；各题新增训练动作见题意，未复制第三方原题。
\begin{exercise}\label{b3:ex:17-draft-01}
核验单位圆盘 \(g(z,0)\) 的 Riesz 测度符号与边界值。
\end{exercise}
\begin{solution}
\(g(z,0)=-\log|z|\)，所以 \(-\Delta g=2\pi\delta_0\)，边界零。作为超调和函数，其 \(\Delta g\) 非正；不能把它的源质量按次调和 Riesz 测度记成正。
\end{solution}

\begin{exercise}\label{b3:ex:17-draft-02}
求上半平面中 \(g(\mathrm{i}y,\mathrm{i}a)\)，\(a,y>0\)、\(a\ne y\)，并考察 \(y\to0\)。
\end{exercise}
\begin{solution}
为 \(\log((y+a)/|y-a|)\)，正，因为 \(y+a>|y-a|\)。当 \(y\to0\) 时比值趋一，核趋零；当 \(y\to a\) 时趋正无穷。
\end{solution}

\begin{exercise}\label{b3:ex:17-draft-03}
设 \(D\subset B(0,R)\) 且 \(0\in D\)。证明 \(g_D(z,0)\le\log(R/|z|)\)。
\end{exercise}
\begin{solution}
右侧为包围圆盘的正 Green 核，限制到 \(D\) 保持奇性，最小性给结论。上界同时证明 \(D\) 的穷竭核有限。
\end{solution}

\begin{exercise}\label{b3:ex:17-draft-04}
解释为何在上半平面给正 Green 核加 \(cy\)、\(c>0\) 后不再是最小 Green 核。
\end{exercise}
\begin{solution}
新增项调和、边界零、在极点无奇性，故局部条件仍成立；但原核逐点严格较小，因而新增核不满足最小性。该例反驳无界域仅以通常边界值刻画唯一性。
\end{solution}

\begin{exercise}\label{b3:ex:17-draft-05}
求环域两边界分别取 \(A,B\) 的径向调和解，并给其共轭周期。
\end{exercise}
\begin{solution}
解为 \(A+(B-A)\log(|z|/r)/\log(R/r)\)，周期为 \(2\pi(B-A)/\log(R/r)\)。只有 \(A=B\) 才有单值共轭。
\end{solution}

\begin{exercise}\label{b3:ex:17-draft-06}
若有限正测度 \(\mu\) 有一个原子，证明其 Green 能量无穷。
\end{exercise}
\begin{solution}
若 \(\mu(\{a\})=c>0\)，双积分在 \((a,a)\) 上有质量 \(c^2\)，而 \(g_D(a,a)=+\infty\)。用截断核 \(\min(g_D,N)\) 下界能量至少为 \(c^2N\)，令 \(N\to\infty\)。
\end{solution}

\begin{exercise}\label{b3:ex:17-draft-07}
证明 \(\sum_{j=0}^m\omega_j=1\)，并说明为何这不造成内边界通量矩阵奇异。
\end{exercise}
\begin{solution}
左侧在每个边界分量取一，唯一性使其等于一。内边界矩阵只含 \(j=1,\ldots,m\)，这些组合在外边界取零，故不能等于非零常数，正定性不受全体权重关系影响。
\end{solution}

\begin{exercise}\label{b3:ex:17-draft-08}
在单位圆盘中，令 \(\mu\) 为半径 \(r\) 圆周的均匀概率测度。计算 \(G\mu(0)\) 与 \(I(\mu)\)。
\end{exercise}
\begin{solution}
中心核为 \(-\log r\)，故 \(G\mu(0)=-\log r\)。显式核中分子的对数圆周平均为零，而 \(\log|re^{\mathrm{i}t}-re^{\mathrm{i}s}|\) 平均为 \(\log r\)，所以圆周上势同为 \(-\log r\)，能量也为该值；对角单点在弧长测度中零测， 对数奇性 可积。
\end{solution}

% !TeX root = ../../Book3.tex
\chapter{Poisson–Jensen 公式与零点分布}
\label{b3:ch:18}
\BookChapterNumber{b3:ch:18}
\begin{chapterabstract}
零点的重数决定对数模的点源质量，零点的位置决定它对圆周平均的贡献。本章先建立整数计数、积分计数及其在全纯扰动下的稳定性，再从圆周对数核和 Green 核推导 Jensen 与 Poisson--Jensen 公式。中心零点、观察圆周上的零点及亚纯函数的极点分别处理。随后以增长界控制零点密度，证明有界圆盘全纯函数的零点判据及 Blaschke 乘积构造，最后把这些公式放入次调和函数的 Riesz 测度框架，并证明非零全纯极限下对数模的局部可积收敛。
\end{chapterabstract}

\begin{learninggoals}
\begin{itemize}
\item 按重数记录零点，区分开、闭圆盘计数与积分计数，并用 Rouch\'e 定理追踪扰动中的零点簇。
\item 推导 Jensen 与 Poisson--Jensen 公式，核对中心系数、Green 项的符号以及边界对数奇性的可积性。
\item 把最大模或圆周幂平均的增长界转为零点限制，识别两半径估计中半径间隔的作用。
\item 证明 Blaschke 条件的必要性与充分性，控制局部乘积收敛及除去零点后的有界余因子。
\item 由对数模的 Riesz 测度解释径向平均和零点稳定性，区分全纯函数的正零点测度与亚纯函数的符号零极点测度。
\end{itemize}
\end{learninggoals}

若 \(f(0)=1\)，且 \(|f|\) 在半径 \(R\) 的圆周上不超过十，那么较小圆盘 \(|z|\le r<R\) 中不能放入任意多的零点。一个位于 \(a\) 的零点对平均对数贡献 \(\log(R/|a|)\)，所以较小圆盘中的每个零点至少贡献 \(\log(R/r)\)。要把这个想法变成可靠估计，既要记录重数，也要说明圆周上的对数如何积分；若原点本身是零点，还须先取出它的幂次，不能直接写下 \(\log|f(0)|\)。

\ProseParagraphBreak
\secref{b3:sec:1801}先把零点整理为局部有限的整数质量。整数计数随半径跳跃，积分计数则给每个非中心零点加上对数权重，在对数半径上连续而分段线性。Rouch\'e 定理保持一个小围圈中的总重数，缩小这些围圈便能追踪合并或分裂的零点簇。这种内部稳定性容许质量移向边界，但要求极限函数不恒为零。

\ProseParagraphBreak
\secref{b3:sec:1802}把计数与函数大小联系起来。先算一个对数核的圆周平均，再有限分解零点，便得到含中心修正的 Jensen 公式。任意内部观察点需要用 Green 核替代中心的对数权重；核的负 Laplace 与 \(\log|f|\) 的正点质量恰好抵消。边界零点仍可处理，但必须保留可积对数奇性，不能把边界无零点版本原样套用。两半径估计和 Blaschke 条件随后说明：增长限制不仅约束零点数量，也约束它们接近边界的速度。

\ProseParagraphBreak
\secref{b3:sec:1803}说明这些计算为何属于势论。零点对应正 Riesz 质量，极点则对应负质量；一般次调和函数的质量还可以分布在曲线或面积上。局部一致收敛的全纯函数在移动零点附近未必有一致收敛的对数，但有限零点分解与对数核的可积性足以给出局部 \(L^1\) 收敛。阅读这部分时应分别留意点态值、面积积分和测试函数所记录的信息。

\ProseParagraphBreak
本章使用前面的局部零点分解、辐角原理与 Rouch\'e 定理，以及\chapref{b3:ch:15}的次调和函数和 Riesz 测度、\chapref{b3:ch:17}的正 Green 核。有关整函数增长与零点分布的进一步阅读，可参考 Levin 的《Lectures on Entire Functions》\cite{Levin1996Lectures}。本章的平均公式与有界零点构造在正文中独立建立，后续整函数乘积理论将继续处理给定零点与全局增长之间的关系。

% !TeX root = ../../Book3.tex
\section{零点计数}
\label{b3:sec:1801}

辐角原理用一条闭曲线计算零点总数，但没有直接显示零点到圆心的距离。对数模的圆周平均会把这项几何信息补进计数。

\subsection{重数}
\label{b3:subsec:180101}

非恒零全纯函数在零点 \(a\) 可唯一写成 \(f(z)=(z-a)^m g(z)\)，其中 \(g(a)\ne0\)，\(m\) 是零点重数。后面的计数均按重数重复零点。紧子集上只有有限个零点；否则有内部聚点，恒等定理迫使函数恒零。

例如 \(f(z)=z^2(z-1)^3\) 在原点和一的质量分别为二、三。零点位置相同而重数不同，会改变对数模的奇性强度；按不同零点个数而不计重数不能恢复解析信息。

\subsection{计数函数}
\label{b3:subsec:180102}

对 \(f(0)\ne0\)，令 \(n(r)\) 为 \(|a|\le r\) 内按重数的零点数，定义积分计数
\[
N(r)=\int_0^r\frac{n(t)}t\,\mathrm{d}t
=\sum_{|a|<r}\log\frac r{|a|}.
\]
第二式由有限和换序得到。零点在半径边界时贡献零，故开闭圆盘约定不影响 \(N\)。若原点重数为 \(m\)，先除去 \(z^m\)，再单独保留 \(m\log r\) 项。\(n\) 是阶梯函数，\(N\) 却连续；每个零点从它进入圆盘时开始累计对数权重。

\subsection{辐角原理再述}
\label{b3:subsec:180103}

边界无零点时，\cref{b3:thm:03-argument-principle}给
\[
n(r)=\frac1{2\pi\mathrm{i}}\int_{|z|=r}\frac{f'(z)}{f(z)}\,\mathrm{d}z.
\]
这里要求 \(f\) 在闭圆盘邻域全纯，或者先在稍小半径应用；边界上有零点时积分不定义。辐角原理显示整数计数的拓扑稳定性，积分计数则更适合平均与增长估计。后面不会把 \(n(r)\) 的跳跃误当作对数平均本身的跳跃。


\ProseParagraphBreak

计数函数的两个版本表达不同信息。若半径内有 \(N\) 个零点全位于 \(|a|=r/2\)，它们给 \(n(r)=N\)、\(N(r)=N\log2\)；若相同个数全位于 \(|a|=(1-\varepsilon)r\)，整数计数仍为 \(N\)，积分计数却为 \(-N\log(1-\varepsilon)\)，随 \(\varepsilon\downarrow0\) 趋零。紧贴观察圆周的零点对该半径的平均影响较小，但一旦把半径扩大，它们会继续累积权重。因此从 \(N(r)\) 控制 \(n(r)\) 必须另取更大的半径，不能同半径直接比较而没有额外系数。

\ProseParagraphBreak

换一个圆心可以避开原点零点。对内部点 \(b\) 且 \(f(b)\ne0\)，在 \(B(b,R)\) 使用变量 \(\zeta=z-b\)，得到完全相同的计数和公式。如果要在几个不同圆心比较，应分别保留 \(\log|f(b)|\)；这个中心常数可能很负，会使零点上界较弱。例如函数在 \(b\) 非零但已非常接近一个零点时，\(|f(b)|\) 小，计数估计仍成立却不一定有效。归一化 \(f(b)=1\) 可消除常数项，但乘法归一化并不改变零点本身。

\ProseParagraphBreak

对含原点重数的函数，用修正积分计数
\[
N_f(r)=m\log r+\sum_{0<|a|<r}m_a\log(r/|a|)
\]
十分方便。它在小于一的半径可能为负，因为 \(m\log r<0\)；这不是非原点零点贡献出现负值，而是已单独记录的中心奇性。若只需要非负零点数的估计，先除去原点幂再应用公式更稳妥。所有符号版本必须在中心奇性项上保持同一约定，尤其进入亚纯函数的零点、极点差计数时不能重复加上原点项。

\ProseParagraphBreak

零点对全纯扰动的稳定性有一个比总数更精细的表述。设 \(f_j\to f\) 在区域 \(D\) 内局部一致，且 \(f\not\equiv0\)。固定一个内部零点 \(a\)，取小圆盘使闭盘中没有其他极限零点，边界上的 \(|f|\) 有正下界。Rouch\'e 定理说明：充分大的 \(j\) 在盘内恰有 \(m_a\) 个零点，按重数计。把圆盘继续缩小，仍有同样结论，所以这些零点全向 \(a\) 集中。紧集上远离极限零点的部分，由 \(|f|\) 的正下界可知最终没有新的零点。
\ProseParagraphBreak
于是离散测度 \(\mu_j=\sum m_{j,a}\delta_a\) 在内部弱收敛到 \(\mu=\sum m_a\delta_a\)：对紧支集连续测试函数，只需用有限个互不相交小圆盘包住相关极限零点，每盘总质量相同，测试函数在盘内的振幅任意小，其余部分最终无质量。这里弱收敛只针对内部紧支集；零点可以逃向边界或无穷远。例如 \(f_j(z)=1-z/j\) 在全平面局部一致趋于一，唯一零点趋向无穷远。极限函数没有零点，与每个近似函数有一个零点并不矛盾。
\ProseParagraphBreak
假设极限不恒零是必要的。\(f_j(z)=e^{-j}\sin(jz)\) 在足够窄的固定水平带中局部一致趋于零，但零点间距趋于零，任何固定实区间内的零点数都趋向无穷。恒等为零的函数不再拥有离散的零点测度，不能给它随意指定一个有限重数来延续上面的结论。

% !TeX root = ../../Book3.tex
\section{Jensen 公式}
\label{b3:sec:1802}

除掉圆盘内的零点，剩余函数的对数模调和。中心平均等式因此会留下每个零点的一个对数距离项，这就是 Jensen 公式。

\subsection{Jensen 公式}
\label{b3:subsec:180201}

\term[jensen-gongshi]{Jensen 公式}[Jensen's formula]的归一化以圆周平均 \(\mathrm{d}t/(2\pi)\) 为准。
\begin{theorem}[Jensen]\label{b3:thm:18-jensen}
设 \(f\) 在闭圆盘 \(\overline{B(0,R)}\) 邻域全纯，\(f(0)\ne0\)，边界无零点。内零点按重数记为 \(a_1,\ldots,a_N\)。则
\[
\frac1{2\pi}\int_0^{2\pi}\log|f(Re^{\mathrm{i}t})|\,\mathrm{d}t
=\log|f(0)|+\sum_{j=1}^N\log\frac R{|a_j|}.
\]
\end{theorem}
\begin{proof}
令 \(g=f/\prod(z-a_j)\)，它在闭圆盘无零点，故 \(\log|g|\) 调和。每个因子满足圆周平均 \(\log|Re^{\mathrm{i}t}-a|=\log R\)：写成 \(\log R+\operatorname{Re}\log(1-(a/R)e^{-\mathrm{i}t})\)，收敛级数的每个非零角频率平均零。于是边界平均为 \(\log|g(0)|+N\log R\)，而 \(\log|g(0)|=\log|f(0)|-\sum\log|a_j|\)，得公式。
\end{proof}
边界零点的对数奇性可积；用相邻半径与局部零点分解可把公式延伸到这样的半径。原点为零点时，用 \(f=z^m g\)、\(g(0)\ne0\)，中心常数改成 \(\log|g(0)|\)，并增加 \(m\log R\)。

\subsection{Poisson–Jensen 公式}
\label{b3:subsec:180202}

\begin{theorem}[Poisson--Jensen]\label{b3:thm:18-poisson-jensen}
在上述条件下，对非零点 \(z\in B(0,R)\)，有
\[
\log|f(z)|=\frac1{2\pi}\int_0^{2\pi}P_R(z,t)\log|f(Re^{\mathrm{i}t})|\,\mathrm{d}t
-\sum_jg_{B(0,R)}(z,a_j).
\]
\end{theorem}
\begin{proof}
加上各个正 Green 核后，函数 \(\log|f|+\sum_jg(\cdot,a_j)\) 在每个零点的 \(m\log|z-a_j|\) 奇性完全抵消，成为调和函数。Green 核边界零，所以边值仍为 \(\log|f|\)。圆盘 Poisson 唯一性给表示。零点处两侧均按极限解释为负无穷。
\end{proof}
亚纯函数的版本在每个极点添加对应的正 Green 项：\(\log|f|\) 的 Laplace 测度为零点质量减极点质量。中心公式因此是零点积分计数减极点积分计数。

\subsection{零点分布估计}
\label{b3:subsec:180203}

\begin{corollary}\label{b3:cor:18-zero-estimate}
若 \(f\) 在 \(B(0,R)\) 全纯，\(f(0)\ne0\)，\(M(S)=\max_{|z|=S}|f(z)|\)，则 \(0<r<S<R\) 时
\[
n(r)\log\frac Sr\le\log M(S)-\log|f(0)|.
\]
\end{corollary}
\begin{proof}
半径 \(r\) 内每个零点向半径 \(S\) 的 Jensen 和贡献至少 \(\log(S/r)\)。平均对数不超过最大对数，代入公式即得。边界零点用稍移半径处理。
\end{proof}
圆盘上有界全纯函数的零点满足 Blaschke 条件 \(\sum(1-|a_j|)<\infty\)：令 \(S\uparrow1\)，Jensen 和有界，再用 \(-\log t\ge1-t\)。有界性因而不允许零点过快地向整个圆周拥挤。


\ProseParagraphBreak

边界零点的延伸可写得更明确。设在 \(|a|=R\) 有一个因子 \((z-a)^m\)，余函数在该点不为零。边界局部参数 \(t=t_0+h\) 时，\(|Re^{\mathrm{i}t}-a|\) 与 \(|h|\) 相比，故 \(\log\) 奇性由常数加 \(m\log|h|\) 控制，而 \(\int_{-\delta}^{\delta}|\log|h||\,\mathrm{d}h<\infty\)。其平均仍为 \(\log R\)：先在 \(|a|<R\) 得到等式，再使 \(a\) 沿半径趋向边界，局部负部的积分可由该对数估计一致控制。Jensen 和中该零点贡献 \(\log(R/|a|)=0\)，因此公式在边界零点半径保持正确。这里需要的是对数的积分收敛，不能只凭逐点极限使用有界收敛定理。

\ProseParagraphBreak

一阶函数已把公式的几何意义展示出来。对 \(f(z)=1-az\)、\(a\ne0\)，有
\[
\frac1{2\pi}\int_0^{2\pi}\log|1-aRe^{\mathrm{i}t}|\,\mathrm{d}t
=\max(0,\log(|a|R)).
\]
当 \(|a|R<1\) 时，圆盘无零点，平均等于中心对数零；当 \(|a|R>1\) 时，一个零点进入圆盘，平均变成 \(\log(|a|R)\)。临界半径仍连续，只有对 \(\log R\) 的斜率发生跳跃。更一般地，对多项式 \(c\prod(z-a_j)\)，边界平均为 \(\log|c|+\sum_j\log\max(R,|a_j|)\)。这可用于不求全部积分而从根的位置直接计算平均。

\ProseParagraphBreak

Jensen 公式控制的是平均对数，不是模的算术平均。由对数凹性，
\[
 \exp\left(\frac1{2\pi}\int_0^{2\pi}\log|f(re^{\mathrm{i}t})|\,\mathrm{d}t\right)
 \le\frac1{2\pi}\int_0^{2\pi}|f(re^{\mathrm{i}t})|\,\mathrm{d}t.
\]
左侧是几何平均。边界上可能有很小的函数值，也可能有很大的值；几何平均仍把所有内部零点贡献准确相加。例如 \(f(z)=e^{az}\) 没有零点，其对数平均为零，而模的平均一般大于一。零点资料完全看不见指数整函数中的调和对数模，正好与上一章 Riesz 测度的自由度相符。

\ProseParagraphBreak

Poisson--Jensen 公式的单个零点项也能显式核验。在单位圆盘，取 \(f(z)=z-a\)，\(|a|<1\)，则
\[
\log|z-a|+g_{\mathbb D}(z,a)=\log|1-\overline az|.
\]
右边无零点且调和，边界上两项对数模相等。因而 Poisson 延拓的是 \(\log|1-\overline az|\)，减去 Green 核后才恢复真正的 \(\log|z-a|\)。若直接对边界 \(\log|f|\) 作 Poisson 积分并宣称得到内部 \(\log|f|\)，就遗漏了零点的负对数奇性。错误在内部零点附近最明显：Poisson 解有限，原对数模却趋负无穷。

\ProseParagraphBreak

亚纯版本可通过有限因子分解完整证明。设 \(f\) 在闭圆盘邻域亚纯，边界无零点、极点，中心既不是零点也不是极点。内零点为 \(a_j\)、极点为 \(b_k\)，均按重数重复。将 \(f\) 乘 \(\prod(z-b_k)\) 再除 \(\prod(z-a_j)\)，所得函数在闭圆盘无零点全纯。对它应用调和平均，分别计算因子平均，得到
\[
\frac1{2\pi}\int\log|f(Re^{\mathrm{i}t})|\,\mathrm{d}t
=\log|f(0)|+\sum_j\log(R/|a_j|)-\sum_k\log(R/|b_k|).
\]
对一般内部点同样加零点 Green 核、减极点 Green 核，奇性抵消后使用 Poisson 表示即可。中心为零点或极点时，先提出整数次幂，指数允许为负。这个公式精确记录零点与极点的差，但它们数量同时很大时可能抵消，所以并不单独控制两类总量。

\ProseParagraphBreak

零点估计中选取半径比会决定常数。取 \(S=2r\) 得 \(n(r)\le(\log M(2r)-\log|f(0)|)/\log2\)；若取 \(S=(1+\varepsilon)r\)，分母 \(\log(1+\varepsilon)\) 很小，估计通常变弱。若 \(f\) 只在固定圆盘全纯，不能为了得到方便的二倍半径超出定义域。靠近外边界时，应选仍在域中的 \(S\)，并保留其实际系数。

\ProseParagraphBreak

作为精确的禁止条件，假设单位圆盘有界非恒零全纯函数的零点为 \(a_n=1-1/n\)、\(n\ge2\)。因为 \(\sum(1-|a_n|)=\sum1/n\) 发散，这种函数不存在；但没有有界性限制时，该点列在圆盘内部无聚点，仍可通过区域上的零点构造实现。当前整平面 Weierstrass 定理尚不能直接用于这类边界聚点，须使用区域版本，或者取 \(\sin(\pi/(1-z))\) 在圆盘上的明确函数，它确有这些零点并在接近边界时无界。这把解析存在性与有界函数的增长限制分开。


\ProseParagraphBreak

有限零点还能给逐点估计。对 \(a\in\mathbb D\setminus\{0\}\) 定义
\[
b_a(z)=\frac{|a|}{a}\frac{a-z}{1-\overline az},
\qquad b_0(z)=z.
\]
它在 \(a\) 有单零点，圆周模一，圆盘内部模小于一。若 \(|f|\le1\)，将任意有限批零点提出，\(F=f/\prod b_{a_j}\) 在圆盘全纯。在 \(|z|=r\uparrow1\) 上，有限乘积的模一致趋一，故最大模原理给 \(|F(z)|\le1\)。因此
\[
|f(z)|\le\prod_j|b_{a_j}(z)|.
\]
原点处即 \(|f(0)|\le\prod|a_j|\)，与 Jensen 完全一致；在零点以外的其他观察点，这个估计保留了角位置，通常比只按半径计数的信息更多。若原点是零点，则先提出原点因子，不能在原点两边同时为零的平凡不等式中读取其他零点数量。

\ProseParagraphBreak

对于任意离散零点列，Blaschke 条件在单位圆盘也是有界实现的充分条件。这可以在当前章独立构造，毋须将下一章的整平面乘积定理倒用。设 \(\sum(1-|a_j|)<\infty\)，只需先分离原点的有限重数。写 \(a=re^{\mathrm{i}\theta}\)，则
\[
1-b_a(z)=(1-r)\frac{1+e^{-\mathrm{i}\theta}z}{1-r e^{-\mathrm{i}\theta}z},
\quad
\sup_{|z|\le s}|1-b_a(z)|\le\frac{1+s}{1-s}(1-r),\ s<1.
\]
尾部因此在每个内部紧集上有可和误差。取误差充分小的尾因子，用原点附近的对数级数 \(\log(1+w)\)，其模不超过 \(2|w|\)，所以对数和局部一致收敛，指数给非零尾乘积。有限头部保留准确零点及重数，所有部分乘积模不超过一，极限也是有界整于圆盘的函数。乘积在原点为 \(\prod|a_j|>0\)：从足够大的 \(j\) 起，\(-\log|a_j|\le2(1-|a_j|)\)，故对数和有限。这证明极限不恒零。必要性与充分性由此在当前范围完整配对。

\ProseParagraphBreak

注意该乘积是单位圆盘全纯函数，一般不能延拓为整个复平面的整函数。其分母的零点在 \(1/\overline a_j\)，可能向圆周聚集；构造只在内部紧集上收敛。称为 Blaschke 乘积时，应把定义域和模界一起记住。条件 \(\sum(1-|a_j|)<\infty\) 也不保证普通边界每一点的连续性，径向边界理论将由后面的函数空间条件进一步研究。

\ProseParagraphBreak

例如 \(a_j=1-2^{-j}\) 满足条件，给一个无穷零点且有界的圆盘全纯函数。点列只有边界聚点一，故不违反恒等定理。若改成 \(a_j=1-j^{-1}\)，条件失去，可在圆盘全纯但必须无界的函数中实现。若将每个 \(1-2^{-j}\) 重复 \(2^j\) 次，按不同零点看仍是原点列，但带重数的缺陷和变成 \(\sum1\)，再次不能有界实现。重数在有界零点判据中的作用不亚于位置。

\ProseParagraphBreak

即使没有统一最大模上界，圆周上的模积分也能控制零点。若某个 \(p>0\) 满足
\[
\sup_{0<r<1}\frac1{2\pi}\int_0^{2\pi}|f(re^{\mathrm{i}t})|^p\,\mathrm{d}t\le A<\infty,
\]
对数凹性给 \((2\pi)^{-1}\int\log|f|\le p^{-1}\log A\)，先以正截断避免零点，再取极限。Jensen 的零点和仍有统一上界，所以非恒零函数继续满足 Blaschke 条件。这一推导解释为何后面的 \(H^p\) 空间即使含有无界函数，零点仍受同一个可和性限制。此处已给所需零点论证，不把后面的空间名称当成未建立的前置。

\ProseParagraphBreak

假设中 \(p>0\) 不能简单换成“平均 \(\log|f|\) 有有限值”：每个固定内部圆周的平均对数本来就是有限的，只有对所有半径的共同上界才能控制边界零点和。量词的差别在 \(\sin(\pi/(1-z))\) 上可看出，它在每个内部圆周的平均都有限，但随着半径逼近一，越来越多零点进入，累计对数和无界。逐个半径的良定义不等于统一增长控制。

\ProseParagraphBreak

有界函数的零点分解还可以连同余因子的模界一起得到。设 \(f\not\equiv0\)、\(|f|\le1\)，取含全部零点的 Blaschke 乘积 \(B\)，包括原点的有限幂。每个有限部分乘积除出后，前面已经证明所得函数的模不超过一。在不含零点的紧集上令部分乘积趋于 \(B\)，便有 \(|f/B|\le1\)；在共同零点处，重数恰好相同，所以商的奇点可去。故
\[
 f=B F,\qquad F\text{ 在圆盘全纯、无零点，且 }|F|\le1.
\]
无零点并不意味着常数。正调和函数 \(-\log|F|\) 记录全部零点之后仍剩下的自由度。
\ProseParagraphBreak
若再设 \(f(0)\ne0\)，Harnack 不等式可把中心值变成内部下界。由于 \(|B(0)|\le1\)，有 \(|F(0)|\ge|f(0)|\)。对 \(|z|\le r<1\)，令 \(v=-\log|F|\ge0\)，便得
\[
 |F(z)|=e^{-v(z)}\ge |F(0)|^{(1+r)/(1-r)}
                  \ge |f(0)|^{(1+r)/(1-r)}.
\]
因此 \(|f(z)|\ge|B(z)|\,|f(0)|^{(1+r)/(1-r)}\)。这个估计保留了每个零点对观察点的伪双曲距离；若 \(z\) 靠近零点，\(|B(z)|\) 自然趋零，估计不会错误地宣称函数处处远离零。若已知某个紧集与零点保持分离，再估计其中的乘积，才能得到实际的正下界。
\ProseParagraphBreak
一个特别有启发性的无零点余因子是
\[
 S(z)=\exp\left(-\frac{1+z}{1-z}\right),\qquad
 -\log|S(z)|=\frac{1-|z|^2}{|1-z|^2}.
\]
右端是边界点1的 Poisson 核，所以在圆盘正，\(|S|<1\)，且无零点。对每个 \(e^{\mathrm{i}\theta}\ne1\)，径向边界模趋于一；沿正实轴趋向1时，函数模却趋于零。Jensen 平均始终为 \(\log|S(0)|=-1\)。如果只把几乎处处的边界模一代入平均，就会得到零，丢失了集中到单个边界点的质量。这说明内部 Poisson 表示与边界极限交换需要可积控制，也说明零点资料不能决定一个有界全纯函数。

% !TeX root = ../../Book3.tex
\section{次调和函数观点}
\label{b3:sec:1803}

Jensen 公式还可以读成一个势论恒等式：圆周平均之差等于圆盘内部 Laplace 测度的对数势。零点只是这一原则的离散实例。

\subsection{\texorpdfstring{\(\log |f|\)}{log|f|}}
\label{b3:subsec:180301}

函数 \(\log|f|\) 在零点外调和，零点处局部为 \(m\log|z-a|\) 加调和函数。这个分解不依赖对数分支选择，因为只取实部。它也使每个局部可积性与奇性符号都可以直接核验。

\subsection{作为 Riesz 测度的零点}
\label{b3:subsec:180302}

\begin{proposition}\label{b3:prop:18-zero-measure}
区域上非恒零全纯函数满足
\[
\frac1{2\pi}\Delta\log|f|=\sum_am_a\,\delta_a,
\]
其中右侧记号 \(m_a\) 表示 \(f\) 在 \(a\) 的零点重数。
\end{proposition}
\begin{proof}
在非零点邻域 Laplace 为零；在零点邻域用局部分解与对数基本解，得到重数乘点质量。零点局部有限，因此各局部恒等式通过有限划分拼成分布等式。
\end{proof}
对亚纯函数，把极点计为负重数，即得相应符号测度。

\subsection{势论解释}
\label{b3:subsec:180303}

若 \(u\in C^2\) 次调和，中心圆周平均 \(m(r)\) 的导数为 \(\mu_u(B(0,r))/r\)，所以
\[
m(R)-u(0)=\int_0^R\frac{\mu_u(B(0,t))}t\,\mathrm{d}t.
\]
证明只需积分圆周平均导数，归一化由 \(\mu_u=\Delta u/(2\pi)\) 确定。光滑逼近可把这一公式推广到适当的次调和代表与非奇中心。对 \(u=\log|f|\)，测度离散，右端正是零点对数计数，因而回到 Jensen。

中心为对数奇点时，\(u(0)=-\infty\)，不能直接相减；必须先分离中心点质量。这是原点重数项出现的原因。

\subsection{与 Nevanlinna 理论的联系}
\label{b3:subsec:180304}

对于亚纯函数，零点与极点可能同时增长。单独看 \(\log|f|\) 的平均会发生抵消，而 \(\log^+|f|\) 的平均与极点积分计数之和能够保留总增长信息。这引向\appref{b3:app:C}的 Nevanlinna 特征。

目前已经得到的是零点的精确势论记录，并没有从 Jensen 一式直接获得第二基本定理。后者还需要对数导数和多个值之间的比较。


\ProseParagraphBreak

一般次调和平均公式可借上一章的局部 Riesz 分解在书内推出。取较大内部圆盘，使 \(u=V+h\)，\(V(z)=\int\log|z-w|\,\mathrm{d}\mu(w)\)，\(h\) 调和。对固定圆周的核平均，有
\[
\frac1{2\pi}\int_0^{2\pi}\log|re^{\mathrm{i}t}-w|\,\mathrm{d}t=\log\max(r,|w|).
\]
若 \(u(0)>-\infty\)，势在中心亦有限，故减去中心值后，核贡献 \(\log^+(r/|w|)\)，这是非负函数。Tonelli 给
\[
\frac1{2\pi}\int u(re^{\mathrm{i}t})\,\mathrm{d}t-u(0)
=\int_{|w|<r}\log(r/|w|)\,\mathrm{d}\mu(w)
=\int_0^r\frac{\mu(B(0,t))}{t}\,\mathrm{d}t.
\]
最后一步将 \(\log(r/|w|)\) 写成 \(\int_{|w|}^r\mathrm{d}t/t\)，再换非负积分次序。圆周上原势的负部可积由此也得到控制；无需未经说明地从分布等式取点值。

\ProseParagraphBreak

这个平均差随 \(r\) 非减，因为 Riesz 测度为正。在对数半径 \(s=\log r\) 下，平均函数的导数为 \(\mu(B(0,e^s))\)，它也非减，所以平均关于 \(s\) 凸。对全纯对数模，导数就是零点整数计数；对一般次调和函数，导数可以连续变化。取 \(u(z)=|z|^2\) 时圆周平均为 \(r^2\)，\(\mu(B(0,r))=2r^2\)，导数为 \(2r\)，恰与公式 \(\mu(B(0,r))/r\) 一致。这同时核验面积归一化和 \(2\pi\) 常数。

\ProseParagraphBreak

若中心是零点，\(u(0)=-\infty\)，则两边不能写作有限数相减。对 \(u=m\log|z|+v\)、\(v(0)\) 有限，圆周平均先保留 \(m\log r\)，再对 \(v\) 使用上述公式；测度中先去掉 \(m\delta_0\)。若任意在原点补一个有限值，Riesz 分布不会因单点修改改变，但次调和平均不等式会失败。这说明分布等价类与点态次调和代表应保持区别，零点处的负无穷值具有数学作用。

\ProseParagraphBreak

对于亚纯函数的 \(\log|f|\)，极点给正无穷的对数奇性，函数本身一般既非次调和也非超调和；它的 Laplace 是符号测度。可以局部写成两个次调和函数之差，对零、极点分别应用平均公式，再取差。符号测度下积分换序需要有限性或正负分解，不能继续把整个表达当成非负积分使用 Tonelli。将正测度论证中的每一步直接复制给符号测度，是从全纯版本过渡到亚纯版本时常见的漏洞。

\ProseParagraphBreak

Nevanlinna 特征保留大小的正部及极点总计数。此处先作一个算术检验：令 \(f(z)=1/z\)，对 \(r>1\)，圆周上 \(\log^+|f|=0\)，但中心极点计数贡献 \(\log r\)，所以特征仍增长；对 \(f(z)=z\)，没有极点，正部平均为 \(\log r\)。两函数互为倒数，粗略增长规模相同，单独的平均 \(\log|f|\) 却一正一负。正部与极点项的结合正是为了防止这种符号抵消丢失增长资料。完整定义及中心归一化将在相应附录中给出，这里只解释为何需要新的量。

\ProseParagraphBreak

零点测度的稳定性也可以从对数模的局部可积收敛看出。下面补上这项分析事实。仍设 \(f_j\to f\not\equiv0\) 局部一致，选一个闭圆盘，其边界无极限零点，内部零点有限。由前面的稳定计数，令 \(p_j\) 为盘中全部零点构成的首一多项式，\(p\) 为极限零点多项式，则 \(p_j\to p\) 一致。商 \(h_j=f_j/p_j\) 在盘内全纯无零点，边界上一致趋于 \(h=f/p\)；最大模原理把收敛传到闭盘。\(h\) 无零点且闭盘上有正下界，故 \(\log|h_j|\to\log|h|\) 一致。
\ProseParagraphBreak
剩下只需处理移动的对数核。若 \(a_j\to a\)，则 \(\log|z-a_j|\to\log|z-a|\) 在固定紧集上依 \(L^1\) 收敛。证明时先删去 \(a\) 附近半径 \(2\delta\) 的小圆，在其外一致收敛；其内，所有充分大的 \(a_j\) 落在半径 \(\delta\) 内，负部积分由
\(2\pi\int_0^{3\delta}r|\log r|\,\mathrm{d}r\) 控制，趋于零。正部在固定有界区域统一有界，所删小圆的面积也趋零。两部分合并即得结论。把有限个移动零点项相加，便得到
\[
 \log|f_j|\longrightarrow\log|f|\quad\text{于 }L^1_{\mathrm{loc}}(D).
\]
对紧支集光滑测试函数 \(\varphi\)，积分 \(\int\log|f_j|\Delta\varphi\) 遂收敛；由零点的 Riesz 测度公式，正好再次得到零点测度的弱收敛。点态对数在移动零点处可能完全没有共同下界，局部可积性才是传递 Laplace 算子的合适层次。
\ProseParagraphBreak
这项结论也说明为什么在正规族中控制函数本身，经常足以控制零点，但前提是先排除恒零极限。若极限恒零，\(\log|f_j|\) 可能整体趋于负无穷，即使零点数并未增加；例如常数 \(1/j\) 的对数在每个点都趋于负无穷，而其 Laplace 始终为零。对数的加性常数与零点测度是两种独立资料，不能从测度收敛反推对数模已经收敛。


\begin{chaptersummary}
零点计数有三个相互联系的表达：辐角积分给出围道内的总重数，积分计数给每个零点赋予径向对数权重，对数模的 Laplace 则把重数写成点质量。整数计数在观察圆周经过零点时跳跃，积分计数与圆周平均保持连续；在对数半径上，斜率的跳跃正好记录该圆周上的质量。

\ProseParagraphBreak
Jensen 公式由圆周对数核平均和有限零点分解得到，中心有零点时保留其幂次及首个非零系数。Poisson--Jensen 公式用正 Green 核表示任意内部点的零点影响：零点项从边界调和延拓中减去，极点项则加回。边界零点的对数奇性可积，公式能否使用取决于所取版本及其边界处理，不能仅凭路径经过零点就否定一切平均表达。

\ProseParagraphBreak
函数增长给零点分布加上定量限制。两半径估计利用每个内部零点的共同正权重；在单位圆盘中让外半径趋于一，就得到计入重数的 Blaschke 条件。归一化因子使尾项趋于一，局部对数级数保证乘积收敛并保留精确重数。除去全部零点后仍有一个有界、无零点的全纯余因子，它可以非常数，零点资料因此不足以恢复原函数。

\ProseParagraphBreak
次调和平均公式把离散零点推广为一般正质量，其单侧导数区分开、闭圆盘。非零全纯极限下，Rouch\'e 定理先稳定局部零点簇，再由移动对数核的可积收敛得到对数模的局部 \(L^1\) 稳定性。极限恒零或同时出现正、负零极点质量时，需要重新检查归一化和测试函数范围；原有的正测度结论不能直接照搬。
\end{chaptersummary}

\BookExercises

% 自拟：把圆盘径向平均推广到圆环，区分整数绕数与盘内零点数，并解释允许负斜率的原因。
\Needspace{6\baselineskip}
\begin{exercise}[圆环上的对数平均]\label{b3:ex:18-draft-01}
设 \(0<\rho<R<\infty\)，\(f\) 在圆环 \(A=\{\rho<|z|<R\}\) 全纯且不恒为零。令
\[
 M(t)=\frac1{2\pi}\int_0^{2\pi}\log|f(e^{t+\mathrm i\theta})|\,\mathrm d\theta,
 \quad \log\rho<t<\log R.
\]
证明 \(M\) 连续，在每个紧区间上分段仿射，并在一维分布意义下满足
\[
 M''=\sum_{a\in Z(f)}m_a\delta_{\log|a|}.
\]
在圆周 \(|z|=e^t\) 无零点时，证明
\[
 M'(t)=\frac1{2\pi\mathrm i}\int_{|z|=e^t}\frac{f'(z)}{f(z)}\,\mathrm dz\in\mathbb Z.
\]
由此写出从任一无零点初始圆周 \(t_0\) 到 \(t>t_0\) 的平均增量公式，并举例说明 \(M\) 虽然凸，却未必非减。
\end{exercise}
\begin{solution}
取任意紧区间 \([t_1,t_2]\subset(\log\rho,\log R)\)，在其对应的闭圆环外再留出一点余量。较大的闭圆环中只有有限个零点，除去相应有限多项式因子后，余函数在这段圆环附近全纯无零点。因此
\[
 \log|f(z)|=h(z)+\sum_a m_a\log|z-a|,
\]
其中 \(h\) 调和，和式有限。即使余函数在整个圆环上未必有单值全纯对数，其对数模仍由局部全纯对数的实部给出，是单值调和函数。

令 \(H(t)\) 为 \(h(e^{t+\mathrm i\theta})\) 的角平均。极坐标下的调和方程给
\[
 \partial_t^2h(e^{t+\mathrm i\theta})
 +\partial_\theta^2h(e^{t+\mathrm i\theta})=0.
\]
对 \(\theta\) 积分时，周期性使第二项积分为零，故 \(H''=0\)。由圆周对数平均公式，在 \([t_1,t_2]\) 上有
\[
 M(t)=\alpha t+\beta+\sum_a m_a\max\{t,\log|a|\}.
\]
这同时证明了平均的有限性、连续性和局部分段仿射性质。函数 \(\max\{t,s\}\) 的二阶分布导数为 \(\delta_s\)，于是得到题中的分布等式；它的右端在每个紧区间上只有有限项。

若观察圆周无零点，可以在邻近圆周一致地对对数模求导，得到
\[
 M'(t)=\operatorname{Re}\frac1{2\pi}
 \int_0^{2\pi}\frac{e^{t+\mathrm i\theta}f'(e^{t+\mathrm i\theta})}
 {f(e^{t+\mathrm i\theta})}\,\mathrm d\theta.
\]
右侧正是题中围道积分的实部。该围道积分是闭曲线 \(f(e^{t+\mathrm i\theta})\) 绕原点的整数绕数，因而本身为实数，得到所需等式。这里 \(f\) 未必定义在圆环内孔中，不能把这个整数直接解释为整张圆盘内的零点数。

记初始圆周的整数为 \(k=M'(t_0)\)。每经过一个零点半径，斜率增加该圆周上的总重数，故
\[
 M(t)=M(t_0)+k(t-t_0)
 +\sum_{e^{t_0}<|a|<e^t}m_a\bigl(t-\log|a|\bigr).
\]
把终点圆周上的零点加进和式不改变值，因为它们的权重为零。二阶导数为正测度，说明 \(M\) 凸；但初始整数 \(k\) 可以为负。例如 \(f(z)=1/z\) 在圆环全纯无零点，\(M(t)=-t\)，在整个区间严格递减。
\end{solution}

% 自拟：求固定中心值下开、闭圆盘的最优整数计数，并完成极值构造与等号分类。
\Needspace{6\baselineskip}
\begin{exercise}[最优零点计数与边界]\label{b3:ex:18-draft-02}
固定 \(0<c<1\)、\(0<r<1\)。在所有满足 \(|f|\le1\)、\(f(0)=c\) 的单位圆盘全纯函数中，分别求闭圆盘 \(|z|\le r\) 与开圆盘 \(|z|<r\) 内零点总重数的最大值。令
\[
 T=\frac{\log c}{\log r},
\]
证明两个答案分别为 \(\lfloor T\rfloor\) 与 \(\lceil T\rceil-1\)，并构造达到它们的函数。若闭圆盘中的零点总重数恰等于实数上界 \(T\)，确定全部这样的函数。
\end{exercise}
\begin{solution}
将闭圆盘中零点按重数列为 \(a_1,\ldots,a_N\)，因 \(f(0)>0\)，这些点都非零。取归一化因子
\[
 b_a(z)=\frac{|a|}{a}\frac{a-z}{1-\overline a z},\quad
 B(z)=\prod_{j=1}^Nb_{a_j}(z).
\]
有限 Blaschke 除法给出 \(f=Bh\)、\(|h|\le1\)，从而
\[
 c=|f(0)|\le B(0)=\prod_{j=1}^N|a_j|\le r^N.
\]
由于 \(\log r<0\)，这意味着 \(N\le T\)，故闭圆盘计数至多 \(\lfloor T\rfloor\)。若所列零点都在开圆盘且 \(N\ge1\)，第二个不等式严格，故 \(N<T\)；\(N=0\) 时同样有 \(N<T\)，因为 \(T>0\)。开圆盘中的整数上界因而为 \(\lceil T\rceil-1\)。

为达到闭圆盘上界，令 \(N=\lfloor T\rfloor\)。若 \(N=0\)，取 \(f=c\)。若 \(N\ge1\)，任选 \(N\) 个模为 \(r\) 的点，允许按重数重复，以它们组成 \(B\)，并取
\[
 f(z)=\frac{c}{r^N}B(z).
\]
这里 \(c\le r^N\)，所以 \(|f|\le1\)，且 \(f(0)=c\)，零点恰为选定的 \(N\) 个。

对于开圆盘，令 \(N=\lceil T\rceil-1\)。\(N=0\) 时仍取常函数。若 \(N\ge1\)，则 \(N<T\) 给出 \(c<r^N\)，可以选择 \(c^{1/N}\le s<r\)。将上面的构造改成 \(N\) 个模为 \(s\) 的零点及系数 \(c/s^N\)，即可达到这个上界。

最后，若闭圆盘计数为 \(N=T\)，则 \(T\) 必须是正整数，且 \(c=r^N\)。前面的整条不等式都取等号，故每个 \(|a_j|=r\)，并有 \(|h(0)|=1\)。最大模原理使 \(h\) 为模一常数；又因归一化后 \(B(0)=r^N=c=f(0)\)，这个常数只能为一。因此全部极值函数正是由 \(N\) 个模为 \(r\) 的零点组成的归一化有限 Blaschke 乘积，不含其他零点。
\end{solution}

% 对原自拟题实质改编：以双重对数修正检验 Blaschke 可和性的两个临界指数。
\Needspace{6\baselineskip}
\begin{exercise}[重数的临界可和性]\label{b3:ex:18-draft-03}
给定 \(\beta,\gamma\in\mathbb R\)，对 \(n\ge3\) 令
\[
 a_n=1-\frac1{n^2},\quad
 m_n=\max\left\{1,
 \left\lfloor\frac{n}{(\log n)^\beta(\log\log n)^\gamma}\right\rfloor\right\}.
\]
确定哪些 \((\beta,\gamma)\) 使这些点以指定重数恰为某个非恒零有界圆盘全纯函数的全部零点。解答须同时给出必要性与存在性。
\end{exercise}
\begin{solution}
对任意固定实数 \(\beta,\gamma\)，有
\[
 x_n=\frac{n}{(\log n)^\beta(\log\log n)^\gamma}\longrightarrow\infty.
\]
这是因为 \(\log x_n=\log n-\beta\log\log n-\gamma\log\log\log n\)，后两项都为 \(o(\log n)\)。因此取整与最大值只改变有限开头及渐近相对误差，\(m_n/x_n\to1\)。带重数的 Blaschke 和与下列级数有相同敛散性：
\[
 \sum_{n\ge3}m_n(1-a_n)
 \quad\text{与}\quad
 \sum_{n\ge3}\frac1{n(\log n)^\beta(\log\log n)^\gamma}.
\]
后一被积型函数最终递减，积分判别法和代换 \(u=\log x\) 将问题化为
\[
 \int_{\log3}^\infty\frac{\mathrm du}{u^\beta(\log u)^\gamma}.
\]
当 \(\beta>1\) 时，\((\log u)^{|\gamma|}\) 的增长慢于任意正幂，故可与某个可积的 \(u^{-p}\)、\(p>1\) 比较。当 \(\beta<1\) 时，同样的比较给出发散。若 \(\beta=1\)，再令 \(v=\log u\)，便得到 \(\int v^{-\gamma}\,\mathrm dv\)，它收敛当且仅当 \(\gamma>1\)。

所求参数范围因而是
\[
 \boxed{\ \beta>1\quad\text{或}\quad\beta=1,\ \gamma>1\ }.
\]
必要性来自\cref{b3:thm:18-blaschke-zeros}。在这个范围内，按重数重复各 \(a_n\) 所得点列在圆盘内部无聚点，且 Blaschke 和有限。同一定理构造的乘积模不超过一、非恒零，并在每个 \(a_n\) 恰有 \(m_n\) 重零点，没有其他零点，从而证明充分性。
\end{solution}

% 自拟：从单因子误差导出乘积截断误差，并计算几何零点列的首阶尾乘积渐近。
\Needspace{6\baselineskip}
\begin{exercise}[Blaschke 尾乘积的误差]\label{b3:ex:18-draft-04}
设非零圆盘点列 \((a_j)_{j\ge1}\) 满足 Blaschke 条件，按重数重复，令
\[
 B_N=\prod_{j=1}^Nb_{a_j},\quad B=\prod_{j=1}^\infty b_{a_j},\quad
 T_N=\prod_{j>N}b_{a_j},\quad
 \delta_N=\sum_{j>N}(1-|a_j|).
\]
这里因子采用 \(b_a(0)=|a|\) 的归一化，\(B=B_NT_N\)。对 \(0<r<1\)，证明
\[
 \sup_{|z|\le r}|T_N(z)-1|
 \le\frac{1+r}{1-r}\delta_N,\quad
 \sup_{|z|\le r}|B(z)-B_N(z)|
 \le\frac{1+r}{1-r}\delta_N.
\]
进一步，若 \(a_j=1-q^j\)、\(0<q<1\)，证明在单位圆盘内局部一致地有
\[
 q^{-N}\bigl(T_N(z)-1\bigr)
 \longrightarrow-\frac{q}{1-q}\frac{1+z}{1-z}.
\]
\end{exercise}
\begin{solution}
对 \(M>N\)，有限乘积逐项相减给出
\[
 \prod_{j=N+1}^Mb_{a_j}-1
 =\sum_{k=N+1}^M(b_{a_k}-1)\prod_{j=N+1}^{k-1}b_{a_j}.
\]
每个因子在圆盘内模不超过一，再用正文的单因子误差界，得到
\[
 \sup_{|z|\le r}\left|\prod_{j=N+1}^Mb_{a_j}(z)-1\right|
 \le\frac{1+r}{1-r}\sum_{k=N+1}^M(1-|a_k|).
\]
令 \(M\to\infty\)，局部一致乘积收敛给出第一个估计。由 \(B-B_N=B_N(T_N-1)\) 及 \(|B_N|\le1\)，第二个估计随即成立。即使 \(B_N\) 在某点为零，乘积恒等式仍成立，无需在那里直接相除。

现在取 \(a_j=1-q^j\)。固定 \(r<1\)，记
\[
 h(z)=\frac{1+z}{1-z},\quad
 d_j(z)=1-b_{a_j}(z)
 =q^j\frac{1+z}{1-(1-q^j)z}.
\]
在 \(|z|\le r\) 上，分母与 \(1-z\) 均有共同正下界。将两分式相减可得
\[
 d_j(z)=q^jh(z)+O_r(q^{2j}),\quad |d_j(z)|=O_r(q^j),
\]
其中 \(q\) 固定，常数不依赖 \(j,z\)。当 \(N\) 足够大时，所有 \(j>N\) 都有 \(|d_j|\le1/2\)，故可以用原点附近的幂级数取 \(\log(1-d_j)\)，并一致地写成
\[
 \log(1-d_j)=-d_j+O_r(|d_j|^2).
\]
这些对数的级数绝对一致收敛，其指数正是 \(T_N\)。于是
\[
 \log T_N(z)
 =-h(z)\frac{q^{N+1}}{1-q}+O_r(q^{2N}).
\]
右侧在闭盘上一致为 \(O_r(q^N)\)，指数展开中的二次余项因而仍为 \(O_r(q^{2N})\)。所以
\[
 T_N(z)-1=-h(z)\frac{q^{N+1}}{1-q}+O_r(q^{2N}).
\]
乘以 \(q^{-N}\) 并令 \(N\to\infty\)，得到所求局部一致极限。
\end{solution}

% 自拟：在恒零极限中归一化零点质量，以圆周对数的绝对积分给出明确的 L1 误差。
\Needspace{6\baselineskip}
\begin{exercise}[趋零函数的归一化零点分布]\label{b3:ex:18-draft-05}
固定 \(0<r<1\)，令
\[
 f_n(z)=2^{-n}(z^n-r^n),\quad
 \nu_n=\frac1n\sum_{k=0}^{n-1}\delta_{r e^{2\pi\mathrm i k/n}},\quad
 u_n(z)=\frac1n\log|f_n(z)|.
\]
证明 \(f_n\) 在闭单位圆盘上一致趋于零，而 \(\nu_n\) 对紧支集连续测试函数收敛于半径 \(r\) 圆周上的均匀概率测度 \(\sigma_r\)。再证明
\[
 u_n\longrightarrow u,\quad
 u(z)=\log\max\{|z|,r\}-\log2
 \quad\text{于 }L^1_{\mathrm{loc}}(\mathbb D),
\]
并对每个 \(0<R<1\) 给出估计
\[
 \|u_n-u\|_{L^1(B(0,R))}\le\frac{2\pi R^2\log2}{n}.
\]
说明这些结论与正文的非零全纯极限条件有何关系。
\end{exercise}
\begin{solution}
在 \(|z|\le1\) 上，\(|f_n(z)|\le2^{-n}(1+r^n)\to0\)。方程 \(z^n=r^n\) 的根恰为题中列出的 \(n\) 个互异点，故 \(\nu_n\) 是零点计数测度除以 \(n\)。对 \(\varphi\in C_c(\mathbb D)\)，函数 \(t\mapsto\varphi(re^{\mathrm it})\) 连续且周期，其等分 Riemann 和给出
\[
 \int\varphi\,\mathrm d\nu_n
 =\frac1n\sum_{k=0}^{n-1}\varphi(re^{2\pi\mathrm i k/n})
 \longrightarrow\frac1{2\pi}\int_0^{2\pi}\varphi(re^{\mathrm it})\,\mathrm dt
 =\int\varphi\,\mathrm d\sigma_r.
\]

为控制对数奇性，先记 \(L_s(t)=\log|1-se^{\mathrm it}|\)，其中 \(0\le s\le1\)。圆周对数平均公式，包括 \(s=1\) 时的可积边界奇性，给出
\[
 \frac1{2\pi}\int_0^{2\pi}L_s(t)\,\mathrm dt=0.
\]
同时 \(L_s^+\le\log2\)。正、负部分积分相等，所以
\[
 \frac1{2\pi}\int_0^{2\pi}|L_s(t)|\,\mathrm dt
 =\frac1\pi\int_0^{2\pi}L_s^+(t)\,\mathrm dt
 \le2\log2.
\]
这是一条对全部 \(s\in[0,1]\) 有效的绝对积分界。

写 \(z=\rho e^{\mathrm i\theta}\)。当 \(0\le\rho\le r\) 时，从 \(z^n-r^n\) 中取出 \(-r^n\)；当 \(\rho>r\) 时取出 \(z^n\)。由此在几乎处处有
\[
 u_n(z)-u(z)=
 \begin{cases}
 n^{-1}\log|1-(\rho/r)^ne^{\mathrm in\theta}|,&0\le\rho\le r,\\
 n^{-1}\log|1-(r/\rho)^ne^{-\mathrm in\theta}|,&\rho>r.
 \end{cases}
\]
两种情形中的实参数都位于 \([0,1]\)。由于角函数为 \(2\pi\) 周期，把 \(\theta\) 换成 \(\pm n\theta\) 不改变一周平均；因此对每个固定 \(\rho\)，有
\[
 \int_0^{2\pi}|u_n(\rho e^{\mathrm i\theta})-u(\rho e^{\mathrm i\theta})|
 \,\mathrm d\theta\le\frac{4\pi\log2}{n}.
\]
对非负函数使用极坐标积分，得到
\[
 \int_{B(0,R)}|u_n-u|\,\mathrm dA
 \le\frac{4\pi\log2}{n}\int_0^R\rho\,\mathrm d\rho
 =\frac{2\pi R^2\log2}{n}.
\]
它也证明各个差值局部绝对可积，并给出所需 \(L^1\) 收敛。

分布公式还给 \((2\pi)^{-1}\Delta u_n=\nu_n\)，因而上述两个极限相容，并有 \((2\pi)^{-1}\Delta u=\sigma_r\)。这里全纯函数的极限是恒零函数，正文的对数模稳定性定理不适用；我们收敛的是另外除以 \(n\) 的对数模与零点质量。未归一化的零点总质量是 \(n\)，在一个等于一于该圆周附近的非负紧支集测试函数上趋于无穷，并没有趋向零测度。
\end{solution}

% 自拟：用高重数相邻零极点区分分布收敛与连续测试下的测度收敛，显式展示总变差失控。
\Needspace{6\baselineskip}
\begin{exercise}[零极点相消与测试函数]\label{b3:ex:18-draft-06}
对 \(n\ge2\)，令 \(\varepsilon_n=n^{-2}\)，并在单位圆盘考虑亚纯函数
\[
 F_n(z)=\left(\frac{z-\varepsilon_n}{z+\varepsilon_n}\right)^n,
 \quad
 \lambda_n=n\delta_{\varepsilon_n}-n\delta_{-\varepsilon_n}.
\]
证明 \(\log|F_n|\to0\) 于 \(L^1_{\mathrm{loc}}(\mathbb D)\)，且 \(\lambda_n\) 在分布意义下趋于零。构造一个 \(\varphi\in C_c(\mathbb D)\)，使 \(\int\varphi\,\mathrm d\lambda_n\) 不趋于零，并证明 \(\lambda_n\) 的正、负部分在原点的每个固定邻域内都没有共同有限质量上界。
\end{exercise}
\begin{solution}
对不在实线段 \([-\varepsilon,\varepsilon]\) 上的 \(z\)，沿实参数积分 \(\log|z-t|\) 的导数可得
\[
 \bigl|\log|z-\varepsilon|-\log|z+\varepsilon|\bigr|
 \le\int_{-\varepsilon}^{\varepsilon}\frac{\mathrm dt}{|z-t|}.
\]
排除的线段面积为零，不影响下面的面积积分。对 \(0<R<1\)、\(|t|\le1\)，平移后将积分区域放大为半径 \(R+1\) 的圆盘，得到
\[
 \int_{B(0,R)}\frac{\mathrm dA(z)}{|z-t|}
 \le\int_{B(0,R+1)}\frac{\mathrm dA(w)}{|w|}
 =2\pi(R+1).
\]
右侧核非负，可以用 Tonelli 定理换序。乘上 \(n\) 并取 \(\varepsilon=\varepsilon_n\)，便有
\[
 \int_{B(0,R)}|\log|F_n(z)||\,\mathrm dA(z)
 \le4\pi(R+1)n\varepsilon_n
 =\frac{4\pi(R+1)}n\longrightarrow0.
\]
每个内部紧集都包含于这样的圆盘，故局部 \(L^1\) 收敛成立。零点和极点处的无穷值按局部可积等价类理解。

由亚纯零极点公式，\((2\pi)^{-1}\Delta\log|F_n|=\lambda_n\)，局部 \(L^1\) 收敛已经推出分布收敛。也可以直接检查：对 \(\psi\in C_c^\infty(\mathbb D)\)，当 \(n\) 足够大时，
\[
 \left|\int\psi\,\mathrm d\lambda_n\right|
 =n|\psi(\varepsilon_n)-\psi(-\varepsilon_n)|
 \le2n\varepsilon_n\sup|\partial_x\psi|
 \longrightarrow0.
\]
这个估计用到了光滑测试函数的一阶变化控制。

现在取非负 \(\chi\in C_c^\infty(\mathbb D)\)，使其在原点附近恒等于一，定义
\[
 \varphi(x+\mathrm iy)=\chi(x+\mathrm iy)\operatorname{sgn}(x)\sqrt{|x|},
 \quad \operatorname{sgn}(0)\sqrt0=0.
\]
该函数连续且有紧支集，但在 \(x=0\) 附近不满足前面的一阶变化估计。对充分大的 \(n\)，有
\[
 \int\varphi\,\mathrm d\lambda_n
 =n\left(\sqrt{\varepsilon_n}-(-\sqrt{\varepsilon_n})\right)=2.
\]
所以光滑测试下的分布收敛并没有推广为全部紧支集连续测试下的测度收敛。

两个支点互异，故 Jordan 分解为
\[
 \lambda_n^+=n\delta_{\varepsilon_n},\quad
 \lambda_n^-=n\delta_{-\varepsilon_n}.
\]
在任何固定原点邻域内，两点最终都被包含，正、负质量各为 \(n\)。特别地，对前面的 \(\chi\)，最终有
\(\int\chi\,\mathrm d\lambda_n^+=\int\chi\,\mathrm d\lambda_n^-=n\)，总变差积分为 \(2n\)。缺少局部总变差上界，便不能把连续函数的光滑逼近误差统一控制；这是两类测试得出不同结论的具体原因。
\end{solution}

% 对原自拟题实质改编：在有限零点除法后加入两个插值条件，以 Schwarz 引理给出可解判据和刚性。
\Needspace{6\baselineskip}
\begin{exercise}[指定零点后的两点插值]\label{b3:ex:18-draft-07}
设 \(B\) 为单位圆盘内的有限 Blaschke 乘积，\(p,q\in\mathbb D\) 互异且 \(B(p)B(q)\ne0\)，给定 \(c,d\in\mathbb C\)。求存在圆盘全纯函数 \(f\) 满足
\[
 |f|\le1,\quad f/B\text{ 在圆盘全纯},\quad f(p)=c,\quad f(q)=d
\]
的充要条件。商条件表示 \(f\) 在指定各点至少具有 \(B\) 的相应零点重数；允许额外零点，也允许 \(f\equiv0\)。记
\[
 A=\frac{c}{B(p)},\quad C=\frac{d}{B(q)},\quad
 T_a(z)=\frac{z-a}{1-\overline a z}\quad(a\in\mathbb D).
\]
分别讨论 \(|A|>1\)、\(|A|=1\) 和 \(|A|<1\)。在最后一种情形下，若所得的 Schwarz 距离不等式取等号，确定解的形式及唯一性。
\end{exercise}
\begin{solution}
设已有这样的 \(f\)，令 \(h=f/B\)。有限 Blaschke 乘积在单位圆周上模为一，并连续到圆周。在包住其全部零点的圆周 \(|z|=s<1\) 上，最大模原理给
\[
 |h(z)|\le\frac1{\min_{|\zeta|=s}|B(\zeta)|}\quad(|z|<s).
\]
令 \(s\uparrow1\)，可得 \(|h|\le1\)。反过来，任意这样的 \(h\) 都使 \(f=Bh\) 满足模界和零点要求。因此问题等价于求圆盘全纯函数 \(h\)，使 \(|h|\le1\)、\(h(p)=A\)、\(h(q)=C\)。

若 \(|A|>1\)，无解。若 \(|A|=1\)，最大模原理强制 \(h\equiv A\)，所以有解当且仅当 \(C=A\)，此时唯一解为 \(f=AB\)。

若 \(|A|<1\)，任何可行 \(h\) 都将圆盘映入开圆盘；否则它在内部取到模一值，便应为模一常数，与 \(h(p)=A\) 矛盾。因而必须有 \(C\in\mathbb D\)。函数
\[
 g=T_A\circ h\circ T_p^{-1}
\]
将圆盘映入自身且 \(g(0)=0\)。Schwarz 引理在点 \(T_p(q)\ne0\) 给出必要条件
\[
 |T_A(C)|\le|T_p(q)|.
\]
反之，若 \(C\in\mathbb D\) 且该不等式成立，令
\[
 \lambda=\frac{T_A(C)}{T_p(q)},\quad
 h(z)=T_A^{-1}\bigl(\lambda T_p(z)\bigr).
\]
因为 \(|\lambda|\le1\)，这个 \(h\) 在圆盘全纯，取值在圆盘内，并满足两个插值条件。所以完整判据为
\[
 \begin{cases}
 \text{无解},&|A|>1,\\
 C=A,&|A|=1,\\
 C\in\mathbb D\ \text{且 }|T_A(C)|\le|T_p(q)|,&|A|<1.
 \end{cases}
\]

在 \(|A|<1\) 的情形，若不等式取等号，Schwarz 引理的等号结论使 \(g(\zeta)=\lambda\zeta\)、\(|\lambda|=1\)。两个插值值已唯一决定 \(\lambda\)，所以唯一的解就是
\[
 f(z)=B(z)T_A^{-1}\bigl(\lambda T_p(z)\bigr).
\]
后一因子为圆盘自同构，恰增加一个简单零点；它也可能落在 \(B\) 的既有零点处，使该点的总重数增加一。题设只要求保留指定零点及至少那些重数，因此这两种情况都允许。
\end{solution}

% 对原自拟题实质改编：从倒反不变性推进到几何平均的双侧系数界、最优常数和等号判别。
\Needspace{6\baselineskip}
\begin{exercise}[多项式的几何平均与系数界]\label{b3:ex:18-draft-08}
设非零多项式 \(P(z)=\sum_{k=0}^dc_kz^k\) 的次数为 \(d\)，定义
\[
 \mathcal M(P)=\exp\left(\frac1{2\pi}\int_0^{2\pi}
 \log|P(e^{\mathrm it})|\,\mathrm dt\right),\quad
 P^*(z)=z^d\overline{P(1/\overline z)},
\]
其中 \(P^*\) 在原点按多项式延拓。将 \(P\) 的根按重数列为 \(a_1,\ldots,a_d\)，证明
\[
 \mathcal M(P)=|c_d|\prod_{j=1}^d\max\{1,|a_j|\},\quad
 \mathcal M(P^*)=\mathcal M(P),
\]
以及
\[
 \max_{0\le k\le d}\frac{|c_k|}{\binom dk}
 \le\mathcal M(P)
 \le\left(\sum_{k=0}^d|c_k|^2\right)^{1/2}.
\]
证明每个系数不等式中的常数 \(\binom dk\) 都不能减小，并确定右端不等式取等号的全部多项式。允许原点零点与单位圆周上的零点。
\end{exercise}
\begin{solution}
因式分解给 \(P(z)=c_d\prod_{j=1}^d(z-a_j)\)。每个因子的圆周对数都绝对可积，平均为 \(\log\max\{1,|a_j|\}\)，其中单位圆周上的对数奇性也包括在内。逐项相加并指数化，得到根的乘积公式。\(d=0\) 时空乘积取一。

展开倒反定义可得
\[
 P^*(z)=\sum_{k=0}^d\overline{c_k}z^{d-k}
 =\overline{c_d}\prod_{j=1}^d(1-\overline{a_j}z).
\]
这是非零多项式；若 \(c_0=0\)，它的实际次数可以小于 \(d\)，但定义仍有效。在单位圆周上 \(|P^*|=|P|\)，两者的可积对数具有同一平均，所以 \(\mathcal M(P^*)=\mathcal M(P)\)。

系数 \(c_k\) 是 \(c_d\) 乘以 \(d-k\) 阶根的基本对称和，再乘一个符号。每个由 \(d-k\) 个根组成的乘积都满足
\[
 \prod_{j\in I}|a_j|\le\prod_{j=1}^d\max\{1,|a_j|\}.
\]
这样的项有 \(\binom d{d-k}=\binom dk\) 个，三角不等式便给
\( |c_k|\le\binom dk\mathcal M(P)\)。取全部 \(k\) 的最大值得到左界。若 \(P(z)=C(z-a)^d\)、\(C\ne0\)、\(|a|=1\)，则 \(\mathcal M(P)=|C|\)，而每个 \(|c_k|=|C|\binom dk\)。因此每个系数的常数都已最优。

为证右界，记
\[
 X(t)=|P(e^{\mathrm it})|^2,\quad
 A=\frac1{2\pi}\int_0^{2\pi}X(t)\,\mathrm dt>0.
\]
由 \(\log x\le x-1\)，在 \(X>0\) 的点有
\(\log(X/A)\le X/A-1\)。有限个圆周零点不影响积分，而 \(\log X\) 已知可积，故
\[
 2\log\mathcal M(P)-\log A
 \le\frac1{2\pi}\int_0^{2\pi}(X/A-1)\,\mathrm dt=0.
\]
再将有限 Fourier 和展开积分，异频项的积分为零，得到
\[
 A=\frac1{2\pi}\int_0^{2\pi}
 \sum_{j,k=0}^dc_j\overline{c_k}e^{\mathrm i(j-k)t}\,\mathrm dt
 =\sum_{k=0}^d|c_k|^2.
\]
这证明了右界。

等号要求 \(X/A-1-\log(X/A)\) 几乎处处为零，即 \(X=A\) 几乎处处。由连续性，\(|P|\) 在整个单位圆周为常数。设 \(m\) 是最小的非零系数指标。若 \(m<d\)，则 \(|P(e^{\mathrm it})|^2\) 的频率 \(d-m\) 项只有一项，其系数为 \(c_d\overline{c_m}\ne0\)，与常数函数矛盾。所以 \(m=d\)，\(P(z)=c_dz^d\)。反过来，任一非零单项式都在单位圆周上有恒定模，确实取等号；其中包括非零常数。
\end{solution}

